
\Data\ID{1003067701}
\Updated{2022-04-23 05:04:57}
\Level{A}
\SubArea{Complex numbers in algebraic and trigonometric form}
\LANGen{
\Question{
Evaluate  \( \mathrm{i}^{100}\cdot\mathrm{i}^{99}\cdot\mathrm{i}^{98}\cdot\mathrm{i}^{97}\cdot\dots\cdot\mathrm{i}^4\cdot\mathrm{i}^3\cdot\mathrm{i}^2\cdot\mathrm{i}\).}
{\( -1 \)}{\( 1 \)}{\( \mathrm{i} \)}{\( -\mathrm{i} \)}{}{}}
\LANGpl{
\Question{
Oblicz  \( \mathrm{i}^{100}\cdot\mathrm{i}^{99}\cdot\mathrm{i}^{98}\cdot\mathrm{i}^{97}\cdot\dots\cdot\mathrm{i}^4\cdot\mathrm{i}^3\cdot\mathrm{i}^2\cdot\mathrm{i}\).}
{\( -1 \)}{\( 1 \)}{\( \mathrm{i} \)}{\( -\mathrm{i} \)}{}{}}
\LANGcs{
\Question{
Vypočítejte:  \( \mathrm{i}^{100}\cdot\mathrm{i}^{99}\cdot\mathrm{i}^{98}\cdot\mathrm{i}^{97}\cdot\dots\cdot\mathrm{i}^4\cdot\mathrm{i}^3\cdot\mathrm{i}^2\cdot\mathrm{i}\).}
{\( -1 \)}{\( 1 \)}{\( \mathrm{i} \)}{\( -\mathrm{i} \)}{}{}}
\LANGsk{
\Question{
Vypočítajte:  \( \mathrm{i}^{100}\cdot\mathrm{i}^{99}\cdot\mathrm{i}^{98}\cdot\mathrm{i}^{97}\cdot\dots\cdot\mathrm{i}^4\cdot\mathrm{i}^3\cdot\mathrm{i}^2\cdot\mathrm{i}\).}
{\( -1 \)}{\( 1 \)}{\( \mathrm{i} \)}{\( -\mathrm{i} \)}{}{}}
\LANGes{
\Question{
Calcula  \( \mathrm{i}^{100}\cdot\mathrm{i}^{99}\cdot\mathrm{i}^{98}\cdot\mathrm{i}^{97}\cdot\dots\cdot\mathrm{i}^4\cdot\mathrm{i}^3\cdot\mathrm{i}^2\cdot\mathrm{i}\).}
{\( -1 \)}{\( 1 \)}{\( \mathrm{i} \)}{\( -\mathrm{i} \)}{}{}}
\End

\Data\ID{1003067702}
\Updated{2022-11-01 09:11:39}
\Level{A}
\SubArea{Complex numbers in algebraic and trigonometric form}
\LANGen{
\Question{
Evaluate  \( \mathrm{i}^{100}+\mathrm{i}^{99}+\mathrm{i}^{98}+\mathrm{i}^{97}+\dots+\mathrm{i}^4+\mathrm{i}^3+\mathrm{i}^2+\mathrm{i}\).}
{\( 0 \)}{\( 1-\mathrm{i} \)}{\( -1 \)}{\( -1+\mathrm{i} \)}{}{}}
\LANGpl{
\Question{
Oblicz  \( \mathrm{i}^{100}+\mathrm{i}^{99}+\mathrm{i}^{98}+\mathrm{i}^{97}+\dots+\mathrm{i}^4+\mathrm{i}^3+\mathrm{i}^2+\mathrm{i}\).}
{\( 0 \)}{\( 1-\mathrm{i} \)}{\( -1 \)}{\( -1+\mathrm{i} \)}{}{}}
\LANGcs{
\Question{
Vypočítejte:  \( \mathrm{i}^{100}+\mathrm{i}^{99}+\mathrm{i}^{98}+\mathrm{i}^{97}+\dots+\mathrm{i}^4+\mathrm{i}^3+\mathrm{i}^2+\mathrm{i} \).}
{\( 0 \)}{\( 1-\mathrm{i} \)}{\( -1 \)}{\( -1+\mathrm{i} \)}{}{}}
\LANGsk{
\Question{
Vypočítajte:  \( \mathrm{i}^{100}+\mathrm{i}^{99}+\mathrm{i}^{98}+\mathrm{i}^{97}+\dots+\mathrm{i}^4+\mathrm{i}^3+\mathrm{i}^2+\mathrm{i} \).}
{\( 0 \)}{\( 1-\mathrm{i} \)}{\( -1 \)}{\( -1+\mathrm{i} \)}{}{}}
\LANGes{
\Question{
Calcula  \( \mathrm{i}^{100}+\mathrm{i}^{99}+\mathrm{i}^{98}+\mathrm{i}^{97}+\dots+\mathrm{i}^4+\mathrm{i}^3+\mathrm{i}^2+\mathrm{i}\).}
{\( 0 \)}{\( 1-\mathrm{i} \)}{\( -1 \)}{\( -1+\mathrm{i} \)}{}{}}
\End

\Data\ID{1003067703}
\Updated{2022-04-23 05:04:57}
\Level{A}
\SubArea{Complex numbers in algebraic and trigonometric form}
\LANGen{
\Question{
Evaluate  \( \mathrm{i}^{10} -  \mathrm{i}^{20} + \mathrm{i}^{30} -  \mathrm{i}^{40} + \mathrm{i}^{50} -  \mathrm{i}^{60} \).}
{\( -6 \)}{\( 0 \)}{\( -1 \)}{\( 6 \)}{}{}}
\LANGpl{
\Question{
Oblicz  \( \mathrm{i}^{10} -  \mathrm{i}^{20} + \mathrm{i}^{30} -  \mathrm{i}^{40} + \mathrm{i}^{50} -  \mathrm{i}^{60} \).}
{\( -6 \)}{\( 0 \)}{\( -1 \)}{\( 6 \)}{}{}}
\LANGcs{
\Question{
Vypočítejte:   \( \mathrm{i}^{10} -  \mathrm{i}^{20} + \mathrm{i}^{30} -  \mathrm{i}^{40} + \mathrm{i}^{50} -  \mathrm{i}^{60} \).}
{\( -6 \)}{\( 0 \)}{\( -1 \)}{\( 6 \)}{}{}}
\LANGsk{
\Question{
Vypočítajte:   \( \mathrm{i}^{10} -  \mathrm{i}^{20} + \mathrm{i}^{30} -  \mathrm{i}^{40} + \mathrm{i}^{50} -  \mathrm{i}^{60} \).}
{\( -6 \)}{\( 0 \)}{\( -1 \)}{\( 6 \)}{}{}}
\LANGes{
\Question{
Calcula  \( \mathrm{i}^{10} -  \mathrm{i}^{20} + \mathrm{i}^{30} -  \mathrm{i}^{40} + \mathrm{i}^{50} -  \mathrm{i}^{60} \).}
{\( -6 \)}{\( 0 \)}{\( -1 \)}{\( 6 \)}{}{}}
\End

\Data\ID{1003067704}
\Updated{2021-12-19 10:12:48}
\Level{A}
\SubArea{Complex numbers in algebraic and trigonometric form}
\LANGen{
\Question{
Evaluate   \( \mathrm{i}^5 - \mathrm{i}^{10} + \mathrm{i}^{15} -  \mathrm{i}^{20} + \mathrm{i}^{25} -  \mathrm{i}^{30} + \mathrm{i}^{35} -  \mathrm{i}^{40} \).}
{\( 0 \)}{\( 1 \)}{\( -1 \)}{\( -\mathrm{i} \)}{}{}}
\LANGpl{
\Question{
Oblicz   \( \mathrm{i}^5 - \mathrm{i}^{10} + \mathrm{i}^{15} -  \mathrm{i}^{20} + \mathrm{i}^{25} -  \mathrm{i}^{30} + \mathrm{i}^{35} -  \mathrm{i}^{40} \).}
{\( 0 \)}{\( 1 \)}{\( -1 \)}{\( -\mathrm{i} \)}{}{}}
\LANGcs{
\Question{
Vypočítejte:   \( \mathrm{i}^5 - \mathrm{i}^{10} + \mathrm{i}^{15} -  \mathrm{i}^{20} + \mathrm{i}^{25} -  \mathrm{i}^{30} + \mathrm{i}^{35} -  \mathrm{i}^{40} \).}
{\( 0 \)}{\( 1 \)}{\( -1 \)}{\( -\mathrm{i} \)}{}{}}
\LANGsk{
\Question{
Vypočítajte:   \( \mathrm{i}^5 - \mathrm{i}^{10} + \mathrm{i}^{15} -  \mathrm{i}^{20} + \mathrm{i}^{25} -  \mathrm{i}^{30} + \mathrm{i}^{35} -  \mathrm{i}^{40} \).}
{\( 0 \)}{\( 1 \)}{\( -1 \)}{\( -\mathrm{i} \)}{}{}}
\LANGes{
\Question{
Calcula   \( \mathrm{i}^5 - \mathrm{i}^{10} + \mathrm{i}^{15} -  \mathrm{i}^{20} + \mathrm{i}^{25} -  \mathrm{i}^{30} + \mathrm{i}^{35} -  \mathrm{i}^{40} \).}
{\( 0 \)}{\( 1 \)}{\( -1 \)}{\( -\mathrm{i} \)}{}{}}
\End

\Data\ID{1003067705}
\Updated{2020-07-18 08:07:56}
\Level{A}
\SubArea{Complex numbers in algebraic and trigonometric form}
\LANGen{
\Question{
Given complex numbers  \(z_1 = -2 + \mathrm{i} \), \( z_2 = 1 - 4\mathrm{i} \) and \( z_3 = -3\mathrm{i} \),  find  \( 2\cdot z_1 + 3\cdot z_2- 4\cdot z_3 \).}
{\( -1 + 2\mathrm{i} \)}{\( 7 + 2\mathrm{i} \)}{\( 11 + 26\mathrm{i} \)}{\( -1 + 26\mathrm{i} \)}{}{}}
\LANGpl{
\Question{
Dano liczby zespolone  \(z_1 = -2 + \mathrm{i} \), \( z_2 = 1 - 4\mathrm{i} \) i \( z_3 = -3\mathrm{i} \),  wskaż \( 2\cdot z_1 + 3\cdot z_2 - 4\cdot z_3 \).}
{\( -1 + 2\mathrm{i} \)}{\( 7 + 2\mathrm{i} \)}{\( 11 + 26\mathrm{i} \)}{\( -1 + 26\mathrm{i} \)}{}{}}
\LANGcs{
\Question{
Jsou daná komplexní čísla:  \(z_1 = -2 + \mathrm{i} \), \( z_2 = 1 - 4\mathrm{i} \) a \( z_3 = -3\mathrm{i} \).  Vypočítejte:  \( 2 z_1 + 3 z_2  -  4 z_3 \).}
{\( -1 + 2\mathrm{i} \)}{\( 7 + 2\mathrm{i} \)}{\( 11 + 26\mathrm{i} \)}{\( -1 + 26\mathrm{i} \)}{}{}}
\LANGsk{
\Question{
Dané sú komplexné čísla:  \(z_1 = -2 + \mathrm{i} \), \( z_2 = 1 - 4\mathrm{i} \) a \( z_3 = -3\mathrm{i} \).  Vypočítajte:  \( 2 z_1 + 3 z_2  -  4 z_3 \).}
{\( -1 + 2\mathrm{i} \)}{\( 7 + 2\mathrm{i} \)}{\( 11 + 26\mathrm{i} \)}{\( -1 + 26\mathrm{i} \)}{}{}}
\LANGes{
\Question{
Dados los números complejos  \(z_1 = -2 + \mathrm{i} \), \( z_2 = 1 - 4\mathrm{i} \) y \( z_3 = -3\mathrm{i} \),  calcula  \( 2\cdot z_1 + 3\cdot z_2- 4\cdot z_3 \).}
{\( -1 + 2\mathrm{i} \)}{\( 7 + 2\mathrm{i} \)}{\( 11 + 26\mathrm{i} \)}{\( -1 + 26\mathrm{i} \)}{}{}}
\End

\Data\ID{1003067706}
\Updated{2020-12-16 10:12:08}
\Level{A}
\SubArea{Complex numbers in algebraic and trigonometric form}
\LANGen{
\Question{
Given complex numbers \( z_1 = -2 + \mathrm{i} \), \( z_2 = 1 - 4\mathrm{i} \) and \( z_3 = -3\mathrm{i} \), find  \(\frac{z_1\cdot z_2}{3\cdot z_3}\).}
{\( -1 + \frac29\mathrm{i} \)}{\( -1 -  \frac29\mathrm{i} \)}{\( 1 + \frac29\mathrm{i} \)}{\( 1 - \frac29\mathrm{i} \)}{}{}}
\LANGpl{
\Question{
Podano liczby zespolone \( z_1 = -2 + \mathrm{i} \), \( z_2 = 1 - 4\mathrm{i} \) i  \( z_3 = -3\mathrm{i} \), wskaż \(\frac{z_1\cdot z_2}{3\cdot z_3}\).}
{\( -1 + \frac29\mathrm{i} \)}{\( -1 -  \frac29\mathrm{i} \)}{\( 1 + \frac29\mathrm{i} \)}{\( 1 - \frac29\mathrm{i} \)}{}{}}
\LANGcs{
\Question{
Jsou daná komplexní čísla: \( z_1 = -2 + \mathrm{i} \), \( z_2 = 1 - 4\mathrm{i} \) a \( z_3 = -3\mathrm{i} \). Vypočítejte  \(\frac{z_1\cdot z_2}{3\cdot z_3}\).}
{\( -1 + \frac29\mathrm{i} \)}{\( -1 -  \frac29\mathrm{i} \)}{\( 1 + \frac29\mathrm{i} \)}{\( 1 - \frac29\mathrm{i} \)}{}{}}
\LANGsk{
\Question{
Dané sú komplexné čísla: \( z_1 = -2 + \mathrm{i} \), \( z_2 = 1 - 4\mathrm{i} \) a \( z_3 = -3\mathrm{i} \). Vypočítajte  \(\frac{z_1\cdot z_2}{3\cdot z_3}\).}
{\( -1 + \frac29\mathrm{i} \)}{\( -1 -  \frac29\mathrm{i} \)}{\( 1 + \frac29\mathrm{i} \)}{\( 1 - \frac29\mathrm{i} \)}{}{}}
\LANGes{
\Question{
Dados los números complejos \( z_1 = -2 + \mathrm{i} \), \( z_2 = 1 - 4\mathrm{i} \) y \( z_3 = -3\mathrm{i} \), calcula \(\frac{z_1\cdot z_2}{3\cdot z_3}\).}
{\( -1 + \frac29\mathrm{i} \)}{\( -1 -  \frac29\mathrm{i} \)}{\( 1 + \frac29\mathrm{i} \)}{\( 1 - \frac29\mathrm{i} \)}{}{}}
\End

\Data\ID{1003067707}
\Updated{2021-04-24 05:04:26}
\Level{A}
\SubArea{Complex numbers in algebraic and trigonometric form}
\LANGen{
\Question{
Find the complex conjugate of the following complex number.
\[ \frac{2-\mathrm{i}}{2+\mathrm{i}}-\frac{2+\mathrm{i}}{2-\mathrm{i}} \]}
{\( \frac85\mathrm{i} \)}{\( -\frac85\mathrm{i} \)}{\( \frac83\mathrm{i} \)}{\( -\frac83\mathrm{i} \)}{}{}}
\LANGpl{
\Question{
Określ sprężoną liczbę zespoloną do
\[ z=\frac{2-\mathrm{i}}{2+\mathrm{i}}-\frac{2+\mathrm{i}}{2-\mathrm{i}} \]}
{\( \frac85\mathrm{i} \)}{\( -\frac85\mathrm{i} \)}{\( \frac83\mathrm{i} \)}{\( -\frac83\mathrm{i} \)}{}{}}
\LANGcs{
\Question{
Určete komplexně sdružené číslo ke komplexnímu číslu \[ z=\frac{2-\mathrm{i}}{2+\mathrm{i}}-\frac{2+\mathrm{i}}{2-\mathrm{i}}\text{.} \]}
{\( \frac85\mathrm{i} \)}{\( -\frac85\mathrm{i} \)}{\( \frac83\mathrm{i} \)}{\( -\frac83\mathrm{i} \)}{}{}}
\LANGsk{
\Question{
Určte komplexne združené číslo ku komplexnému číslu \[z= \frac{2-\mathrm{i}}{2+\mathrm{i}}-\frac{2+\mathrm{i}}{2-\mathrm{i}}\text{.} \]}
{\( \frac85\mathrm{i} \)}{\( -\frac85\mathrm{i} \)}{\( \frac83\mathrm{i} \)}{\( -\frac83\mathrm{i} \)}{}{}}
\LANGes{
\Question{
Determina el conjugado del siguiente número complejo.
\[z= \frac{2-\mathrm{i}}{2+\mathrm{i}}-\frac{2+\mathrm{i}}{2-\mathrm{i}} \]}
{\( \frac85\mathrm{i} \)}{\( -\frac85\mathrm{i} \)}{\( \frac83\mathrm{i} \)}{\( -\frac83\mathrm{i} \)}{}{}}
\End

\Data\ID{1003067708}
\Updated{2020-07-18 08:07:20}
\Level{A}
\SubArea{Complex numbers in algebraic and trigonometric form}
\LANGen{
\Question{
Let \( z=\frac{1-\mathrm{i}}{1+\mathrm{i}}+\frac{1+\mathrm{i}}{1-\mathrm{i}} \).  If \( z \) is a complex number, find its absolute value.}
{\( 0 \)}{\( 2 \)}{\( 4 \)}{Number \( z \) is not a complex number.}{}{}}
\LANGpl{
\Question{
Podano \( z=\frac{1-\mathrm{i}}{1+\mathrm{i}}+\frac{1+\mathrm{i}}{1-\mathrm{i}} \).  Jeśli \( z \) to liczba zespolona,  określ jej wartość bezwzględną.}
{\( 0 \)}{\( 2 \)}{\( 4 \)}{Liczbą \( z \) nie jest liczbą zespoloną.}{}{}}
\LANGcs{
\Question{
Nechť \( z=\frac{1 - \mathrm{i}}{1 + \mathrm{i}}+\frac{1 + \mathrm{i}}{1 - \mathrm{i}} \).  Pokud je  \( z \) komplexní číslo, určete jeho absolutní hodnotu.}
{\( 0 \)}{\( 2 \)}{\( 4 \)}{Číslo \( z \) není  komplexní číslo.}{}{}}
\LANGsk{
\Question{
Nech \( z=\frac{1 - \mathrm{i}}{1 + \mathrm{i}}+\frac{1 + \mathrm{i}}{1 - \mathrm{i}} \).  Ak je  \( z \) komplexné číslo, určte jeho absolútnu hodnotu.}
{\( 0 \)}{\( 2 \)}{\( 4 \)}{Číslo \( z \) nie je komplexné číslo.}{}{}}
\LANGes{
\Question{
Sea \( z=\frac{1-\mathrm{i}}{1+\mathrm{i}}+\frac{1+\mathrm{i}}{1-\mathrm{i}} \).  Si \( z \) es un número complejo, calcula su valor absoluto.}
{\( 0 \)}{\( 2 \)}{\( 4 \)}{El número \( z \) no es un número complejo.}{}{}}
\End

\Data\ID{1103067709}
\Updated{2020-07-18 08:07:54}
\Level{A}
\SubArea{Complex numbers in algebraic and trigonometric form}
\LANGen{
\Question{
Choose a diagram, which shows in red all the complex numbers whose absolute value is three.}
{}{}{}{}{}{}}
\LANGpl{
\Question{
Wskaż diagram, na którym kolorem czerwonym zaznaczono liczby zespolone, których wartość bezwzględna wynosi trzy.}
{}{}{}{}{}{}}
\LANGcs{
\Question{
Vyberte obrázek, na němž jsou červeně znázorněny obrazy všech komplexních čísel \( z \), pro která platí  \( |z| = 3 \).}
{}{}{}{}{}{}}
\LANGsk{
\Question{
Vyberte obrázok, na ktorom sú červeno znázornené obrazy všetkých komplexných čísel \(z\), pro ktoré platí \( |z| = 3 \).}
{}{}{}{}{}{}}
\LANGes{
\Question{
Elige un diagrama que muestra en rojo todos los números complejos cuyo valor absoluto es tres.}
{}{}{}{}{}{}}

\IMGanswer{1103067709a.pdf}
\IMGanswer{1103067709b.pdf}
\IMGanswer{1103067709c.pdf}
\IMGanswer{1103067709d.pdf}\End

\Data\ID{1103067710}
\Updated{2020-07-18 08:07:04}
\Level{A}
\SubArea{Complex numbers in algebraic and trigonometric form}
\LANGen{
\Question{
Assuming \( |z - 1| \leq 2 \), choose a diagram, which shows in red all the complex numbers \( z \).}
{}{}{}{}{}{}}
\LANGpl{
\Question{
Podano \( |z - 1| \leq 2 \), wybierz diagram, który pokazuje na czerwono wszystkie liczby zespolone \( z \).}
{}{}{}{}{}{}}
\LANGcs{
\Question{
Vyberte obrázek, na němž jsou červeně znázorněny obrazy všech komplexních čísel \( z \), pro která platí \( |z - 1| \leq 2 \).}
{}{}{}{}{}{}}
\LANGsk{
\Question{
Vyberte obrázok, na ktorom sú červeno znázornené obrazy všetkých komplexných čísel \(z\), pro ktoré platí \( |z - 1| \leq 2 \).}
{}{}{}{}{}{}}
\LANGes{
\Question{
Suponiendo que \( |z - 1| \leq 2 \), elige un diagrama  que muestra en rojo todos los números complejos  \( z \).}
{}{}{}{}{}{}}

\IMGanswer{1103067710a.pdf}
\IMGanswer{1103067710b.pdf}
\IMGanswer{1103067710c.pdf}
\IMGanswer{1103067710d.pdf}\End

\Data\ID{1103079901}
\Updated{2021-04-24 05:04:04}
\Level{A}
\SubArea{Complex numbers in algebraic and trigonometric form}
\IMG{1103079901.pdf}
\LANGen{
\Question{
The diagram shows in red all the complex numbers \( z \) such that:}
{\( |z + \mathrm{i}| \geq 2 \)}{\( |z - \mathrm{i}| \geq 2 \)}{\( |z + 1| \geq 2 \)}{\( |z - 1| \geq 2 \)}{}{}}
\LANGpl{
\Question{
Na diagramie kolorem czerwonym zaznaczono wszystkie liczby zespolone \( z \), które spełniają warunek:}
{\( |z + \mathrm{i}| \geq 2 \)}{\( |z - \mathrm{i}| \geq 2 \)}{\( |z + 1| \geq 2 \)}{\( |z - 1| \geq 2 \)}{}{}}
\LANGcs{
\Question{
Na obrázku jsou červeně označeny obrazy všech komplexních čísel  \( z \), pro která platí:}
{\( |z + \mathrm{i}| \geq 2 \)}{\( |z - \mathrm{i}| \geq 2 \)}{\( |z + 1| \geq 2 \)}{\( |z - 1| \geq 2 \)}{}{}}
\LANGsk{
\Question{
Na obrázku sú červene označené obrazy všetkých komplexných čísel \( z \), pre ktoré platí:}
{\( |z + \mathrm{i}| \geq 2 \)}{\( |z - \mathrm{i}| \geq 2 \)}{\( |z + 1| \geq 2 \)}{\( |z - 1| \geq 2 \)}{}{}}
\LANGes{
\Question{
El diagrama muestra en rojo todos los números complejos  \( z \) que cumplen que:}
{\( |z + \mathrm{i}| \geq 2 \)}{\( |z - \mathrm{i}| \geq 2 \)}{\( |z + 1| \geq 2 \)}{\( |z - 1| \geq 2 \)}{}{}}
\End

\Data\ID{1103079902}
\Updated{2021-04-24 05:04:38}
\Level{A}
\SubArea{Complex numbers in algebraic and trigonometric form}
\IMG{1103079902.pdf}
\LANGen{
\Question{
The diagram shows in red all the complex numbers \( z \) such that:}
{\( |z + 1 + 2\mathrm{i}| < 1 \)}{\( |z - 1 - 2\mathrm{i}| < 1 \)}{\( |z + 1 - 2\mathrm{i}| < 1 \)}{\( |z - 1 + 2\mathrm{i}| < 1 \)}{}{}}
\LANGpl{
\Question{
Na diagramie kolorem czerwonym zaznaczono wszystkie liczby zespolone \( z \), które spełniają warunek:}
{\( |z + 1 + 2\mathrm{i}| < 1 \)}{\( |z - 1 - 2\mathrm{i}| < 1 \)}{\( |z + 1 - 2\mathrm{i}| < 1 \)}{\( |z - 1 + 2\mathrm{i}| < 1 \)}{}{}}
\LANGcs{
\Question{
Na obrázku jsou červeně označeny obrazy všech komplexních čísel  \( z \), pro která platí:}
{\( |z + 1 + 2\mathrm{i}| < 1 \)}{\( |z - 1 - 2\mathrm{i}| < 1 \)}{\( |z + 1 - 2\mathrm{i}| < 1 \)}{\( |z - 1 + 2\mathrm{i}| < 1 \)}{}{}}
\LANGsk{
\Question{
Na obrázku sú červene označené obrazy všetkých komplexných čísel \( z \), pre ktoré platí:}
{\( |z + 1 + 2\mathrm{i}| < 1 \)}{\( |z - 1 - 2\mathrm{i}| < 1 \)}{\( |z + 1 - 2\mathrm{i}| < 1 \)}{\( |z - 1 + 2\mathrm{i}| < 1 \)}{}{}}
\LANGes{
\Question{
El diagrama muestra en rojo todos los números complejos \( z \) que cumplen que:}
{\( |z + 1 + 2\mathrm{i}| < 1 \)}{\( |z - 1 - 2\mathrm{i}| < 1 \)}{\( |z + 1 - 2\mathrm{i}| < 1 \)}{\( |z - 1 + 2\mathrm{i}| < 1 \)}{}{}}
\End

\Data\ID{1103079903}
\Updated{2021-04-24 05:04:06}
\Level{A}
\SubArea{Complex numbers in algebraic and trigonometric form}
\IMG{1103079903.pdf}
\LANGen{
\Question{
The diagram shows in red all the complex numbers \( z \) such that:}
{\( |z- 1 + \mathrm{i}| = 2 \)}{\( |z- 1 - \mathrm{i}| = 2 \)}{\( |z + 1 - \mathrm{i}| = 2 \)}{\( |z + 1 + \mathrm{i}| = 2 \)}{}{}}
\LANGpl{
\Question{
Na diagramie kolorem czerwonym zaznaczono wszystkie liczby zespolone \( z \), które spełniają równanie:}
{\( |z-1 + \mathrm{i}| = 2 \)}{\( |z -1 - \mathrm{i}| = 2 \)}{\( |z + 1 - \mathrm{i}| = 2 \)}{\( |z + 1 + \mathrm{i}| = 2 \)}{}{}}
\LANGcs{
\Question{
Na obrázku jsou červeně označeny obrazy všech komplexních čísel  \( z \), pro která platí:}
{\( |z - 1 + \mathrm{i}| = 2 \)}{\( |z - 1 - \mathrm{i}| = 2 \)}{\( |z + 1 - \mathrm{i}| = 2 \)}{\( |z + 1 + \mathrm{i}| = 2 \)}{}{}}
\LANGsk{
\Question{
Na obrázku sú červene označené obrazy všetkých komplexných čísel \( z \), pre ktoré platí:}
{\( |z - 1 + \mathrm{i}| = 2 \)}{\( |z - 1 - \mathrm{i}| = 2 \)}{\( |z + 1 - \mathrm{i}| = 2 \)}{\( |z + 1 + \mathrm{i}| = 2 \)}{}{}}
\LANGes{
\Question{
El diagrama muestra en rojo todos los números complejos  \( z \) que cumplen que:}
{\( |z- 1 + \mathrm{i}| = 2 \)}{\( |z- 1 - \mathrm{i}| = 2 \)}{\( |z + 1 - \mathrm{i}| = 2 \)}{\( |z + 1 + \mathrm{i}| = 2 \)}{}{}}
\End

\Data\ID{1103079904}
\Updated{2020-07-18 08:07:22}
\Level{A}
\SubArea{Complex numbers in algebraic and trigonometric form}
\LANGen{
\Question{
Given the complex numbers \( u = 1 + 2\mathrm{i} \) and \( v = 2 -\mathrm{i} \), choose the diagram which shows the complex number \( z \),  such that \( z = u^2  - v^2 \).}
{}{}{}{}{}{}}
\LANGpl{
\Question{
Dane są liczby zespolone \( u = 1 + 2\mathrm{i} \) i \( v = 2 -\mathrm{i} \), wybierz diagram, na którym zaznaczono liczbę zespoloną \( z \), tak, aby  \( z = u^2  - v^2 \).}
{}{}{}{}{}{}}
\LANGcs{
\Question{
Jsou daná komplexní čísla: \( u = 1 + 2\mathrm{i} \)  a  \( v = 2 -\mathrm{i} \). Vyberte obrázek, na kterém je v Gaussově rovině zobrazené komplexní číslo  \( z = u^2  - v^2 \).}
{}{}{}{}{}{}}
\LANGsk{
\Question{
Dané sú komplexné čísla: \( u = 1 + 2\mathrm{i} \)  a  \( v = 2 -\mathrm{i} \). Vyberte obrázok, na ktorom je v Gaussovej rovine zobrazené komplexné číslo  \( z = u^2  - v^2 \).}
{}{}{}{}{}{}}
\LANGes{
\Question{
Dados los números complejos \( u = 1 + 2\mathrm{i} \) y \( v = 2 -\mathrm{i} \), elige un diagrama que muestra el número complejo \( z \),  suponiendo que \( z = u^2  - v^2 \).}
{}{}{}{}{}{}}

\IMGanswer{1103079904a_0.pdf}
\IMGanswer{1103079904b_0.pdf}
\IMGanswer{1103079904c_0.pdf}
\IMGanswer{1103079904d_0.pdf}\End

\Data\ID{2010010401}
\Updated{2022-11-01 09:11:03}
\Level{A}
\SubArea{Complex numbers in algebraic and trigonometric form}
\LANGen{
\Question{
Given complex numbers \(z_{1} = 10 -7\mathrm{i}\) 
and \(z_{2} = 5- 4\mathrm{i}\),
find \(z_{1} - z_{2}\).}
{\(5 -3 \mathrm{i}\)}{\(5-11 \mathrm{i}\)}{\(3 -5 \mathrm{i}\)}{\(3 - \mathrm{i}\)}{}{}}
\LANGpl{
\Question{
Mając dane liczby zespolone \(z_{1} = 10 -7\mathrm{i}\) 
i  \(z_{2} = 5- 4\mathrm{i}\),
znajdź \(z_{1} - z_{2}\).}
{\(5 -3 \mathrm{i}\)}{\(5-11 \mathrm{i}\)}{\(3 -5 \mathrm{i}\)}{\(3 - \mathrm{i}\)}{}{}}
\LANGcs{
\Question{
Jsou dána komplexní čísla \(z_{1} = 10 -7\mathrm{i}\) 
a \(z_{2} = 5- 4\mathrm{i}\).
Vypočtěte \(z_{1} - z_{2}\).}
{\(5 -3 \mathrm{i}\)}{\(5-11 \mathrm{i}\)}{\(3 -5 \mathrm{i}\)}{\(3 - \mathrm{i}\)}{}{}}
\LANGsk{
\Question{
Sú dané komplexné čísla \(z_{1} = 10 -7\mathrm{i}\) 
a \(z_{2} = 5- 4\mathrm{i}\).
Vypočítajte \(z_{1} - z_{2}\).}
{\(5 -3 \mathrm{i}\)}{\(5-11 \mathrm{i}\)}{\(3 -5 \mathrm{i}\)}{\(3 - \mathrm{i}\)}{}{}}
\LANGes{
\Question{
Dados los números complejos \(z_{1} = 10 -7\mathrm{i}\) 
y \(z_{2} = 5- 4\mathrm{i}\),
determina \(z_{1} - z_{2}\).}
{\(5 -3 \mathrm{i}\)}{\(5-11 \mathrm{i}\)}{\(3 -5 \mathrm{i}\)}{\(3 - \mathrm{i}\)}{}{}}
\End

\Data\ID{2010010402}
\Updated{2022-08-25 10:08:41}
\Level{A}
\SubArea{Complex numbers in algebraic and trigonometric form}
\LANGen{
\Question{
Given complex numbers 
\[ 
 z_1 = \sqrt{5} - \sqrt{3}\mathrm{i},\quad z_2 = \sqrt{5} +\sqrt{3}\mathrm{i},
\]
find \(z_1 \cdot z_2\).}
{\(8\)}{\(2\)}{\(5-3 \mathrm{i}\)}{\(5+3\mathrm{i}\)}{}{}}
\LANGpl{
\Question{
Mając dane liczby zespolone  
\[ 
 z_1 = \sqrt{5} - \sqrt{3}\mathrm{i},\quad z_2 = \sqrt{5} +\sqrt{3}\mathrm{i},
\]
znajdź  \(z_1 \cdot z_2\).}
{\(8\)}{\(2\)}{\(5-3 \mathrm{i}\)}{\(5+3\mathrm{i}\)}{}{}}
\LANGcs{
\Question{
Jsou dána komplexní čísla 
\[ 
 z_1 = \sqrt{5} - \sqrt{3}\mathrm{i},\quad z_2 = \sqrt{5} +\sqrt{3}\mathrm{i}.
\]
Vypočtěte \(z_1 \cdot z_2\).}
{\(8\)}{\(2\)}{\(5-3 \mathrm{i}\)}{\(5+3\mathrm{i}\)}{}{}}
\LANGsk{
\Question{
Sú dané komplexné čísla 
\[ 
 z_1 = \sqrt{5} - \sqrt{3}\mathrm{i},\quad z_2 = \sqrt{5} +\sqrt{3}\mathrm{i}.
\]
Vypočítajte \(z_1 \cdot z_2\).}
{\(8\)}{\(2\)}{\(5-3 \mathrm{i}\)}{\(5+3\mathrm{i}\)}{}{}}
\LANGes{
\Question{
Dados los números complejos
\[ 
 z_1 = \sqrt{5} - \sqrt{3}\mathrm{i},\quad z_2 = \sqrt{5} +\sqrt{3}\mathrm{i},
\]
determina \(z_1 \cdot z_2\).}
{\(8\)}{\(2\)}{\(5-3 \mathrm{i}\)}{\(5+3\mathrm{i}\)}{}{}}
\End

\Data\ID{2010010403}
\Updated{2022-08-25 10:08:41}
\Level{A}
\SubArea{Complex numbers in algebraic and trigonometric form}
\LANGen{
\Question{
Given complex numbers 
\[ 
 a = \sqrt{2} + \mathrm{i},\ \quad b = {2} -\sqrt{3}\mathrm{i},\ 
\]
find \(\frac{a}
{b}\).}
{\(\frac{2\sqrt{2}-\sqrt{3}}
 {7} + \mathrm{i}\frac{\sqrt{6}+2} 
 {7} \)}{\(\frac{2\sqrt{2}-\sqrt{3}}
 {7} - \mathrm{i}\frac{\sqrt{6}+2} 
 {7} \)}{\(\frac{2\sqrt{2}+\sqrt{3}}
 {4} + \mathrm{i}\frac{\sqrt{6}-2} 
 {4} \)}{\(\frac{\sqrt{2}-\sqrt{6}}
 {3} + \mathrm{i}\frac{\sqrt{3}+2} 
 {3} \)}{}{}}
\LANGpl{
\Question{
Mając dane liczby zespolone 
\[ 
 a = \sqrt{2} + \mathrm{i},\ \quad b = {2} -\sqrt{3}\mathrm{i},\ 
\]
znajdź  \(\frac{a}
{b}\).}
{\(\frac{2\sqrt{2}-\sqrt{3}}
 {7} + \mathrm{i}\frac{\sqrt{6}+2} 
 {7} \)}{\(\frac{2\sqrt{2}-\sqrt{3}}
 {7} - \mathrm{i}\frac{\sqrt{6}+2} 
 {7} \)}{\(\frac{2\sqrt{2}+\sqrt{3}}
 {4} + \mathrm{i}\frac{\sqrt{6}-2} 
 {4} \)}{\(\frac{\sqrt{2}-\sqrt{6}}
 {3} + \mathrm{i}\frac{\sqrt{3}+2} 
 {3} \)}{}{}}
\LANGcs{
\Question{
Jsou dána komplexní čísla 
\[ 
 a = \sqrt{2} + \mathrm{i},\ \quad b = {2} -\sqrt{3}\mathrm{i}\text{. }
\]
Vypočtěte \(\frac{a}
{b}\).}
{\(\frac{2\sqrt{2}-\sqrt{3}}
 {7} + \mathrm{i}\frac{\sqrt{6}+2} 
 {7} \)}{\(\frac{2\sqrt{2}-\sqrt{3}}
 {7} - \mathrm{i}\frac{\sqrt{6}+2} 
 {7} \)}{\(\frac{2\sqrt{2}+\sqrt{3}}
 {4} + \mathrm{i}\frac{\sqrt{6}-2} 
 {4} \)}{\(\frac{\sqrt{2}-\sqrt{6}}
 {3} + \mathrm{i}\frac{\sqrt{3}+2} 
 {3} \)}{}{}}
\LANGsk{
\Question{
Sú dané komplexné čísla
\[ 
 a = \sqrt{2} + \mathrm{i},\ \quad b = {2} -\sqrt{3}\mathrm{i},\ 
\]
nájdite \(\frac{a}
{b}\).}
{\(\frac{2\sqrt{2}-\sqrt{3}}
 {7} + \mathrm{i}\frac{\sqrt{6}+2} 
 {7} \)}{\(\frac{2\sqrt{2}-\sqrt{3}}
 {7} - \mathrm{i}\frac{\sqrt{6}+2} 
 {7} \)}{\(\frac{2\sqrt{2}+\sqrt{3}}
 {4} + \mathrm{i}\frac{\sqrt{6}-2} 
 {4} \)}{\(\frac{\sqrt{2}-\sqrt{6}}
 {3} + \mathrm{i}\frac{\sqrt{3}+2} 
 {3} \)}{}{}}
\LANGes{
\Question{
Dados los números complejos
\[ 
 a = \sqrt{2} + \mathrm{i},\ \quad b = {2} -\sqrt{3}\mathrm{i},\ 
\]
determina \(\frac{a}
{b}\).}
{\(\frac{2\sqrt{2}-\sqrt{3}}
 {7} + \mathrm{i}\frac{\sqrt{6}+2} 
 {7} \)}{\(\frac{2\sqrt{2}-\sqrt{3}}
 {7} - \mathrm{i}\frac{\sqrt{6}+2} 
 {7} \)}{\(\frac{2\sqrt{2}+\sqrt{3}}
 {4} + \mathrm{i}\frac{\sqrt{6}-2} 
 {4} \)}{\(\frac{\sqrt{2}-\sqrt{6}}
 {3} + \mathrm{i}\frac{\sqrt{3}+2} 
 {3} \)}{}{}}
\End

\Data\ID{2010010404}
\Updated{2023-01-23 10:01:54}
\Level{A}
\SubArea{Complex numbers in algebraic and trigonometric form}
\LANGen{
\Question{
Find the complex conjugate of \(z = \mathrm{i}^{10} +3\mathrm{i}^{5}\).}
{\(-1 -3\mathrm{i}\)}{\(-1 +3\mathrm{i}\)}{\(-3-\mathrm{i}\)}{\(3-\mathrm{i}\)}{}{}}
\LANGpl{
\Question{
Znajdź liczbę sprzężoną do liczby \(z = \mathrm{i}^{10} +3\mathrm{i}^{5}\).}
{\(-1 -3\mathrm{i}\)}{\(-1 +3\mathrm{i}\)}{\(-3-\mathrm{i}\)}{\(3-\mathrm{i}\)}{}{}}
\LANGcs{
\Question{
Určete číslo komplexně sdružené k číslu \(z = \mathrm{i}^{10} +3\mathrm{i}^{5}\).}
{\(-1 -3\mathrm{i}\)}{\(-1 +3\mathrm{i}\)}{\(-3-\mathrm{i}\)}{\(3-\mathrm{i}\)}{}{}}
\LANGsk{
\Question{
Nájdite komplexne združené číslo \(z = \mathrm{i}^{10} +3\mathrm{i}^{5}\).}
{\(-1 -3\mathrm{i}\)}{\(-1 +3\mathrm{i}\)}{\(-3-\mathrm{i}\)}{\(3-\mathrm{i}\)}{}{}}
\LANGes{
\Question{
Determina el conjugado del complejo \(z = \mathrm{i}^{10} +3\mathrm{i}^{5}\).}
{\(-1 -3\mathrm{i}\)}{\(-1 +3\mathrm{i}\)}{\(-3-\mathrm{i}\)}{\(3-\mathrm{i}\)}{}{}}
\End

\Data\ID{9000031201}
\Updated{2020-07-18 08:07:28}
\Level{A}
\SubArea{Complex numbers in algebraic and trigonometric form}
\LANGen{
\Question{
Given complex numbers \(z_{1} = 1 - 2\mathrm{i}\) 
and \(z_{2} = 3 + 5\mathrm{i}\),
find \(z_{1}z_{2}\).}
{\(13 -\mathrm{i}\)}{\(13 + \mathrm{i}\)}{\(- 7 -\mathrm{i}\)}{\(13 + 11\mathrm{i}\)}{}{}}
\LANGpl{
\Question{
Podano liczby zespolone \(z_{1} = 1 - 2\mathrm{i}\) 
i  \(z_{2} = 3 + 5\mathrm{i}\),
wskaż  \(z_{1}z_{2}\).}
{\(13 -\mathrm{i}\)}{\(13 + \mathrm{i}\)}{\(- 7 -\mathrm{i}\)}{\(13 + 11\mathrm{i}\)}{}{}}
\LANGcs{
\Question{
Jsou dána komplexní čísla \(z_{1} = 1 - 2\mathrm{i}\),
\(z_{2} = 3 + 5\mathrm{i}\).
Určete jejich součin v algebraickém tvaru.}
{\(13 -\mathrm{i}\)}{\(13 + \mathrm{i}\)}{\(- 7 -\mathrm{i}\)}{\(13 + 11\mathrm{i}\)}{}{}}
\LANGsk{
\Question{
Sú dané komplexné čísla \(z_{1} = 1 - 2\mathrm{i}\) 
a \(z_{2} = 3 + 5\mathrm{i}\).
Určte súčin \(z_{1}z_{2}\) v algebraickom tvare.}
{\(13 -\mathrm{i}\)}{\(13 + \mathrm{i}\)}{\(- 7 -\mathrm{i}\)}{\(13 + 11\mathrm{i}\)}{}{}}
\LANGes{
\Question{
Dados los números complejos \(z_{1} = 1 - 2\mathrm{i}\) 
y \(z_{2} = 3 + 5\mathrm{i}\),
calcula \(z_{1}z_{2}\).}
{\(13 -\mathrm{i}\)}{\(13 + \mathrm{i}\)}{\(- 7 -\mathrm{i}\)}{\(13 + 11\mathrm{i}\)}{}{}}
\End

\Data\ID{9000031202}
\Updated{2022-04-23 05:04:05}
\Level{A}
\SubArea{Complex numbers in algebraic and trigonometric form}
\LANGen{
\Question{
Find the imaginary part of the complex number
\(z = \frac{2+\mathrm{i}} 
{1-\mathrm{i}}\).}
{\(\frac{3}
{2}\)}{\(-\frac{3} 
{2}\)}{\(\frac{1}
{2}\)}{\(-\frac{1} 
{2}\)}{}{}}
\LANGpl{
\Question{
Wskaż część urojoną liczby zespolonej
\(z = \frac{2+\mathrm{i}} 
{1-\mathrm{i}}\).}
{\(\frac{3}
{2}\)}{\(-\frac{3} 
{2}\)}{\(\frac{1}
{2}\)}{\(-\frac{1} 
{2}\)}{}{}}
\LANGcs{
\Question{
Určete imaginární část komplexního čísla
\(z = \frac{2+\mathrm{i}} 
{1-\mathrm{i}}\).}
{\(\frac{3}
{2}\)}{\(-\frac{3} 
{2}\)}{\(\frac{1}
{2}\)}{\(-\frac{1} 
{2}\)}{}{}}
\LANGsk{
\Question{
Určte imaginárnu časť komplexného čísla 
\(z = \frac{2+\mathrm{i}} 
{1-\mathrm{i}}\).}
{\(\frac{3}
{2}\)}{\(-\frac{3} 
{2}\)}{\(\frac{1}
{2}\)}{\(-\frac{1} 
{2}\)}{}{}}
\LANGes{
\Question{
Determina la unidad imaginaria del número complejo \(z = \frac{2+\mathrm{i}} 
{1-\mathrm{i}}\).}
{\(\frac{3}
{2}\)}{\(-\frac{3} 
{2}\)}{\(\frac{1}
{2}\)}{\(-\frac{1} 
{2}\)}{}{}}
\End

\Data\ID{9000031203}
\Updated{2021-12-19 10:12:48}
\Level{A}
\SubArea{Complex numbers in algebraic and trigonometric form}
\LANGen{
\Question{
Find the real part of the complex number
\(z = \frac{2-\mathrm{i}} 
{2+\mathrm{i}}\).}
{\(0.6\)}{\(0.8\)}{\(- 0.8\)}{\(1\)}{}{}}
\LANGpl{
\Question{
Wskaż część rzeczywistą  liczby zespolonej
\(z = \frac{2-\mathrm{i}} 
{2+\mathrm{i}}\).}
{\(0{,}6\)}{\(0{,}8\)}{\(- 0{,}8\)}{\(1\)}{}{}}
\LANGcs{
\Question{
Určete reálnou část komplexního čísla
\(z = \frac{2-\mathrm{i}} 
{2+\mathrm{i}}\).}
{\(0{,}6\)}{\(0{,}8\)}{\(- 0{,}8\)}{\(1\)}{}{}}
\LANGsk{
\Question{
Určte reálnu časť komplexného čísla
\(z = \frac{2-\mathrm{i}} 
{2+\mathrm{i}}\).}
{\(0{,}6\)}{\(0{,}8\)}{\(- 0{,}8\)}{\(1\)}{}{}}
\LANGes{
\Question{
Determina la parte real del número complejo 
\(z = \frac{2-\mathrm{i}} 
{2+\mathrm{i}}\).}
{\(0.6\)}{\(0.8\)}{\(- 0.8\)}{\(1\)}{}{}}
\End

\Data\ID{9000031204}
\Updated{2022-04-23 05:04:05}
\Level{A}
\SubArea{Complex numbers in algebraic and trigonometric form}
\LANGen{
\Question{
Find the absolute value of the complex number
\(z = \frac{2-\mathrm{i}} 
{2+\mathrm{i}}\).}
{\(1\)}{\(5\)}{\(\frac{\sqrt{7}}
 {5} \)}{\(\frac{\sqrt{5}}
 {5} \)}{}{}}
\LANGpl{
\Question{
Określ wartość bezwzględną podanej liczby zespolonej
\(z = \frac{2-\mathrm{i}} 
{2+\mathrm{i}}\).}
{\(1\)}{\(5\)}{\(\frac{\sqrt{7}}
 {5} \)}{\(\frac{\sqrt{5}}
 {5} \)}{}{}}
\LANGcs{
\Question{
Určete absolutní hodnotu komplexního čísla
\(z = \frac{2-\mathrm{i}} 
{2+\mathrm{i}}\).}
{\(1\)}{\(5\)}{\(\frac{\sqrt{7}}
 {5} \)}{\(\frac{\sqrt{5}}
 {5} \)}{}{}}
\LANGsk{
\Question{
Určte absolútnu hodnotu komplexného čísla
\(z = \frac{2-\mathrm{i}} 
{2+\mathrm{i}}\).}
{\(1\)}{\(5\)}{\(\frac{\sqrt{7}}
 {5} \)}{\(\frac{\sqrt{5}}
 {5} \)}{}{}}
\LANGes{
\Question{
Determina el valor absoluto del número complejo
\(z = \frac{2-\mathrm{i}} 
{2+\mathrm{i}}\).}
{\(1\)}{\(5\)}{\(\frac{\sqrt{7}}
 {5} \)}{\(\frac{\sqrt{5}}
 {5} \)}{}{}}
\End

\Data\ID{9000031205}
\Updated{2022-04-23 05:04:48}
\Level{A}
\SubArea{Complex numbers in algebraic and trigonometric form}
\LANGen{
\Question{
Find the complex conjugate of \(z = \mathrm{i}^{5} - 3\mathrm{i}^{10}\).}
{\(3 -\mathrm{i}\)}{\(- 3 -\mathrm{i}\)}{\(- 3 + \mathrm{i}\)}{\(3 + \mathrm{i}\)}{}{}}
\LANGpl{
\Question{
Określ sprężoną liczbę zespoloną do  \(z = \mathrm{i}^{5} - 3\mathrm{i}^{10}\).}
{\(3 -\mathrm{i}\)}{\(- 3 -\mathrm{i}\)}{\(- 3 + \mathrm{i}\)}{\(3 + \mathrm{i}\)}{}{}}
\LANGcs{
\Question{
Určete číslo komplexně sdružené k číslu
\(z = \mathrm{i}^{5} - 3\mathrm{i}^{10}\).}
{\(3 -\mathrm{i}\)}{\(- 3 -\mathrm{i}\)}{\(- 3 + \mathrm{i}\)}{\(3 + \mathrm{i}\)}{}{}}
\LANGsk{
\Question{
Určte číslo komplexne združené k číslu \(z = \mathrm{i}^{5} - 3\mathrm{i}^{10}\).}
{\(3 -\mathrm{i}\)}{\(- 3 -\mathrm{i}\)}{\(- 3 + \mathrm{i}\)}{\(3 + \mathrm{i}\)}{}{}}
\LANGes{
\Question{
Determina el conjugado del complejo \(z = \mathrm{i}^{5} - 3\mathrm{i}^{10}\).}
{\(3 -\mathrm{i}\)}{\(- 3 -\mathrm{i}\)}{\(- 3 + \mathrm{i}\)}{\(3 + \mathrm{i}\)}{}{}}
\End

\Data\ID{9000031206}
\Updated{2021-04-24 05:04:16}
\Level{A}
\SubArea{Complex numbers in algebraic and trigonometric form}
\LANGen{
\Question{
Find the opposite number to the complex number
\(z = \frac{1+\mathrm{i}} 
{1-\mathrm{i}}\).}
{\(-\mathrm{i}\)}{\(1\)}{\(- 1\)}{\(\mathrm{i}\)}{}{}}
\LANGpl{
\Question{
Określ liczbę przeciwną do liczby zespolonej
\(z = \frac{1+\mathrm{i}} 
{1-\mathrm{i}}\).}
{\(-\mathrm{i}\)}{\(1\)}{\(- 1\)}{\(\mathrm{i}\)}{}{}}
\LANGcs{
\Question{
Určete číslo opačné ke komplexnímu číslu
\(z = \frac{1+\mathrm{i}} 
{1-\mathrm{i}}\).}
{\(-\mathrm{i}\)}{\(1\)}{\(- 1\)}{\(\mathrm{i}\)}{}{}}
\LANGsk{
\Question{
Určte opačné číslo ku komplexnému číslu 
\(z = \frac{1+\mathrm{i}} 
{1-\mathrm{i}}\).}
{\(-\mathrm{i}\)}{\(1\)}{\(- 1\)}{\(\mathrm{i}\)}{}{}}
\LANGes{
\Question{
Determina el opuesto del número complejo
\(z = \frac{1+\mathrm{i}} 
{1-\mathrm{i}}\).}
{\(-\mathrm{i}\)}{\(1\)}{\(- 1\)}{\(\mathrm{i}\)}{}{}}
\End

\Data\ID{9000034801}
\Updated{2022-04-23 05:04:48}
\Level{A}
\SubArea{Complex numbers in algebraic and trigonometric form}
\LANGen{
\Question{
Given complex numbers \(z_{1} = 4 -\mathrm{i}\) 
and \(z_{2} = 1 - 2\mathrm{i}\),
find \(z_{1} - z_{2}\).}
{\(3 + \mathrm{i}\)}{\(3 - 3\mathrm{i}\)}{\(5 - 3\mathrm{i}\)}{\(3 -\mathrm{i}\)}{}{}}
\LANGpl{
\Question{
Dane są liczby zespolone  \(z_{1} = 4 -\mathrm{i}\) 
i \(z_{2} = 1 - 2\mathrm{i}\),
określ  \(z_{1} - z_{2}\).}
{\(3 + \mathrm{i}\)}{\(3 - 3\mathrm{i}\)}{\(5 - 3\mathrm{i}\)}{\(3 -\mathrm{i}\)}{}{}}
\LANGcs{
\Question{
Jsou dána komplexní čísla \(z_{1} = 4 -\mathrm{i}\),
\(z_{2} = 1 - 2\mathrm{i}\). Určete
jejich rozdíl \(z_{1} - z_{2}\) 
v algebraickém tvaru.}
{\(3 + \mathrm{i}\)}{\(3 - 3\mathrm{i}\)}{\(5 - 3\mathrm{i}\)}{\(3 -\mathrm{i}\)}{}{}}
\LANGsk{
\Question{
Sú dané komplexné čísla \(z_{1} = 4 -\mathrm{i}\) 
a \(z_{2} = 1 - 2\mathrm{i}\). Určte ich rozdiel \(z_{1} - z_{2}\) v algebraickom tvare.}
{\(3 + \mathrm{i}\)}{\(3 - 3\mathrm{i}\)}{\(5 - 3\mathrm{i}\)}{\(3 -\mathrm{i}\)}{}{}}
\LANGes{
\Question{
Dados los números complejos \(z_{1} = 4 -\mathrm{i}\) 
y \(z_{2} = 1 - 2\mathrm{i}\),
calcula \(z_{1} - z_{2}\).}
{\(3 + \mathrm{i}\)}{\(3 - 3\mathrm{i}\)}{\(5 - 3\mathrm{i}\)}{\(3 -\mathrm{i}\)}{}{}}
\End

\Data\ID{9000034802}
\Updated{2020-07-18 09:07:32}
\Level{A}
\SubArea{Complex numbers in algebraic and trigonometric form}
\LANGen{
\Question{
Find the opposite number to the complex number
\(z = 3 -\mathrm{i}\).}
{\(- 3 + \mathrm{i}\)}{\(- 3 -\mathrm{i}\)}{\(3 + \mathrm{i}\)}{\(3 -\mathrm{i}\)}{}{}}
\LANGpl{
\Question{
Określ liczbę przeciwną do liczby zespolonej
\(z = 3 -\mathrm{i}\).}
{\(- 3 + \mathrm{i}\)}{\(- 3 -\mathrm{i}\)}{\(3 + \mathrm{i}\)}{\(3 -\mathrm{i}\)}{}{}}
\LANGcs{
\Question{
Určete číslo opačné ke komplexnímu číslu
\(z = 3 -\mathrm{i}\).}
{\(- 3 + \mathrm{i}\)}{\(- 3 -\mathrm{i}\)}{\(3 + \mathrm{i}\)}{\(3 -\mathrm{i}\)}{}{}}
\LANGsk{
\Question{
Určte opačné číslo ku komplexnému číslu 
\(z = 3 -\mathrm{i}\).}
{\(- 3 + \mathrm{i}\)}{\(- 3 -\mathrm{i}\)}{\(3 + \mathrm{i}\)}{\(3 -\mathrm{i}\)}{}{}}
\LANGes{
\Question{
Determina el opuesto del número complejo 
\(z = 3 -\mathrm{i}\).}
{\(- 3 + \mathrm{i}\)}{\(- 3 -\mathrm{i}\)}{\(3 + \mathrm{i}\)}{\(3 -\mathrm{i}\)}{}{}}
\End

\Data\ID{9000034803}
\Updated{2020-07-18 09:07:00}
\Level{A}
\SubArea{Complex numbers in algebraic and trigonometric form}
\LANGen{
\Question{
Find the complex conjugate of \(z = 1 - 3\mathrm{i}\).}
{\(1 + 3\mathrm{i}\)}{\(- 1 - 3\mathrm{i}\)}{\(- 1 + 3\mathrm{i}\)}{\(1 - 3\mathrm{i}\)}{}{}}
\LANGpl{
\Question{
Określ sprężoną liczbę zespoloną do  \(z = 1 - 3\mathrm{i}\).}
{\(1 + 3\mathrm{i}\)}{\(- 1 - 3\mathrm{i}\)}{\(- 1 + 3\mathrm{i}\)}{\(1 - 3\mathrm{i}\)}{}{}}
\LANGcs{
\Question{
Určete číslo komplexně sdružené k číslu
\(z = 1 - 3\mathrm{i}\).}
{\(1 + 3\mathrm{i}\)}{\(- 1 - 3\mathrm{i}\)}{\(- 1 + 3\mathrm{i}\)}{\(1 - 3\mathrm{i}\)}{}{}}
\LANGsk{
\Question{
Určte číslo komplexne združené k číslu \(z = 1 - 3\mathrm{i}\).}
{\(1 + 3\mathrm{i}\)}{\(- 1 - 3\mathrm{i}\)}{\(- 1 + 3\mathrm{i}\)}{\(1 - 3\mathrm{i}\)}{}{}}
\LANGes{
\Question{
Determina el conjugado de un complejo \(z = 1 - 3\mathrm{i}\).}
{\(1 + 3\mathrm{i}\)}{\(- 1 - 3\mathrm{i}\)}{\(- 1 + 3\mathrm{i}\)}{\(1 - 3\mathrm{i}\)}{}{}}
\End

\Data\ID{9000034804}
\Updated{2021-04-24 05:04:04}
\Level{A}
\SubArea{Complex numbers in algebraic and trigonometric form}
\LANGen{
\Question{
Find the absolute value of the complex number
\(z = 3 -\mathrm{i}\).}
{\(\sqrt{10}\)}{\(2\)}{\(2\sqrt{2}\)}{\(\sqrt{2}\)}{}{}}
\LANGpl{
\Question{
Określ wartość bezwzględną podanej liczby zespolonej
\(z = 3 -\mathrm{i}\).}
{\(\sqrt{10}\)}{\(2\)}{\(2\sqrt{2}\)}{\(\sqrt{2}\)}{}{}}
\LANGcs{
\Question{
Určete absolutní hodnotu komplexního čísla
\(z = 3 -\mathrm{i}\).}
{\(\sqrt{10}\)}{\(2\)}{\(2\sqrt{2}\)}{\(\sqrt{2}\)}{}{}}
\LANGsk{
\Question{
Určte absolútnu hodnotu komplexného čísla
\(z = 3 -\mathrm{i}\).}
{\(\sqrt{10}\)}{\(2\)}{\(2\sqrt{2}\)}{\(\sqrt{2}\)}{}{}}
\LANGes{
\Question{
Calcula el valor absoluto del número complejo \(z = 3 -\mathrm{i}\).}
{\(\sqrt{10}\)}{\(2\)}{\(2\sqrt{2}\)}{\(\sqrt{2}\)}{}{}}
\End

\Data\ID{9000034805}
\Updated{2020-07-18 10:07:50}
\Level{A}
\SubArea{Complex numbers in algebraic and trigonometric form}
\LANGen{
\Question{
Find the complex number \(z\) 
which satisfies \(2z = 2 - 3\mathrm{i}\).}
{\(1 -\frac{3} 
{2}\mathrm{i}\)}{\(- 3\mathrm{i}\)}{\(4 - 6\mathrm{i}\)}{\(- 1 + \frac{3} 
{2}\mathrm{i}\)}{}{}}
\LANGpl{
\Question{
Określ  liczbę zespoloną  \(z\) 
spełniającą równanie \(2z = 2 - 3\mathrm{i}\).}
{\(1 -\frac{3} 
{2}\mathrm{i}\)}{\(- 3\mathrm{i}\)}{\(4 - 6\mathrm{i}\)}{\(- 1 + \frac{3} 
{2}\mathrm{i}\)}{}{}}
\LANGcs{
\Question{
Najděte komplexní číslo \(z\) 
v algebraickém tvaru, jestliže platí
\(2z = 2 - 3\mathrm{i}\).}
{\(1 -\frac{3} 
{2}\mathrm{i}\)}{\(- 3\mathrm{i}\)}{\(4 - 6\mathrm{i}\)}{\(- 1 + \frac{3} 
{2}\mathrm{i}\)}{}{}}
\LANGsk{
\Question{
Nájdite komplexné číslo \(z\) v algebraickom tvare, ak platí
 \(2z = 2 - 3\mathrm{i}\).}
{\(1 -\frac{3} 
{2}\mathrm{i}\)}{\(- 3\mathrm{i}\)}{\(4 - 6\mathrm{i}\)}{\(- 1 + \frac{3} 
{2}\mathrm{i}\)}{}{}}
\LANGes{
\Question{
Halla el número complejo \(z\) 
suponiendo que  \(2z = 2 - 3\mathrm{i}\).}
{\(1 -\frac{3} 
{2}\mathrm{i}\)}{\(- 3\mathrm{i}\)}{\(4 - 6\mathrm{i}\)}{\(- 1 + \frac{3} 
{2}\mathrm{i}\)}{}{}}
\End

\Data\ID{9000034806}
\Updated{2020-07-18 10:07:09}
\Level{A}
\SubArea{Complex numbers in algebraic and trigonometric form}
\LANGen{
\Question{
Simplify \(\mathrm{i}^{15}\).}
{\(-\mathrm{i}\)}{\(- 1\)}{\(1\)}{\(\mathrm{i}\)}{}{}}
\LANGpl{
\Question{
Uprość \(\mathrm{i}^{15}\).}
{\(-\mathrm{i}\)}{\(- 1\)}{\(1\)}{\(\mathrm{i}\)}{}{}}
\LANGcs{
\Question{
Vypočítejte \(\mathrm{i}^{15}\).}
{\(-\mathrm{i}\)}{\(- 1\)}{\(1\)}{\(\mathrm{i}\)}{}{}}
\LANGsk{
\Question{
Vypočítajte \(\mathrm{i}^{15}\).}
{\(-\mathrm{i}\)}{\(- 1\)}{\(1\)}{\(\mathrm{i}\)}{}{}}
\LANGes{
\Question{
Simplifica \(\mathrm{i}^{15}\).}
{\(-\mathrm{i}\)}{\(- 1\)}{\(1\)}{\(\mathrm{i}\)}{}{}}
\End

\Data\ID{9000035701}
\Updated{2020-07-18 11:07:35}
\Level{A}
\SubArea{Complex numbers in algebraic and trigonometric form}
\IMG{000357_01-figure0.pdf}
\LANGen{
\Question{
What is the algebraic form of the complex number \( A \) graphed in the complex
plane (as shown in the picture)?}
{\( -3 + 2\mathrm{i}\)}{\( 2 - 3\mathrm{i}\)}{\( 2 + 3\mathrm{i}\)}{\( -3 - 2\mathrm{i}\)}{}{}}
\LANGpl{
\Question{
Wskaż postać algebraiczną liczby zespolonej zaznaczonej  na płaszczyźnie zespolonej.}
{\(-3 + 2\mathrm{i}\)}{\(2 - 3\mathrm{i}\)}{\( 2 + 3\mathrm{i}\)}{\( -3 - 2\mathrm{i}\)}{}{}}
\LANGcs{
\Question{
Bod \( A \) (viz obrázek) je obrazem komplexního čísla:}
{\(-3 + 2\mathrm{i}\)}{\( 2 - 3\mathrm{i}\)}{\( 2 + 3\mathrm{i}\)}{\(-3 - 2\mathrm{i}\)}{}{}}
\LANGsk{
\Question{
Bod \( A \) (viz obrázok) je obrazom komplexného čísla:}
{\(-3 + 2\mathrm{i}\)}{\( 2 - 3\mathrm{i}\)}{\(2 + 3\mathrm{i}\)}{\( -3 - 2\mathrm{i}\)}{}{}}
\LANGes{
\Question{
El punto \(A \) (mira la imagen) es una representación de un número complejo:}
{\( -3 + 2\mathrm{i}\)}{\( 2 - 3\mathrm{i}\)}{\( 2 + 3\mathrm{i}\)}{\( -3 - 2\mathrm{i}\)}{}{}}
\End

\Data\ID{9000035702}
\Updated{2020-07-18 11:07:12}
\Level{A}
\SubArea{Complex numbers in algebraic and trigonometric form}
\IMG{000357_02-figure0.pdf}
\LANGen{
\Question{
Find the absolute value of the complex number \( A \) graphed in the complex
plane as shown in the picture.}
{\(5\)}{\(\sqrt{5}\)}{\(3\)}{\(4\)}{}{}}
\LANGpl{
\Question{
Wskaż wartość bezwzględną  liczby zespolonej zaznaczonej  na płaszczyźnie zespolonej.}
{\(5\)}{\(\sqrt{5}\)}{\(3\)}{\(4\)}{}{}}
\LANGcs{
\Question{
Jaká je absolutní hodnota komplexního čísla znázorněného v komplexní rovině bodem \(A\) (viz obrázek)?}
{\(5\)}{\(\sqrt{5}\)}{\(3\)}{\(4\)}{}{}}
\LANGsk{
\Question{
Absolútna hodnota čísla znázorneného bodom A sa rovná:}
{\(5\)}{\(\sqrt{5}\)}{\(3\)}{\(4\)}{}{}}
\LANGes{
\Question{
¿Cuál es el valor absoluto de un número complejo representado en el plano complejo por el punto \(A \) (mira la imagen)?}
{\(5\)}{\(\sqrt{5}\)}{\(3\)}{\(4\)}{}{}}
\End

\Data\ID{9000035703}
\Updated{2020-07-18 11:07:51}
\Level{A}
\SubArea{Complex numbers in algebraic and trigonometric form}
\IMG{000357_03-figure0.pdf}
\LANGen{
\Question{
Find the absolute value of the complex number \( A \) graphed in the complex
plane as shown in the picture.}
{\(2\sqrt{5}\)}{\(2\sqrt{3}\)}{\(4\)}{\(\sqrt{6}\)}{}{}}
\LANGpl{
\Question{
Wskaż wartość bezwzględną  liczby zespolonej zaznaczonej  na płaszczyźnie zespolonej.}
{\(2\sqrt{5}\)}{\(2\sqrt{3}\)}{\(4\)}{\(\sqrt{6}\)}{}{}}
\LANGcs{
\Question{
Jaká je absolutní hodnota komplexního čísla znázorněného v komplexní rovině bodem \(A\) (viz obrázek)?}
{\(2\sqrt{5}\)}{\(2\sqrt{3}\)}{\(4\)}{\(\sqrt{6}\)}{}{}}
\LANGsk{
\Question{
Absolútna hodnota čísla znázorneného bodom A sa rovná:}
{\(2\sqrt{5}\)}{\(2\sqrt{3}\)}{\(4\)}{\(\sqrt{6}\)}{}{}}
\LANGes{
\Question{
¿Cuál es el valor absoluto de un número complejo representado en el plano complejo por el punto \(A \) (mira la imagen)?}
{\(2\sqrt{5}\)}{\(2\sqrt{3}\)}{\(4\)}{\(\sqrt{6}\)}{}{}}
\End

\Data\ID{9000035705}
\Updated{2020-07-18 11:07:02}
\Level{A}
\SubArea{Complex numbers in algebraic and trigonometric form}
\LANGen{
\Question{
Find the absolute value of the complex number
\(z = (1 - 2\mathrm{i})(2 + \mathrm{i})\).}
{\(5\)}{\(3\)}{\(\sqrt{10}\)}{\(2\sqrt{2}\)}{}{}}
\LANGpl{
\Question{
Określ wartość bezwzględną podanej liczby zespolonej
\(z = (1 - 2\mathrm{i})(2 + \mathrm{i})\).}
{\(5\)}{\(3\)}{\(\sqrt{10}\)}{\(2\sqrt{2}\)}{}{}}
\LANGcs{
\Question{
Absolutní hodnota komplexního čísla
\(z = (1 - 2\mathrm{i})(2 + \mathrm{i})\) je
rovna:}
{\(5\)}{\(3\)}{\(\sqrt{10}\)}{\(2\sqrt{2}\)}{}{}}
\LANGsk{
\Question{
Absolútna hodnota komplexného čísla
\(z = (1 - 2\mathrm{i})(2 + \mathrm{i})\) sa rovná:}
{\(5\)}{\(3\)}{\(\sqrt{10}\)}{\(2\sqrt{2}\)}{}{}}
\LANGes{
\Question{
Determina el valor absoluto del número complejo
\(z = (1 - 2\mathrm{i})(2 + \mathrm{i})\).}
{\(5\)}{\(3\)}{\(\sqrt{10}\)}{\(2\sqrt{2}\)}{}{}}
\End

\Data\ID{9000035706}
\Updated{2021-12-19 10:12:04}
\Level{A}
\SubArea{Complex numbers in algebraic and trigonometric form}
\LANGen{
\Question{
Find the absolute value of the complex number
\(z = \frac{2+6\mathrm{i}} 
{1-2\mathrm{i}}\).}
{\(2\sqrt{2}\)}{\(2\sqrt{5}\)}{\(2\)}{\(2\sqrt{3}\)}{}{}}
\LANGpl{
\Question{
Określ wartość bezwzględną podanej liczby zespolonej
\(z = \frac{2+6\mathrm{i}} 
{1-2\mathrm{i}}\).}
{\(2\sqrt{2}\)}{\(2\sqrt{5}\)}{\(2\)}{\(2\sqrt{3}\)}{}{}}
\LANGcs{
\Question{
Absolutní hodnota komplexního čísla
\(z = \frac{2+6\mathrm{i}} 
{1-2\mathrm{i}}\) je
rovna:}
{\(2\sqrt{2}\)}{\(2\sqrt{5}\)}{\(2\)}{\(2\sqrt{3}\)}{}{}}
\LANGsk{
\Question{
Absolútna hodnota komplexného čísla 
\(z = \frac{2+6\mathrm{i}} 
{1-2\mathrm{i}}\) sa rovná:}
{\(2\sqrt{2}\)}{\(2\sqrt{5}\)}{\(2\)}{\(2\sqrt{3}\)}{}{}}
\LANGes{
\Question{
Determina el valor absoluto del número complejo
\(z = \frac{2+6\mathrm{i}} 
{1-2\mathrm{i}}\).}
{\(2\sqrt{2}\)}{\(2\sqrt{5}\)}{\(2\)}{\(2\sqrt{3}\)}{}{}}
\End

\Data\ID{9000035707}
\Updated{2022-04-23 05:04:05}
\Level{A}
\SubArea{Complex numbers in algebraic and trigonometric form}
\LANGen{
\Question{
Find the real part of the complex number
\(z= 2 + 2\mathrm{i}^{2} + \mathrm{i}^{3} -\mathrm{i}^{4} + 2\mathrm{i}^{5}\).}
{\(- 1\)}{\(1\)}{\(5\)}{\(- 3\)}{}{}}
\LANGpl{
\Question{
Wskaż część rzeczywistą podanej  liczby zespolonej
\(z=2 + 2\mathrm{i}^{2} + \mathrm{i}^{3} -\mathrm{i}^{4} + 2\mathrm{i}^{5}\).}
{\(- 1\)}{\(1\)}{\(5\)}{\(- 3\)}{}{}}
\LANGcs{
\Question{
Reálná část komplexního čísla
\(z=2 + 2\mathrm{i}^{2} + \mathrm{i}^{3} -\mathrm{i}^{4} + 2\mathrm{i}^{5}\) se
rovná:}
{\(- 1\)}{\(1\)}{\(5\)}{\(- 3\)}{}{}}
\LANGsk{
\Question{
Reálna časť komplexného čísla
\(z=2 + 2\mathrm{i}^{2} + \mathrm{i}^{3} -\mathrm{i}^{4} + 2\mathrm{i}^{5}\) sa rovná:}
{\(- 1\)}{\(1\)}{\(5\)}{\(- 3\)}{}{}}
\LANGes{
\Question{
Determina la parte real del número coplejo
\(z= 2 + 2\mathrm{i}^{2} + \mathrm{i}^{3} -\mathrm{i}^{4} + 2\mathrm{i}^{5}\).}
{\(- 1\)}{\(1\)}{\(5\)}{\(- 3\)}{}{}}
\End

\Data\ID{9000035708}
\Updated{2021-12-19 10:12:36}
\Level{A}
\SubArea{Complex numbers in algebraic and trigonometric form}
\LANGen{
\Question{
Find the imaginary part of the complex number
\(z=1 + 2\mathrm{i}^{12} + 3\mathrm{i}^{19} -\mathrm{i}^{22} + 2\mathrm{i}^{105}\).}
{\(- 1\)}{\(- 5\)}{\(1\)}{\(4\)}{}{}}
\LANGpl{
\Question{
Wskaż część urojoną  podanej  liczby zespolonej
\(z=1 + 2\mathrm{i}^{12} + 3\mathrm{i}^{19} -\mathrm{i}^{22} + 2\mathrm{i}^{105}\).}
{\(- 1\)}{\(- 5\)}{\(1\)}{\(4\)}{}{}}
\LANGcs{
\Question{
Imaginární část komplexního čísla
\(z=1 + 2\mathrm{i}^{12} + 3\mathrm{i}^{19} -\mathrm{i}^{22} + 2\mathrm{i}^{105}\) se
rovná:}
{\(- 1\)}{\(- 5\)}{\(1\)}{\(4\)}{}{}}
\LANGsk{
\Question{
Imaginárna časť komplexného čísla
\(z=1 + 2\mathrm{i}^{12} + 3\mathrm{i}^{19} -\mathrm{i}^{22} + 2\mathrm{i}^{105}\) sa rovná:}
{\(- 1\)}{\(- 5\)}{\(1\)}{\(4\)}{}{}}
\LANGes{
\Question{
Determina la parte imaginaria del número complejo 
\(z=1 + 2\mathrm{i}^{12} + 3\mathrm{i}^{19} -\mathrm{i}^{22} + 2\mathrm{i}^{105}\).}
{\(- 1\)}{\(- 5\)}{\(1\)}{\(4\)}{}{}}
\End

\Data\ID{9000035710}
\Updated{2021-04-24 05:04:37}
\Level{A}
\SubArea{Complex numbers in algebraic and trigonometric form}
\LANGen{
\Question{
Find the complex conjugate of \(z=\frac{3+\mathrm{i}}
{2-\mathrm{i}} + (\mathrm{i} + 1)(2 + \mathrm{i})\).}
{\(2 - 4\mathrm{i}\)}{\(2 + 4\mathrm{i}\)}{\(- 2 - 4\mathrm{i}\)}{\(- 2 + 4\mathrm{i}\)}{}{}}
\LANGpl{
\Question{
Określ sprężoną liczbę zespoloną do  \(z=\frac{3+\mathrm{i}}
{2-\mathrm{i}} + (\mathrm{i} + 1)(2 + \mathrm{i})\).}
{\(2 - 4\mathrm{i}\)}{\(2 + 4\mathrm{i}\)}{\(- 2 - 4\mathrm{i}\)}{\(- 2 + 4\mathrm{i}\)}{}{}}
\LANGcs{
\Question{
Číslo komplexně sdružené s číslem
\(z=\frac{3+\mathrm{i}}
{2-\mathrm{i}} + (\mathrm{i} + 1)(2 + \mathrm{i})\) má
tvar:}
{\(2 - 4\mathrm{i}\)}{\(2 + 4\mathrm{i}\)}{\(- 2 - 4\mathrm{i}\)}{\(- 2 + 4\mathrm{i}\)}{}{}}
\LANGsk{
\Question{
Číslo komplexne združené s číslom \(z=\frac{3+\mathrm{i}}
{2-\mathrm{i}} + (\mathrm{i} + 1)(2 + \mathrm{i})\) má tvar:}
{\(2 - 4\mathrm{i}\)}{\(2 + 4\mathrm{i}\)}{\(- 2 - 4\mathrm{i}\)}{\(- 2 + 4\mathrm{i}\)}{}{}}
\LANGes{
\Question{
Determina el conjugado del número complejo \(z=\frac{3+\mathrm{i}}
{2-\mathrm{i}} + (\mathrm{i} + 1)(2 + \mathrm{i})\).}
{\(2 - 4\mathrm{i}\)}{\(2 + 4\mathrm{i}\)}{\(- 2 - 4\mathrm{i}\)}{\(- 2 + 4\mathrm{i}\)}{}{}}
\End

\Data\ID{9000035801}
\Updated{2020-07-18 11:07:56}
\Level{A}
\SubArea{Complex numbers in algebraic and trigonometric form}
\LANGen{
\Question{
Find the complex conjugate of the following complex number. 
\[ 
 \mathrm{i} + 3\mathrm{i}(2 -\mathrm{i})^{2} - 4(1 -\mathrm{i})^{3}
\]}
{\(20 - 18\mathrm{i}\)}{\(20 - 24\mathrm{i}\)}{\(20 + 18\mathrm{i}\)}{\(- 8 + 26\mathrm{i}\)}{}{}}
\LANGpl{
\Question{
Określ sprężoną liczbę zespoloną do podanej liczby zespolonej.
\[ 
 \mathrm{i} + 3\mathrm{i}(2 -\mathrm{i})^{2} - 4(1 -\mathrm{i})^{3}
\]}
{\(20 - 18\mathrm{i}\)}{\(20 - 24\mathrm{i}\)}{\(20 + 18\mathrm{i}\)}{\(- 8 + 26\mathrm{i}\)}{}{}}
\LANGcs{
\Question{
Číslo komplexně sdružené k číslu 
\[ 
 z = \mathrm{i} + 3\mathrm{i}(2 -\mathrm{i})^{2} - 4(1 -\mathrm{i})^{3}
\]
je:}
{\(20 - 18\mathrm{i}\)}{\(20 - 24\mathrm{i}\)}{\(20 + 18\mathrm{i}\)}{\(- 8 + 26\mathrm{i}\)}{}{}}
\LANGsk{
\Question{
Číslo komplexne združené k číslu
\[ 
 \mathrm{i} + 3\mathrm{i}(2 -\mathrm{i})^{2} - 4(1 -\mathrm{i})^{3}
\] je:}
{\(20 - 18\mathrm{i}\)}{\(20 - 24\mathrm{i}\)}{\(20 + 18\mathrm{i}\)}{\(- 8 + 26\mathrm{i}\)}{}{}}
\LANGes{
\Question{
Calcula el conjugado del siguinte número complejo
\[ 
 \mathrm{i} + 3\mathrm{i}(2 -\mathrm{i})^{2} - 4(1 -\mathrm{i})^{3}
\]}
{\(20 - 18\mathrm{i}\)}{\(20 - 24\mathrm{i}\)}{\(20 + 18\mathrm{i}\)}{\(- 8 + 26\mathrm{i}\)}{}{}}
\End

\Data\ID{9000035803}
\Updated{2020-07-18 11:07:44}
\Level{A}
\SubArea{Complex numbers in algebraic and trigonometric form}
\LANGen{
\Question{
Given the complex number \(z = -1 + 2\mathrm{i}\),
find the imaginary part of the complex number
\(\frac{1}
{z}\).}
{\(-\frac{2} 
{5}\)}{\(\frac{1}
{2}\)}{\(\frac{2}
{5}\)}{\(-\frac{1} 
{2}\)}{}{}}
\LANGpl{
\Question{
Dane są liczby zespolone \(z = -1 + 2\mathrm{i}\),
wskaż część urojoną liczby zespolonej
\(\frac{1}
{z}\).}
{\(-\frac{2} 
{5}\)}{\(\frac{1}
{2}\)}{\(\frac{2}
{5}\)}{\(-\frac{1} 
{2}\)}{}{}}
\LANGcs{
\Question{
Imaginární část komplexního čísla
\(\frac{1}
{z}\), kde
\(z = -1 + 2\mathrm{i}\), je
rovna číslu:}
{\(-\frac{2} 
{5}\)}{\(\frac{1}
{2}\)}{\(\frac{2}
{5}\)}{\(-\frac{1} 
{2}\)}{}{}}
\LANGsk{
\Question{
Je dané komplexné číslo \(z = -1 + 2\mathrm{i}\). Určte imaginárnu časť komplexného čísla 
\(\frac{1}
{z}\).}
{\(-\frac{2} 
{5}\)}{\(\frac{1}
{2}\)}{\(\frac{2}
{5}\)}{\(-\frac{1} 
{2}\)}{}{}}
\LANGes{
\Question{
Dado el número complejo \(z = -1 + 2\mathrm{i}\),
Determina la parte imaginaria del número complejo
\(\frac{1}
{z}\).}
{\(-\frac{2} 
{5}\)}{\(\frac{1}
{2}\)}{\(\frac{2}
{5}\)}{\(-\frac{1} 
{2}\)}{}{}}
\End

\Data\ID{9000035804}
\Updated{2020-07-18 11:07:34}
\Level{A}
\SubArea{Complex numbers in algebraic and trigonometric form}
\LANGen{
\Question{
Find the algebraic form of the following complex number. By \(\overline{z }\) the complex conjugate of \(z \) is denoted.
\[ 
 \overline{\overline{(2 + \mathrm{i}) }\; \overline{(3 + 2\mathrm{i}) } }
\]}
{\(4 + 7\mathrm{i}\)}{\(8 + 7\mathrm{i}\)}{\(8 - 7\mathrm{i}\)}{\(4 - 7\mathrm{i}\)}{}{}}
\LANGpl{
\Question{
Wskaż postać algebraiczną podanej  liczby zespolonej.
\[ 
 \overline{\overline{(2 + \mathrm{i}) }\; \overline{(3 + 2\mathrm{i}) } }
\]}
{\(4 + 7\mathrm{i}\)}{\(8 + 7\mathrm{i}\)}{\(8 - 7\mathrm{i}\)}{\(4 - 7\mathrm{i}\)}{}{}}
\LANGcs{
\Question{
Určete algebraický tvar komplexního čísla \(\overline{\overline{(2 + \mathrm{i}) }\; \overline{(3 + 2\mathrm{i}) } }\).}
{\(4 + 7\mathrm{i}\)}{\(8 + 7\mathrm{i}\)}{\(8 - 7\mathrm{i}\)}{\(4 - 7\mathrm{i}\)}{}{}}
\LANGsk{
\Question{
Určte algebraickú formu daného komplexného čísla. 
\[ 
 \overline{\overline{(2 + \mathrm{i}) }\; \overline{(3 + 2\mathrm{i}) } }
\]}
{\(4 + 7\mathrm{i}\)}{\(8 + 7\mathrm{i}\)}{\(8 - 7\mathrm{i}\)}{\(4 - 7\mathrm{i}\)}{}{}}
\LANGes{
\Question{
Determina la forma algebraica del siguiente número complejo \(\overline{\overline{(2 + \mathrm{i}) }\; \overline{(3 + 2\mathrm{i}) } }\).}
{\(4 + 7\mathrm{i}\)}{\(8 + 7\mathrm{i}\)}{\(8 - 7\mathrm{i}\)}{\(4 - 7\mathrm{i}\)}{}{}}
\End

\Data\ID{9000035807}
\Updated{2021-12-19 10:12:41}
\Level{A}
\SubArea{Complex numbers in algebraic and trigonometric form}
\LANGen{
\Question{
Given the complex numbers \(a = 2 - 3\mathrm{i}\),
\(b = 1 + 2\mathrm{i}\), find the
quotient \(\frac{a}
{b}\).}
{\(-\frac{4} 
{5} -\frac{7} 
{5}\mathrm{i}\)}{\(2 -\frac{3} 
{2}\mathrm{i}\)}{\(\frac{8}
{5} -\frac{7} 
{5}\mathrm{i}\)}{\(\frac{4}
{3} + \frac{7} 
{3}\mathrm{i}\)}{}{}}
\LANGpl{
\Question{
Dane są liczby zespolone \(a = 2 - 3\mathrm{i}\),
\(b = 1 + 2\mathrm{i}\), oblicz iloraz \(\frac{a}
{b}\).}
{\(-\frac{4} 
{5} -\frac{7} 
{5}\mathrm{i}\)}{\(2 -\frac{3} 
{2}\mathrm{i}\)}{\(\frac{8}
{5} -\frac{7} 
{5}\mathrm{i}\)}{\(\frac{4}
{3} + \frac{7} 
{3}\mathrm{i}\)}{}{}}
\LANGcs{
\Question{
Jsou dána komplexní čísla \(a = 2 - 3\mathrm{i}\),
\(b = 1 + 2\mathrm{i}\). Podíl
\(\frac{a}
{b}\) je
roven:}
{\(-\frac{4} 
{5} -\frac{7} 
{5}\mathrm{i}\)}{\(2 -\frac{3} 
{2}\mathrm{i}\)}{\(\frac{8}
{5} -\frac{7} 
{5}\mathrm{i}\)}{\(\frac{4}
{3} + \frac{7} 
{3}\mathrm{i}\)}{}{}}
\LANGsk{
\Question{
Sú dané komplexné čísla \(a = 2 - 3\mathrm{i}\),
\(b = 1 + 2\mathrm{i}\). Podiel \(\frac{a}
{b}\) sa rovná:}
{\(-\frac{4} 
{5} -\frac{7} 
{5}\mathrm{i}\)}{\(2 -\frac{3} 
{2}\mathrm{i}\)}{\(\frac{8}
{5} -\frac{7} 
{5}\mathrm{i}\)}{\(\frac{4}
{3} + \frac{7} 
{3}\mathrm{i}\)}{}{}}
\LANGes{
\Question{
Dados los números complejos \(a = 2 - 3\mathrm{i}\),
\(b = 1 + 2\mathrm{i}\), determina el cociente  \(\frac{a}
{b}\).}
{\(-\frac{4} 
{5} -\frac{7} 
{5}\mathrm{i}\)}{\(2 -\frac{3} 
{2}\mathrm{i}\)}{\(\frac{8}
{5} -\frac{7} 
{5}\mathrm{i}\)}{\(\frac{4}
{3} + \frac{7} 
{3}\mathrm{i}\)}{}{}}
\End

\Data\ID{9000037401}
\Updated{2020-07-18 11:07:10}
\Level{A}
\SubArea{Complex numbers in algebraic and trigonometric form}
\LANGen{
\Question{
Evaluate \(\mathrm{i}^{50}\).}
{\(- 1\)}{\(\mathrm{i}\)}{\(1\)}{\(-\mathrm{i}\)}{}{}}
\LANGpl{
\Question{
Oblicz \(\mathrm{i}^{50}\).}
{\(- 1\)}{\(\mathrm{i}\)}{\(1\)}{\(-\mathrm{i}\)}{}{}}
\LANGcs{
\Question{
Určete hodnotu výrazu \(\mathrm{i}^{50}\).}
{\(- 1\)}{\(\mathrm{i}\)}{\(1\)}{\(-\mathrm{i}\)}{}{}}
\LANGsk{
\Question{
Určte hodnotu výrazu \(\mathrm{i}^{50}\).}
{\(- 1\)}{\(\mathrm{i}\)}{\(1\)}{\(-\mathrm{i}\)}{}{}}
\LANGes{
\Question{
Calcula \(\mathrm{i}^{50}\).}
{\(- 1\)}{\(\mathrm{i}\)}{\(1\)}{\(-\mathrm{i}\)}{}{}}
\End

\Data\ID{9000037402}
\Updated{2020-07-18 11:07:19}
\Level{A}
\SubArea{Complex numbers in algebraic and trigonometric form}
\LANGen{
\Question{
Evaluate \(\mathrm{i}^{7}\).}
{\(-\mathrm{i}\)}{\(- 1\)}{\(\mathrm{i}\)}{\(1\)}{}{}}
\LANGpl{
\Question{
Oblicz \(\mathrm{i}^{7}\).}
{\(-\mathrm{i}\)}{\(- 1\)}{\(\mathrm{i}\)}{\(1\)}{}{}}
\LANGcs{
\Question{
Určete hodnotu výrazu \(\mathrm{i}^{7}\).}
{\(-\mathrm{i}\)}{\(- 1\)}{\(\mathrm{i}\)}{\(1\)}{}{}}
\LANGsk{
\Question{
Určte hodnotu výrazu \(\mathrm{i}^{7}\).}
{\(-\mathrm{i}\)}{\(- 1\)}{\(\mathrm{i}\)}{\(1\)}{}{}}
\LANGes{
\Question{
Calcula \(\mathrm{i}^{7}\).}
{\(-\mathrm{i}\)}{\(- 1\)}{\(\mathrm{i}\)}{\(1\)}{}{}}
\End

\Data\ID{9000037501}
\Updated{2020-07-18 11:07:31}
\Level{A}
\SubArea{Complex numbers in algebraic and trigonometric form}
\LANGen{
\Question{
Find the absolute value of the following complex number. 
\[ 
 3 + \sqrt{2}\mathrm{i}
\]}
{\(\sqrt{11}\)}{\(\sqrt{13}\)}{\(3\)}{\(3\sqrt{2}\)}{}{}}
\LANGpl{
\Question{
Określ wartość bezwzględną podanej liczby zespolonej.
\[ 
 3 + \sqrt{2}\mathrm{i}
\]}
{\(\sqrt{11}\)}{\(\sqrt{13}\)}{\(3\)}{\(3\sqrt{2}\)}{}{}}
\LANGcs{
\Question{
Určete absolutní hodnotu komplexního čísla
\(a.\) 
\[ 
 a = 3 + \sqrt{2}\mathrm{i}
\]}
{\(\sqrt{11}\)}{\(\sqrt{13}\)}{\(3\)}{\(3\sqrt{2}\)}{}{}}
\LANGsk{
\Question{
Určte absolútnu hodnotu daného komplexného čísla. 
\[ 
 3 + \sqrt{2}\mathrm{i}
\]}
{\(\sqrt{11}\)}{\(\sqrt{13}\)}{\(3\)}{\(3\sqrt{2}\)}{}{}}
\LANGes{
\Question{
Determina el valor absoluto del siguiente número complejo
\[ 
 3 + \sqrt{2}\mathrm{i}
\]}
{\(\sqrt{11}\)}{\(\sqrt{13}\)}{\(3\)}{\(3\sqrt{2}\)}{}{}}
\End

\Data\ID{9000037502}
\Updated{2020-07-18 11:07:37}
\Level{A}
\SubArea{Complex numbers in algebraic and trigonometric form}
\LANGen{
\Question{
Find the total sum of the complex numbers
\(a\),
\(b\) and
\(c\).
\[ 
 a = 3 + \sqrt{2}\mathrm{i},\quad b = 1 - 4\mathrm{i},\quad c = \sqrt{3} - 3\mathrm{i}
\]}
{\(4 + \sqrt{3} + \mathrm{i}(\sqrt{2} - 7)\)}{\(4 + \mathrm{i}\sqrt{3}\)}{\(4 + \sqrt{2} + \mathrm{i}(\sqrt{3} - 3)\)}{\(4 + \sqrt{3} -\mathrm{i}(\sqrt{2} - 7)\)}{}{}}
\LANGpl{
\Question{
Wskaż całkowitą sumę podanych liczb zespolonych
\(a\),
\(b\) i
\(c\).
\[ 
 a = 3 + \sqrt{2}\mathrm{i},\quad b = 1 - 4\mathrm{i},\quad c = \sqrt{3} - 3\mathrm{i}
\]}
{\(4 + \sqrt{3} + \mathrm{i}(\sqrt{2} - 7)\)}{\(4 + \mathrm{i}\sqrt{3}\)}{\(4 + \sqrt{2} + \mathrm{i}(\sqrt{3} - 3)\)}{\(4 + \sqrt{3} -\mathrm{i}(\sqrt{2} - 7)\)}{}{}}
\LANGcs{
\Question{
Určete součet komplexních čísel
\(a\),
\(b\),
\(c\).
\[ 
 a = 3 + \sqrt{2}\mathrm{i},\quad b = 1 - 4\mathrm{i},\quad c = \sqrt{3} - 3\mathrm{i}
\]}
{\(4 + \sqrt{3} + \mathrm{i}(\sqrt{2} - 7)\)}{\(4 + \mathrm{i}\sqrt{3}\)}{\(4 + \sqrt{2} + \mathrm{i}(\sqrt{3} - 3)\)}{\(4 + \sqrt{3} -\mathrm{i}(\sqrt{2} - 7)\)}{}{}}
\LANGsk{
\Question{
Určte súčet komplexných čísel 
\(a\),
\(b\) a
\(c\).
\[ 
 a = 3 + \sqrt{2}\mathrm{i},\quad b = 1 - 4\mathrm{i},\quad c = \sqrt{3} - 3\mathrm{i}
\]}
{\(4 + \sqrt{3} + \mathrm{i}(\sqrt{2} - 7)\)}{\(4 + \mathrm{i}\sqrt{3}\)}{\(4 + \sqrt{2} + \mathrm{i}(\sqrt{3} - 3)\)}{\(4 + \sqrt{3} -\mathrm{i}(\sqrt{2} - 7)\)}{}{}}
\LANGes{
\Question{
Calcula la suma total de números complejos.
\(a\),
\(b\) y
\(c\).
\[ 
 a = 3 + \sqrt{2}\mathrm{i},\quad b = 1 - 4\mathrm{i},\quad c = \sqrt{3} - 3\mathrm{i}
\]}
{\(4 + \sqrt{3} + \mathrm{i}(\sqrt{2} - 7)\)}{\(4 + \mathrm{i}\sqrt{3}\)}{\(4 + \sqrt{2} + \mathrm{i}(\sqrt{3} - 3)\)}{\(4 + \sqrt{3} -\mathrm{i}(\sqrt{2} - 7)\)}{}{}}
\End

\Data\ID{9000037503}
\Updated{2022-04-23 05:04:48}
\Level{A}
\SubArea{Complex numbers in algebraic and trigonometric form}
\LANGen{
\Question{
Given complex numbers 
\[ 
 a = \sqrt{2} + \sqrt{3}\mathrm{i},\quad b = \sqrt{2} -\sqrt{3}\mathrm{i},
\]
find the product \(ab\).}
{\(5\)}{\(2\)}{\(\sqrt{2} + \mathrm{i}\sqrt{3}\)}{\(\sqrt{2} -\mathrm{i}\sqrt{3}\)}{}{}}
\LANGpl{
\Question{
Dane są liczby zespolone 
\[ 
 a = \sqrt{2} + \sqrt{3}\mathrm{i},\quad b = \sqrt{2} -\sqrt{3}\mathrm{i},
\]
oblicz iloczyn \(ab\).}
{\(5\)}{\(2\)}{\(\sqrt{2} + \mathrm{i}\sqrt{3}\)}{\(\sqrt{2} -\mathrm{i}\sqrt{3}\)}{}{}}
\LANGcs{
\Question{
Určete součin komplexních čísel
\(a\),
\(b\).
\[ 
 a = \sqrt{2} + \sqrt{3}\mathrm{i},\quad b = \sqrt{2} -\sqrt{3}\mathrm{i}
\]}
{\(5\)}{\(2\)}{\(\sqrt{2} + \mathrm{i}\sqrt{3}\)}{\(\sqrt{2} -\mathrm{i}\sqrt{3}\)}{}{}}
\LANGsk{
\Question{
Sú dané komplexné čísla  
\[ 
 a = \sqrt{2} + \sqrt{3}\mathrm{i},\quad b = \sqrt{2} -\sqrt{3}\mathrm{i}\text{.}
\]
Určte súčin \(ab\).}
{\(5\)}{\(2\)}{\(\sqrt{2} + \mathrm{i}\sqrt{3}\)}{\(\sqrt{2} -\mathrm{i}\sqrt{3}\)}{}{}}
\LANGes{
\Question{
Dados los números complejos
\[ 
 a = \sqrt{2} + \sqrt{3}\mathrm{i},\quad b = \sqrt{2} -\sqrt{3}\mathrm{i},
\]
calcula el producto \(ab\).}
{\(5\)}{\(2\)}{\(\sqrt{2} + \mathrm{i}\sqrt{3}\)}{\(\sqrt{2} -\mathrm{i}\sqrt{3}\)}{}{}}
\End

\Data\ID{9000037504}
\Updated{2020-07-18 11:07:42}
\Level{A}
\SubArea{Complex numbers in algebraic and trigonometric form}
\LANGen{
\Question{
Given complex numbers 
\[ 
 a = 5 + 2\mathrm{i},\quad b = 3 -\mathrm{i},\quad c = \mathrm{i}\text{,}
\]
find the product \(abc\).}
{\(- 1 + 17\mathrm{i}\)}{\(1 - 17\mathrm{i}\)}{\(- 1 - 17\mathrm{i}\)}{\(1 + 17\mathrm{i}\)}{}{}}
\LANGpl{
\Question{
Dane są liczby zespolone
\[ 
 a = 5 + 2\mathrm{i},\quad b = 3 -\mathrm{i},\quad c = \mathrm{i}\text{,}
\]
oblicz iloczyn \(abc\).}
{\(- 1 + 17\mathrm{i}\)}{\(1 - 17\mathrm{i}\)}{\(- 1 - 17\mathrm{i}\)}{\(1 + 17\mathrm{i}\)}{}{}}
\LANGcs{
\Question{
Určete součin komplexních čísel
\(a\),
\(b\),
\(c\).
\[ 
 a = 5 + 2\mathrm{i},\quad b = 3 -\mathrm{i},\quad c = \mathrm{i}
\]}
{\(- 1 + 17\mathrm{i}\)}{\(1 - 17\mathrm{i}\)}{\(- 1 - 17\mathrm{i}\)}{\(1 + 17\mathrm{i}\)}{}{}}
\LANGsk{
\Question{
Sú dané komplexné čísla
\[ 
 a = 5 + 2\mathrm{i},\quad b = 3 -\mathrm{i},\quad c = \mathrm{i}\text{.}
\]
Určte súčin \(abc\).}
{\(- 1 + 17\mathrm{i}\)}{\(1 - 17\mathrm{i}\)}{\(- 1 - 17\mathrm{i}\)}{\(1 + 17\mathrm{i}\)}{}{}}
\LANGes{
\Question{
Dados los números complejos
\[ 
 a = 5 + 2\mathrm{i},\quad b = 3 -\mathrm{i},\quad c = \mathrm{i}\text{,}
\]
calcula el producto \(abc\).}
{\(- 1 + 17\mathrm{i}\)}{\(1 - 17\mathrm{i}\)}{\(- 1 - 17\mathrm{i}\)}{\(1 + 17\mathrm{i}\)}{}{}}
\End

\Data\ID{9000037505}
\Updated{2020-07-18 11:07:06}
\Level{A}
\SubArea{Complex numbers in algebraic and trigonometric form}
\LANGen{
\Question{
Find the complex conjugate of the following complex number. 
\[ 
 -2\sqrt{3} -\mathrm{i}
\]}
{\(- 2\sqrt{3} + \mathrm{i}\)}{\(2\sqrt{3} -\mathrm{i}\)}{\(11\)}{\(10\mathrm{i}\)}{}{}}
\LANGpl{
\Question{
Określ sprężoną liczbę zespoloną do podanej liczby zespolonej.
\[ 
 -2\sqrt{3} -\mathrm{i}
\]}
{\(- 2\sqrt{3} + \mathrm{i}\)}{\(2\sqrt{3} -\mathrm{i}\)}{\(11\)}{\(10\mathrm{i}\)}{}{}}
\LANGcs{
\Question{
Určete číslo komplexně sdružené k číslu
\(a\).
\[ 
 a = -2\sqrt{3} -\mathrm{i}
\]}
{\(- 2\sqrt{3} + \mathrm{i}\)}{\(2\sqrt{3} -\mathrm{i}\)}{\(11\)}{\(10\mathrm{i}\)}{}{}}
\LANGsk{
\Question{
Určte číslo komplexne združené k danému komplexnému číslu. 
\[ 
 -2\sqrt{3} -\mathrm{i}
\]}
{\(- 2\sqrt{3} + \mathrm{i}\)}{\(2\sqrt{3} -\mathrm{i}\)}{\(11\)}{\(10\mathrm{i}\)}{}{}}
\LANGes{
\Question{
Determina el conjugado del número complejo. 
\[ 
 -2\sqrt{3} -\mathrm{i}
\]}
{\(- 2\sqrt{3} + \mathrm{i}\)}{\(2\sqrt{3} -\mathrm{i}\)}{\(11\)}{\(10\mathrm{i}\)}{}{}}
\End

\Data\ID{9000037506}
\Updated{2022-04-23 05:04:05}
\Level{A}
\SubArea{Complex numbers in algebraic and trigonometric form}
\LANGen{
\Question{
Given complex numbers 
\[ 
 a = 3 + 5\mathrm{i}\text{, }\quad b = 2 -\mathrm{i}\text{, }
\]
find the quotient \(\frac{a}
{b}\).}
{\(\frac{1}
{5} + \mathrm{i}\frac{13} 
 {5} \)}{\(\frac{1}
{3} + \mathrm{i}\frac{13} 
 {3} \)}{\(\frac{1}
{5} + \mathrm{i}\frac{7} 
{5}\)}{\(\frac{1}
{3} + \mathrm{i}\frac{7} 
{3}\)}{}{}}
\LANGpl{
\Question{
Dane są liczby zespolone
\[ 
 a = 3 + 5\mathrm{i}\text{, }\quad b = 2 -\mathrm{i}\text{, }
\]
oblicz iloraz \(\frac{a}
{b}\).}
{\(\frac{1}
{5} + \mathrm{i}\frac{13} 
 {5} \)}{\(\frac{1}
{3} + \mathrm{i}\frac{13} 
 {3} \)}{\(\frac{1}
{5} + \mathrm{i}\frac{7} 
{5}\)}{\(\frac{1}
{3} + \mathrm{i}\frac{7} 
{3}\)}{}{}}
\LANGcs{
\Question{
Určete podíl \(\frac{a}
{b}\) 
komplexních čísel \(a\),
\(b\).
\[ 
 a = 3 + 5\mathrm{i},\quad b = 2 -\mathrm{i}
\]}
{\(\frac{1}
{5} + \mathrm{i}\frac{13} 
 {5} \)}{\(\frac{1}
{3} + \mathrm{i}\frac{13} 
 {3} \)}{\(\frac{1}
{5} + \mathrm{i}\frac{7} 
{5}\)}{\(\frac{1}
{3} + \mathrm{i}\frac{7} 
{3}\)}{}{}}
\LANGsk{
\Question{
Sú dané komplexné čísla 
\[ 
 a = 3 + 5\mathrm{i}\text{, }\quad b = 2 -\mathrm{i}\text{.}
\]
Určte podiel \(\frac{a}
{b}\).}
{\(\frac{1}
{5} + \mathrm{i}\frac{13} 
 {5} \)}{\(\frac{1}
{3} + \mathrm{i}\frac{13} 
 {3} \)}{\(\frac{1}
{5} + \mathrm{i}\frac{7} 
{5}\)}{\(\frac{1}
{3} + \mathrm{i}\frac{7} 
{3}\)}{}{}}
\LANGes{
\Question{
Dados los números complejos
\[ 
 a = 3 + 5\mathrm{i}\text{, }\quad b = 2 -\mathrm{i}\text{, }
\]
determina el cociente \(\frac{a}
{b}\).}
{\(\frac{1}
{5} + \mathrm{i}\frac{13} 
 {5} \)}{\(\frac{1}
{3} + \mathrm{i}\frac{13} 
 {3} \)}{\(\frac{1}
{5} + \mathrm{i}\frac{7} 
{5}\)}{\(\frac{1}
{3} + \mathrm{i}\frac{7} 
{3}\)}{}{}}
\End

\Data\ID{9000037507}
\Updated{2022-04-23 05:04:48}
\Level{A}
\SubArea{Complex numbers in algebraic and trigonometric form}
\LANGen{
\Question{
Given complex numbers 
\[ 
 a = \sqrt{3} + 2\mathrm{i}\text{, }\quad b = \sqrt{2} -\mathrm{i}\text{, }
\]
find the quotient \(\frac{a}
{b}\).}
{\(\frac{\sqrt{6}-2}
 {3} + \mathrm{i}\frac{2\sqrt{2}+\sqrt{3}} 
 {3} \)}{\(\frac{\sqrt{6}-2}
 {3} -\mathrm{i}\frac{2\sqrt{2}+\sqrt{3}} 
 {3} \)}{\(\frac{\sqrt{6}-3}
 {2} + \mathrm{i}\frac{2\sqrt{2}+\sqrt{3}} 
 {2} \)}{\(\frac{\sqrt{6}-2}
 {2} -\mathrm{i}\frac{2\sqrt{2}+\sqrt{3}} 
 {2} \)}{}{}}
\LANGpl{
\Question{
Dane są liczby zespolone
\[ 
 a = \sqrt{3} + 2\mathrm{i}\text{, }\quad b = \sqrt{2} -\mathrm{i}\text{, }
\]
oblicz iloraz \(\frac{a}
{b}\).}
{\(\frac{\sqrt{6}-2}
 {3} + \mathrm{i}\frac{2\sqrt{2}+\sqrt{3}} 
 {3} \)}{\(\frac{\sqrt{6}-2}
 {3} -\mathrm{i}\frac{2\sqrt{2}+\sqrt{3}} 
 {3} \)}{\(\frac{\sqrt{6}-3}
 {2} + \mathrm{i}\frac{2\sqrt{2}+\sqrt{3}} 
 {2} \)}{\(\frac{\sqrt{6}-2}
 {2} -\mathrm{i}\frac{2\sqrt{2}+\sqrt{3}} 
 {2} \)}{}{}}
\LANGcs{
\Question{
Určete podíl \(\frac{a}
{b}\) 
komplexních čísel \(a\),
\(b\).
\[ 
 a = \sqrt{3} + 2\mathrm{i},\quad b = \sqrt{2} -\mathrm{i}
\]}
{\(\frac{\sqrt{6}-2}
 {3} + \mathrm{i}\frac{2\sqrt{2}+\sqrt{3}} 
 {3} \)}{\(\frac{\sqrt{6}-2}
 {3} -\mathrm{i}\frac{2\sqrt{2}+\sqrt{3}} 
 {3} \)}{\(\frac{\sqrt{6}-3}
 {2} + \mathrm{i}\frac{2\sqrt{2}+\sqrt{3}} 
 {2} \)}{\(\frac{\sqrt{6}-2}
 {2} -\mathrm{i}\frac{2\sqrt{2}+\sqrt{3}} 
 {2} \)}{}{}}
\LANGsk{
\Question{
Sú dané komplexné čísla  
\[ 
 a = \sqrt{3} + 2\mathrm{i}\text{, }\quad b = \sqrt{2} -\mathrm{i}\text{.}
\]
Určte podiel \(\frac{a}
{b}\).}
{\(\frac{\sqrt{6}-2}
 {3} + \mathrm{i}\frac{2\sqrt{2}+\sqrt{3}} 
 {3} \)}{\(\frac{\sqrt{6}-2}
 {3} -\mathrm{i}\frac{2\sqrt{2}+\sqrt{3}} 
 {3} \)}{\(\frac{\sqrt{6}-3}
 {2} + \mathrm{i}\frac{2\sqrt{2}+\sqrt{3}} 
 {2} \)}{\(\frac{\sqrt{6}-2}
 {2} -\mathrm{i}\frac{2\sqrt{2}+\sqrt{3}} 
 {2} \)}{}{}}
\LANGes{
\Question{
Dados los números complejos
\[ 
 a = \sqrt{3} + 2\mathrm{i}\text{, }\quad b = \sqrt{2} -\mathrm{i}\text{, }
\]
determina el cociente\(\frac{a}
{b}\).}
{\(\frac{\sqrt{6}-2}
 {3} + \mathrm{i}\frac{2\sqrt{2}+\sqrt{3}} 
 {3} \)}{\(\frac{\sqrt{6}-2}
 {3} -\mathrm{i}\frac{2\sqrt{2}+\sqrt{3}} 
 {3} \)}{\(\frac{\sqrt{6}-3}
 {2} + \mathrm{i}\frac{2\sqrt{2}+\sqrt{3}} 
 {2} \)}{\(\frac{\sqrt{6}-2}
 {2} -\mathrm{i}\frac{2\sqrt{2}+\sqrt{3}} 
 {2} \)}{}{}}
\End

\Data\ID{1003082301}
\Updated{2022-04-23 05:04:17}
\Level{B}
\SubArea{Complex numbers in algebraic and trigonometric form}
\LANGen{
\Question{
Given the complex numbers  \( a=\sqrt2\left(\cos160^{\circ}+\mathrm{i}\cdot\sin160^{\circ}\right) \), \( b=3\sqrt2\left(\cos150^{\circ}+\mathrm{i}\cdot\sin150^{\circ}\right) \) and \( c=2\left(\cos240^{\circ}+\mathrm{i}\cdot\sin240^{\circ}\right) \), evaluate \( a\cdot b\cdot c \).}
{\(12\left(\cos190^{\circ}+\mathrm{i}\cdot\sin190^{\circ}\right) \)}{\(12\left(\cos10^{\circ}+\mathrm{i}\cdot\sin10^{\circ}\right) \)}{\(12\left(\cos10^{\circ}-\mathrm{i}\cdot\sin10^{\circ}\right) \)}{\(12\left(\cos190^{\circ}-\mathrm{i}\cdot\sin190^{\circ}\right) \)}{}{}}
\LANGpl{
\Question{
Dane są liczby zespolone \( a=\sqrt2\left(\cos160^{\circ}+\mathrm{i}\cdot\sin160^{\circ}\right) \), \( b=3\sqrt2\left(\cos150^{\circ}+\mathrm{i}\cdot\sin150^{\circ}\right) \) i \( c=2\left(\cos240^{\circ}+\mathrm{i}\cdot\sin240^{\circ}\right) \), oblicz \( a\cdot b\cdot c \).}
{\(12\left(\cos190^{\circ}+\mathrm{i}\cdot\sin190^{\circ}\right) \)}{\(12\left(\cos10^{\circ}+\mathrm{i}\cdot\sin10^{\circ}\right) \)}{\(12\left(\cos10^{\circ}-\mathrm{i}\cdot\sin10^{\circ}\right) \)}{\(12\left(\cos190^{\circ}-\mathrm{i}\cdot\sin190^{\circ}\right) \)}{}{}}
\LANGcs{
\Question{
Jsou daná komplexní čísla   \( a=\sqrt2\left(\cos160^{\circ}+\mathrm{i}\cdot\sin160^{\circ}\right) \), \( b=3\sqrt2\left(\cos150^{\circ}+\mathrm{i}\cdot\sin150^{\circ}\right) \) a \( c=2\left(\cos240^{\circ}+\mathrm{i}\cdot\sin240^{\circ}\right) \). Vypočítejte \( a\cdot b\cdot c \).}
{\(12\left(\cos190^{\circ}+\mathrm{i}\cdot\sin190^{\circ}\right) \)}{\(12\left(\cos10^{\circ}+\mathrm{i}\cdot\sin10^{\circ}\right) \)}{\(12\left(\cos10^{\circ}-\mathrm{i}\cdot\sin10^{\circ}\right) \)}{\(12\left(\cos190^{\circ}-\mathrm{i}\cdot\sin190^{\circ}\right) \)}{}{}}
\LANGsk{
\Question{
Sú dané komplexné čísla   \( a=\sqrt2\left(\cos160^{\circ}+\mathrm{i}\cdot\sin160^{\circ}\right) \), \( b=3\sqrt2\left(\cos150^{\circ}+\mathrm{i}\cdot\sin150^{\circ}\right) \) a \( c=2\left(\cos240^{\circ}+\mathrm{i}\cdot\sin240^{\circ}\right) \). Vypočítajte \( a\cdot b\cdot c \).}
{\(12\left(\cos190^{\circ}+\mathrm{i}\cdot\sin190^{\circ}\right) \)}{\(12\left(\cos10^{\circ}+\mathrm{i}\cdot\sin10^{\circ}\right) \)}{\(12\left(\cos10^{\circ}-\mathrm{i}\cdot\sin10^{\circ}\right) \)}{\(12\left(\cos190^{\circ}-\mathrm{i}\cdot\sin190^{\circ}\right) \)}{}{}}
\LANGes{
\Question{
Dados los números complejos  \( a=\sqrt2\left(\cos160^{\circ}+\mathrm{i}\cdot\sin160^{\circ}\right) \), \( b=3\sqrt2\left(\cos150^{\circ}+\mathrm{i}\cdot\sin150^{\circ}\right) \) y \( c=2\left(\cos240^{\circ}+\mathrm{i}\cdot\sin240^{\circ}\right) \), calcula \( a\cdot b\cdot c \).}
{\(12\left(\cos190^{\circ}+\mathrm{i}\cdot\sin190^{\circ}\right) \)}{\(12\left(\cos10^{\circ}+\mathrm{i}\cdot\sin10^{\circ}\right) \)}{\(12\left(\cos10^{\circ}-\mathrm{i}\cdot\sin10^{\circ}\right) \)}{\(12\left(\cos190^{\circ}-\mathrm{i}\cdot\sin190^{\circ}\right) \)}{}{}}
\End

\Data\ID{1003082302}
\Updated{2020-07-18 08:07:19}
\Level{B}
\SubArea{Complex numbers in algebraic and trigonometric form}
\LANGen{
\Question{
Given the complex numbers \( a=10\left(\cos\frac43\pi+\mathrm{i}\cdot\sin\frac43\pi\right) \), \( b=7\left(\cos150^{\circ}+\mathrm{i}\cdot\sin150^{\circ}\right) \) and \( c=5\left(\cos\frac74\pi+\mathrm{i}\cdot\sin\frac74\pi \right) \), evaluate \( \frac{a\cdot b}c \).}
{\( 14\left(\cos\frac5{12}\pi+\mathrm{i}\cdot\sin\frac5{12}\pi\right) \)}{\( 14\left(\cos\frac14\pi+\mathrm{i}\cdot\sin\frac14\pi\right) \)}{\( 14\left(\cos\frac{23}{12}\pi+\mathrm{i}\cdot\sin\frac{23}{12}\pi\right) \)}{\( 14\left(\cos\frac54\pi+\mathrm{i}\cdot\sin\frac54\pi\right) \)}{}{}}
\LANGpl{
\Question{
Dane są liczby zespolone \( a=10\left(\cos\frac43\pi+\mathrm{i}\cdot\sin\frac43\pi\right) \), \( b=7\left(\cos150^{\circ}+\mathrm{i}\cdot\sin150^{\circ}\right) \) i \( c=5\left(\cos\frac74\pi+\mathrm{i}\cdot\sin\frac74\pi \right) \), oblicz \( \frac{a\cdot b}c \).}
{\( 14\left(\cos\frac5{12}\pi+\mathrm{i}\cdot\sin\frac5{12}\pi\right) \)}{\( 14\left(\cos\frac14\pi+\mathrm{i}\cdot\sin\frac14\pi\right) \)}{\( 14\left(\cos\frac{23}{12}\pi+\mathrm{i}\cdot\sin\frac{23}{12}\pi\right) \)}{\( 14\left(\cos\frac54\pi+\mathrm{i}\cdot\sin\frac54\pi\right) \)}{}{}}
\LANGcs{
\Question{
Jsou daná komplexní čísla \( a=10\left(\cos\frac43\pi+\mathrm{i}\cdot\sin\frac43\pi\right) \), \( b=7\left(\cos150^{\circ}+\mathrm{i}\cdot\sin150^{\circ}\right) \) a \( c=5\left(\cos\frac74\pi+\mathrm{i}\cdot\sin\frac74\pi \right) \). Vypočítejte \( \frac{a\cdot b}c \).}
{\( 14\left(\cos\frac5{12}\pi+\mathrm{i}\cdot\sin\frac5{12}\pi\right) \)}{\( 14\left(\cos\frac14\pi+\mathrm{i}\cdot\sin\frac14\pi\right) \)}{\( 14\left(\cos\frac{23}{12}\pi+\mathrm{i}\cdot\sin\frac{23}{12}\pi\right) \)}{\( 14\left(\cos\frac54\pi+\mathrm{i}\cdot\sin\frac54\pi\right) \)}{}{}}
\LANGsk{
\Question{
Sú dané komplexné čísla \( a=10\left(\cos\frac43\pi+\mathrm{i}\cdot\sin\frac43\pi\right) \), \( b=7\left(\cos150^{\circ}+\mathrm{i}\cdot\sin150^{\circ}\right) \) a \( c=5\left(\cos\frac74\pi+\mathrm{i}\cdot\sin\frac74\pi \right) \). Vypočítajte \( \frac{a\cdot b}c \).}
{\( 14\left(\cos\frac5{12}\pi+\mathrm{i}\cdot\sin\frac5{12}\pi\right) \)}{\( 14\left(\cos\frac14\pi+\mathrm{i}\cdot\sin\frac14\pi\right) \)}{\( 14\left(\cos\frac{23}{12}\pi+\mathrm{i}\cdot\sin\frac{23}{12}\pi\right) \)}{\( 14\left(\cos\frac54\pi+\mathrm{i}\cdot\sin\frac54\pi\right) \)}{}{}}
\LANGes{
\Question{
Dados los números complejos \( a=10\left(\cos\frac43\pi+\mathrm{i}\cdot\sin\frac43\pi\right) \), \( b=7\left(\cos150^{\circ}+\mathrm{i}\cdot\sin150^{\circ}\right) \) y \( c=5\left(\cos\frac74\pi+\mathrm{i}\cdot\sin\frac74\pi \right) \) calcula \( \frac{a\cdot b}c \).}
{\( 14\left(\cos\frac5{12}\pi+\mathrm{i}\cdot\sin\frac5{12}\pi\right) \)}{\( 14\left(\cos\frac14\pi+\mathrm{i}\cdot\sin\frac14\pi\right) \)}{\( 14\left(\cos\frac{23}{12}\pi+\mathrm{i}\cdot\sin\frac{23}{12}\pi\right) \)}{\( 14\left(\cos\frac54\pi+\mathrm{i}\cdot\sin\frac54\pi\right) \)}{}{}}
\End

\Data\ID{1003082303}
\Updated{2022-04-23 05:04:17}
\Level{B}
\SubArea{Complex numbers in algebraic and trigonometric form}
\LANGen{
\Question{
Given the complex numbers \( a=6\sqrt2\left(\cos\frac{\pi}3+\mathrm{i}\cdot\sin\frac{\pi}3\right) \), \( b=3\sqrt2\left(\cos\frac56\pi+\mathrm{i}\cdot\sin\frac56\pi\right) \) and \( c=2\left(\cos240^{\circ}+\mathrm{i}\cdot\sin240^{\circ}\right) \), evaluate \( \frac a{b\cdot c} \).}
{\( \cos\frac{\pi}6+\mathrm{i}\cdot\sin\frac{\pi}6 \)}{\(  \cos\frac{11}6\pi+\mathrm{i}\cdot\sin\frac{11}6\pi \)}{\( 4\left(\cos\frac{\pi}6\pi+\mathrm{i}\cdot\sin\frac{\pi}6\pi\right) \)}{\( 4\left(\cos\frac{11}6\pi+\mathrm{i}\cdot\sin\frac{11}6\pi\right) \)}{}{}}
\LANGpl{
\Question{
Dane są liczby zespolone \( a=6\sqrt2\left(\cos\frac{\pi}3+\mathrm{i}\cdot\sin\frac{\pi}3\right) \), \( b=3\sqrt2\left(\cos\frac56\pi+\mathrm{i}\cdot\sin\frac56\pi\right) \) i \( c=2\left(\cos240^{\circ}+\mathrm{i}\cdot\sin240^{\circ}\right) \), oblicz \( \frac a{b\cdot c} \).}
{\( \cos\frac{\pi}6+\mathrm{i}\cdot\sin\frac{\pi}6 \)}{\(  \cos\frac{11}6\pi+\mathrm{i}\cdot\sin\frac{11}6\pi \)}{\( 4\left(\cos\frac{\pi}6\pi+\mathrm{i}\cdot\sin\frac{\pi}6\pi\right) \)}{\( 4\left(\cos\frac{11}6\pi+\mathrm{i}\cdot\sin\frac{11}6\pi\right) \)}{}{}}
\LANGcs{
\Question{
Jsou daná komplexní čísla \( a=6\sqrt2\left(\cos\frac{\pi}3+\mathrm{i}\cdot\sin\frac{\pi}3\right) \), \( b=3\sqrt2\left(\cos\frac56\pi+\mathrm{i}\cdot\sin\frac56\pi\right) \) a \( c=2\left(\cos240^{\circ}+\mathrm{i}\cdot\sin240^{\circ}\right) \). Vypočítejte \( \frac a{b\cdot c} \).}
{\( \cos\frac{\pi}6+\mathrm{i}\cdot\sin\frac{\pi}6 \)}{\(  \cos\frac{11}6\pi+\mathrm{i}\cdot\sin\frac{11}6\pi \)}{\( 4\left(\cos\frac{\pi}6\pi+\mathrm{i}\cdot\sin\frac{\pi}6\pi\right) \)}{\( 4\left(\cos\frac{11}6\pi+\mathrm{i}\cdot\sin\frac{11}6\pi\right) \)}{}{}}
\LANGsk{
\Question{
Sú dané komplexné čísla \( a=6\sqrt2\left(\cos\frac{\pi}3+\mathrm{i}\cdot\sin\frac{\pi}3\right) \), \( b=3\sqrt2\left(\cos\frac56\pi+\mathrm{i}\cdot\sin\frac56\pi\right) \) a \( c=2\left(\cos240^{\circ}+\mathrm{i}\cdot\sin240^{\circ}\right) \). Vypočítajte \( \frac a{b\cdot c} \).}
{\( \cos\frac{\pi}6+\mathrm{i}\cdot\sin\frac{\pi}6 \)}{\(  \cos\frac{11}6\pi+\mathrm{i}\cdot\sin\frac{11}6\pi \)}{\( 4\left(\cos\frac{\pi}6\pi+\mathrm{i}\cdot\sin\frac{\pi}6\pi\right) \)}{\( 4\left(\cos\frac{11}6\pi+\mathrm{i}\cdot\sin\frac{11}6\pi\right) \)}{}{}}
\LANGes{
\Question{
Dados los números complejos  \( a=6\sqrt2\left(\cos\frac{\pi}3+\mathrm{i}\cdot\sin\frac{\pi}3\right) \), \( b=3\sqrt2\left(\cos\frac56\pi+\mathrm{i}\cdot\sin\frac56\pi\right) \) y \( c=2\left(\cos240^{\circ}+\mathrm{i}\cdot\sin240^{\circ}\right) \), calcula  \( \frac a{b\cdot c} \).}
{\( \cos\frac{\pi}6+\mathrm{i}\cdot\sin\frac{\pi}6 \)}{\(  \cos\frac{11}6\pi+\mathrm{i}\cdot\sin\frac{11}6\pi \)}{\( 4\left(\cos\frac{\pi}6\pi+\mathrm{i}\cdot\sin\frac{\pi}6\pi\right) \)}{\( 4\left(\cos\frac{11}6\pi+\mathrm{i}\cdot\sin\frac{11}6\pi\right) \)}{}{}}
\End

\Data\ID{1003082305}
\Updated{2022-04-23 05:04:17}
\Level{B}
\SubArea{Complex numbers in algebraic and trigonometric form}
\LANGen{
\Question{
Let \( [x;y]\in\mathbb{R}\times\mathbb{R} \), \( z_1   = 5 + xy\,\mathrm{i} \)  and \( z_2  = x + y - 4\,\mathrm{i} \). Find all \( [x;y] \) such that \( z_1 \) and \( z_2 \) are  the complex conjugates.}
{\( [x;y] \in\left\{[4;1],[1;4]\right\} \)}{\( [x;y]\in\left\{[6;1],[9;4]\right\} \)}{\( [x;y]\in\left\{[4;9],[1;6]\right\} \)}{\([x;y]\in\left\{[-4;9],[-1;6]\right\} \)}{\(  [x;y]\in\left\{[6;-1],[9;-4]\right\}  \)}{}}
\LANGpl{
\Question{
Niech \( [x;y]\in\mathbb{R}\times\mathbb{R} \), \( z_1   = 5 + xy\,\mathrm{i} \)  i \( z_2  = x + y - 4\,\mathrm{i} \). Wskaż wszystkie \( [x;y] \) tak, aby \( z_1 \) i \( z_2 \) były sprzężonymi liczbami zespolonymi.}
{\( [x;y] \in\left\{[4;1],[1;4]\right\} \)}{\( [x;y]\in\left\{[6;1],[9;4]\right\} \)}{\( [x;y]\in\left\{[4;9],[1;6]\right\} \)}{\([x;y]\in\left\{[-4;9],[-1;6]\right\} \)}{\(  [x;y]\in\left\{[6;-1],[9;-4]\right\}  \)}{}}
\LANGcs{
\Question{
Určete, pro která \( [x,y] \in \mathbb{R}\times\mathbb{R}\), jsou čísla \(z_1   = 5 + xy\,\mathrm{i} \)  a \( z_2  = x + y - 4\,\mathrm{i} \) komplexně sdružená.}
{\( [x;y] \in\left\{[4;1],[1;4]\right\} \)}{\( [x;y]\in\left\{[6;1],[9;4]\right\} \)}{\( [x;y]\in\left\{[4;9],[1;6]\right\} \)}{\([x;y]\in\left\{[-4;9],[-1;6]\right\} \)}{\(  [x;y]\in\left\{[6;-1],[9;-4]\right\}  \)}{}}
\LANGsk{
\Question{
Nech \( [x;y]\in\mathbb{R}\times\mathbb{R} \), \( z_1   = 5 + xy\,\mathrm{i} \)  a \( z_2  = x + y - 4\,\mathrm{i} \). Určte všetky \( [x;y] \), ak \( z_1 \) a \( z_2 \) sú čísla komplexne združené.}
{\( [x;y] \in\left\{[4;1],[1;4]\right\} \)}{\( [x;y]\in\left\{[6;1],[9;4]\right\} \)}{\( [x;y]\in\left\{[4;9],[1;6]\right\} \)}{\([x;y]\in\left\{[-4;9],[-1;6]\right\} \)}{\(  [x;y]\in\left\{[6;-1],[9;-4]\right\}  \)}{}}
\LANGes{
\Question{
Sean \( [x;y]\in\mathbb{R}\times\mathbb{R} \), \( z_1   = 5 + xy\,\mathrm{i} \)  y \( z_2  = x + y - 4\,\mathrm{i} \). Calcula todos \( [x;y] \) suponiendo que \( z_1 \) y \( z_2 \) son números complejos conjugados.}
{\( [x;y] \in\left\{[4;1],[1;4]\right\} \)}{\( [x;y]\in\left\{[6;1],[9;4]\right\} \)}{\( [x;y]\in\left\{[4;9],[1;6]\right\} \)}{\([x;y]\in\left\{[-4;9],[-1;6]\right\} \)}{\(  [x;y]\in\left\{[6;-1],[9;-4]\right\}  \)}{}}
\End

\Data\ID{1003082401}
\Updated{2021-04-24 05:04:20}
\Level{B}
\SubArea{Complex numbers in algebraic and trigonometric form}
\LANGen{
\Question{
Find the argument of the complex number \(\sqrt2\left(\cos160^{\circ}+\mathrm{i}\cdot\sin160^{\circ}\right) \).}
{\( 160^{\circ} \)}{\( 200^{\circ} \)}{\( \sqrt{2}  \)}{\( 340^{\circ} \)}{}{}}
\LANGpl{
\Question{
Wyznacz argument liczby zespolonej \(\sqrt2\left(\cos160^{\circ}+\mathrm{i}\cdot\sin160^{\circ}\right) \).}
{\( 160^{\circ} \)}{\( 200^{\circ} \)}{\( \sqrt{2}  \)}{\( 340^{\circ} \)}{}{}}
\LANGcs{
\Question{
Určete argument komplexního čísla \(\sqrt2\left(\cos160^{\circ}+\mathrm{i}\cdot\sin160^{\circ}\right) \).}
{\( 160^{\circ} \)}{\( 200^{\circ} \)}{\( \sqrt{2}  \)}{\( 340^{\circ} \)}{}{}}
\LANGsk{
\Question{
Určte argument komplexného čísla \(\sqrt2\left(\cos160^{\circ}+\mathrm{i}\cdot\sin160^{\circ}\right) \).}
{\( 160^{\circ} \)}{\( 200^{\circ} \)}{\( \sqrt{2}  \)}{\( 340^{\circ} \)}{}{}}
\LANGes{
\Question{
Determina el argumento del número complejo \(\sqrt2\left(\cos160^{\circ}+\mathrm{i}\cdot\sin160^{\circ}\right) \).}
{\( 160^{\circ} \)}{\( 200^{\circ} \)}{\( \sqrt{2}  \)}{\( 340^{\circ} \)}{}{}}
\End

\Data\ID{1003082402}
\Updated{2021-04-24 05:04:08}
\Level{B}
\SubArea{Complex numbers in algebraic and trigonometric form}
\LANGen{
\Question{
Find the argument of the complex number \( 4\left(\cos\frac{\pi}3+\mathrm{i}\cdot\sin\frac{\pi}3\right) \).}
{\( \frac{\pi}3  \)}{\( \frac{4\pi}3 \)}{\( 4 \)}{\( \frac{5\pi}3 \)}{}{}}
\LANGpl{
\Question{
Wyznacz argument liczby zespolonej \( 4\left(\cos\frac{\pi}3+\mathrm{i}\cdot\sin\frac{\pi}3\right) \).}
{\( \frac{\pi}3  \)}{\( \frac{4\pi}3 \)}{\( 4 \)}{\( \frac{5\pi}3 \)}{}{}}
\LANGcs{
\Question{
Určete argument komplexního čísla \( 4\left(\cos\frac{\pi}3+\mathrm{i}\cdot\sin\frac{\pi}3\right) \).}
{\( \frac{\pi}3  \)}{\( \frac{4\pi}3 \)}{\( 4 \)}{\( \frac{5\pi}3 \)}{}{}}
\LANGsk{
\Question{
Určte argument komplexného čísla \( 4\left(\cos\frac{\pi}3+\mathrm{i}\cdot\sin\frac{\pi}3\right) \).}
{\( \frac{\pi}3  \)}{\( \frac{4\pi}3 \)}{\( 4 \)}{\( \frac{5\pi}3 \)}{}{}}
\LANGes{
\Question{
Determina el argumento del número complejo  \( 4\left(\cos\frac{\pi}3+\mathrm{i}\cdot\sin\frac{\pi}3\right) \).}
{\( \frac{\pi}3  \)}{\( \frac{4\pi}3 \)}{\( 4 \)}{\( \frac{5\pi}3 \)}{}{}}
\End

\Data\ID{1003082403}
\Updated{2021-04-24 05:04:26}
\Level{B}
\SubArea{Complex numbers in algebraic and trigonometric form}
\LANGen{
\Question{
Find the argument of the complex number \( 2\left(\cos60^{\circ}-\mathrm{i}\cdot\sin60^{\circ}\right) \).}
{\( 300^{\circ} \)}{\( 60^{\circ} \)}{\( 120^{\circ} \)}{\(  240^{\circ} \)}{}{}}
\LANGpl{
\Question{
Wyznacz argument liczby zespolonej \( 2\left(\cos60^{\circ}-\mathrm{i}\cdot\sin60^{\circ}\right) \).}
{\( 300^{\circ} \)}{\( 60^{\circ} \)}{\( 120^{\circ} \)}{\(  240^{\circ} \)}{}{}}
\LANGcs{
\Question{
Určete argument komplexního čísla: \( 2\left(\cos60^{\circ}-\mathrm{i}\cdot\sin60^{\circ}\right) \).}
{\( 300^{\circ} \)}{\( 60^{\circ} \)}{\( 120^{\circ} \)}{\(  240^{\circ} \)}{}{}}
\LANGsk{
\Question{
Určte argument komplexného čísla: \( 2\left(\cos60^{\circ}-\mathrm{i}\cdot\sin60^{\circ}\right) \).}
{\( 300^{\circ} \)}{\( 60^{\circ} \)}{\( 120^{\circ} \)}{\(  240^{\circ} \)}{}{}}
\LANGes{
\Question{
Determina el argumento del número complejo  \( 2\left(\cos60^{\circ}-\mathrm{i}\cdot\sin60^{\circ}\right) \).}
{\( 300^{\circ} \)}{\( 60^{\circ} \)}{\( 120^{\circ} \)}{\(  240^{\circ} \)}{}{}}
\End

\Data\ID{1003082404}
\Updated{2021-04-24 05:04:53}
\Level{B}
\SubArea{Complex numbers in algebraic and trigonometric form}
\LANGen{
\Question{
Find the argument of the complex number \(4\left(\cos\frac{\pi}6-\mathrm{i}\cdot\sin\frac{\pi}6\right) \).}
{\( \frac{11\pi}6 \)}{\( \frac{\pi}6 \)}{\(  \frac{5\pi}6 \)}{\( \frac{7\pi}6 \)}{}{}}
\LANGpl{
\Question{
Wyznacz argument liczby zespolonej \(4\left(\cos\frac{\pi}6-\mathrm{i}\cdot\sin\frac{\pi}6\right) \).}
{\( \frac{11\pi}6 \)}{\( \frac{\pi}6 \)}{\(  \frac{5\pi}6 \)}{\( \frac{7\pi}6 \)}{}{}}
\LANGcs{
\Question{
Určete argument komplexního čísla \(4\left(\cos\frac{\pi}6-\mathrm{i}\cdot\sin\frac{\pi}6\right) \).}
{\( \frac{11\pi}6 \)}{\( \frac{\pi}6 \)}{\(  \frac{5\pi}6 \)}{\( \frac{7\pi}6 \)}{}{}}
\LANGsk{
\Question{
Určte argument komplexného čísla \(4\left(\cos\frac{\pi}6-\mathrm{i}\cdot\sin\frac{\pi}6\right) \).}
{\( \frac{11\pi}6 \)}{\( \frac{\pi}6 \)}{\(  \frac{5\pi}6 \)}{\( \frac{7\pi}6 \)}{}{}}
\LANGes{
\Question{
Determina el argumento del número complejo  \(4\left(\cos\frac{\pi}6-\mathrm{i}\cdot\sin\frac{\pi}6\right) \).}
{\( \frac{11\pi}6 \)}{\( \frac{\pi}6 \)}{\(  \frac{5\pi}6 \)}{\( \frac{7\pi}6 \)}{}{}}
\End

\Data\ID{2010013101}
\Updated{2023-01-23 10:01:02}
\Level{B}
\SubArea{Complex numbers in algebraic and trigonometric form}
\LANGen{
\Question{
Given \(z_{1} = 4\left (\cos \frac{7\pi}
{3} + \mathrm{i}\sin \frac{7\pi}
{3} \right )\) 
and \(z_{2} = \frac12\left (\cos \frac{\pi}
{6} + \mathrm{i}\sin \frac{\pi}
{6} \right )\),
evaluate \(z_{1}\cdot z_{2} \).}
{\( 2\mathrm{i}\)}{\(- 2\mathrm{i}\)}{\(2\)}{\(-2\)}{}{}}
\LANGpl{
\Question{
Mając dane liczby  \(z_{1} = 4\left (\cos \frac{7\pi}
{3} + \mathrm{i}\sin \frac{7\pi}
{3} \right )\) 
oraz  \(z_{2} = \frac12\left (\cos \frac{\pi}
{6} + \mathrm{i}\sin \frac{\pi}
{6} \right )\),
oblicz \(z_{1}\cdot z_{2} \).}
{\( 2\mathrm{i}\)}{\(- 2\mathrm{i}\)}{\(2\)}{\(-2\)}{}{}}
\LANGcs{
\Question{
Jsou dána komplexní čísla \(z_{1} = 4\left (\cos \frac{7\pi}
{3} + \mathrm{i}\sin \frac{7\pi}
{3} \right )\) 
a \(z_{2} = \frac12\left (\cos \frac{\pi}
{6} + \mathrm{i}\sin \frac{\pi}
{6} \right )\). Výraz \(z_{1}\cdot z_{2} \) je roven:}
{\( 2\mathrm{i}\)}{\(- 2\mathrm{i}\)}{\(2\)}{\(-2\)}{}{}}
\LANGsk{
\Question{
Nech \(z_{1} = 4\left (\cos \frac{7\pi}
{3} + \mathrm{i}\sin \frac{7\pi}
{3} \right )\) 
a \(z_{2} = \frac12\left (\cos \frac{\pi}
{6} + \mathrm{i}\sin \frac{\pi}
{6} \right )\),
vypočítajte \(z_{1}\cdot z_{2} \).}
{\( 2\mathrm{i}\)}{\(- 2\mathrm{i}\)}{\(2\)}{\(-2\)}{}{}}
\LANGes{
\Question{
Sean \(z_{1} = 4\left (\cos \frac{7\pi}
{3} + \mathrm{i}\sin \frac{7\pi}
{3} \right )\) 
y \(z_{2} = \frac12\left (\cos \frac{\pi}
{6} + \mathrm{i}\sin \frac{\pi}
{6} \right )\),
calcula \(z_{1}\cdot z_{2} \).}
{\( 2\mathrm{i}\)}{\(- 2\mathrm{i}\)}{\(2\)}{\(-2\)}{}{}}
\End

\Data\ID{2010013102}
\Updated{2023-01-23 10:01:18}
\Level{B}
\SubArea{Complex numbers in algebraic and trigonometric form}
\LANGen{
\Question{
Given the complex numbers  \( a=2\left(\cos \frac{\pi}{3}+\mathrm{i}\sin \frac{\pi}{3}\right) \), \( b=\sqrt{2}\left(\cos \frac{5\pi}{4}+\mathrm{i}\sin \frac{5\pi}{4}\right) \) and \( c=2\sqrt{2}\left(\cos \left(-\frac{\pi}{6}\right)+\mathrm{i}\sin \left(-\frac{\pi}{6}\right)\right) \), evaluate \( a\cdot b\cdot c \).}
{\(8\left(\cos \frac{17\pi}{12}+\mathrm{i}\sin \frac{17\pi}{12} \right) \)}{\(8\left(\cos \frac{17\pi}{12}-\mathrm{i}\sin \frac{17\pi}{12} \right) \)}{\(8\left(\cos \frac{7\pi}{4}+\mathrm{i}\sin \frac{7\pi}{4} \right) \)}{\(4\sqrt{2}\left(\cos \frac{17\pi}{12}+\mathrm{i}\sin \frac{17\pi}{12} \right) \)}{}{}}
\LANGpl{
\Question{
Mając dane liczby zespolone   \( a=2\left(\cos \frac{\pi}{3}+\mathrm{i}\sin \frac{\pi}{3}\right) \), \( b=\sqrt{2}\left(\cos \frac{5\pi}{4}+\mathrm{i}\sin \frac{5\pi}{4}\right) \)  oraz  \( c=2\sqrt{2}\left(\cos \left(-\frac{\pi}{6}\right)+\mathrm{i}\sin \left(-\frac{\pi}{6}\right)\right) \), oblicz  \( a\cdot b\cdot c \).}
{\(8\left(\cos \frac{17\pi}{12}+\mathrm{i}\sin \frac{17\pi}{12} \right) \)}{\(8\left(\cos \frac{17\pi}{12}-\mathrm{i}\sin \frac{17\pi}{12} \right) \)}{\(8\left(\cos \frac{7\pi}{4}+\mathrm{i}\sin \frac{7\pi}{4} \right) \)}{\(4\sqrt{2}\left(\cos \frac{17\pi}{12}+\mathrm{i}\sin \frac{17\pi}{12} \right) \)}{}{}}
\LANGcs{
\Question{
Jsou daná komplexní čísla \( a=2\left(\cos \frac{\pi}{3}+\mathrm{i}\sin \frac{\pi}{3}\right) \), \( b=\sqrt{2}\left(\cos \frac{5\pi}{4}+\mathrm{i}\sin \frac{5\pi}{4}\right) \) a \( c=2\sqrt{2}\left(\cos \left(-\frac{\pi}{6}\right)+\mathrm{i}\sin \left(-\frac{\pi}{6}\right)\right) \). Vypočítejte \( a\cdot b\cdot c \).}
{\(8\left(\cos \frac{17\pi}{12}+\mathrm{i}\sin \frac{17\pi}{12} \right) \)}{\(8\left(\cos \frac{17\pi}{12}-\mathrm{i}\sin \frac{17\pi}{12} \right) \)}{\(8\left(\cos \frac{7\pi}{4}+\mathrm{i}\sin \frac{7\pi}{4} \right) \)}{\(4\sqrt{2}\left(\cos \frac{17\pi}{12}+\mathrm{i}\sin \frac{17\pi}{12} \right) \)}{}{}}
\LANGsk{
\Question{
Vzhľadom na komplexné čísla \( a=2\left(\cos \frac{\pi}{3}+\mathrm{i}\sin \frac{\pi}{3}\right) \), \( b=\sqrt{2}\left(\cos \frac{5\pi}{4}+\mathrm{i}\sin \frac{5\pi}{4}\right) \) a \( c =2\sqrt{2}\left(\cos \left(-\frac{\pi}{6}\right)+\mathrm{i}\sin \left(-\frac{\pi}{6}\right)\right) \), zistite \( a\cdot b\cdot c \).}
{\(8\left(\cos \frac{17\pi}{12}+\mathrm{i}\sin \frac{17\pi}{12} \right) \)}{\(8\left(\cos \frac{17\pi}{12}-\mathrm{i}\sin \frac{17\pi}{12} \right) \)}{\(8\left(\cos \frac{7\pi}{4}+\mathrm{i}\sin \frac{7\pi}{4} \right) \)}{\(4\sqrt{2}\left(\cos \frac{17\pi}{12}+\mathrm{i}\sin \frac{17\pi}{12} \right) \)}{}{}}
\LANGes{
\Question{
Dados los números complejos \( a=2\left(\cos \frac{\pi}{3}+\mathrm{i}\sin \frac{\pi}{3}\right) \), \( b=\sqrt{2}\left(\cos \frac{5\pi}{4}+\mathrm{i}\sin \frac{5\pi}{4}\right) \) y \( c=2\sqrt{2}\left(\cos \left(-\frac{\pi}{6}\right)+\mathrm{i}\sin \left(-\frac{\pi}{6}\right)\right) \), calcula \( a\cdot b\cdot c \).}
{\(8\left(\cos \frac{17\pi}{12}+\mathrm{i}\sin \frac{17\pi}{12} \right) \)}{\(8\left(\cos \frac{17\pi}{12}-\mathrm{i}\sin \frac{17\pi}{12} \right) \)}{\(8\left(\cos \frac{7\pi}{4}+\mathrm{i}\sin \frac{7\pi}{4} \right) \)}{\(4\sqrt{2}\left(\cos \frac{17\pi}{12}+\mathrm{i}\sin \frac{17\pi}{12} \right) \)}{}{}}
\End

\Data\ID{2010013103}
\Updated{2023-01-23 10:01:18}
\Level{B}
\SubArea{Complex numbers in algebraic and trigonometric form}
\LANGen{
\Question{
Given the complex numbers \( a=3\sqrt{2}\left(\cos120^{\circ}+\mathrm{i}\sin120^{\circ}\right) \),  \( b=2\left(\cos\frac{4\pi}3+\mathrm{i}\sin\frac{4\pi}{3}\right) \) and  \( c=\frac{\sqrt{2}}{2}\left(\cos\frac{\pi}4+\mathrm{i}\sin\frac{\pi}{4}\right) \), evaluate \( \frac{a\cdot b}{c} \).}
{\( 12\left(\cos\frac{7\pi}4+\mathrm{i}\sin\frac{7\pi}4\right) \)}{\( 12\left(\sin\frac{7\pi}4+\mathrm{i}\cos\frac{7\pi}4\right) \)}{\( 6\left(\cos\frac{7\pi}4+\mathrm{i}\sin\frac{7\pi}4\right) \)}{\( 12\left(\cos\frac{9\pi}4+\mathrm{i}\sin\frac{9\pi}4\right) \)}{}{}}
\LANGpl{
\Question{
Mając dane liczby zespolone  \( a=3\sqrt{2}\left(\cos120^{\circ}+\mathrm{i}\sin120^{\circ}\right) \),  \( b=2\left(\cos\frac{4\pi}3+\mathrm{i}\sin\frac{4\pi}{3}\right) \)  oraz   \( c=\frac{\sqrt{2}}{2}\left(\cos\frac{\pi}4+\mathrm{i}\sin\frac{\pi}{4}\right) \),  oblicz  \( \frac{a\cdot b}{c} \).}
{\( 12\left(\cos\frac{7\pi}4+\mathrm{i}\sin\frac{7\pi}4\right) \)}{\( 12\left(\sin\frac{7\pi}4+\mathrm{i}\cos\frac{7\pi}4\right) \)}{\( 6\left(\cos\frac{7\pi}4+\mathrm{i}\sin\frac{7\pi}4\right) \)}{\( 12\left(\cos\frac{9\pi}4+\mathrm{i}\sin\frac{9\pi}4\right) \)}{}{}}
\LANGcs{
\Question{
Jsou daná komplexní čísla \( a=3\sqrt{2}\left(\cos120^{\circ}+\mathrm{i}\sin120^{\circ}\right) \),  \( b=2\left(\cos\frac{4\pi}3+\mathrm{i}\sin\frac{4\pi}{3}\right) \) a \( c=\frac{\sqrt{2}}{2}\left(\cos\frac{\pi}4+\mathrm{i}\sin\frac{\pi}{4}\right) \). Vypočítejte \( \frac{a\cdot b}{c} \).}
{\( 12\left(\cos\frac{7\pi}4+\mathrm{i}\sin\frac{7\pi}4\right) \)}{\( 12\left(\sin\frac{7\pi}4+\mathrm{i}\cos\frac{7\pi}4\right) \)}{\( 6\left(\cos\frac{7\pi}4+\mathrm{i}\sin\frac{7\pi}4\right) \)}{\( 12\left(\cos\frac{9\pi}4+\mathrm{i}\sin\frac{9\pi}4\right) \)}{}{}}
\LANGsk{
\Question{
Vzhľadom na komplexné čísla \( a=3\sqrt{2}\left(\cos120^{\circ}+\mathrm{i}\sin120^{\circ}\right) \), \( b=2\left (\cos\frac{4\pi}3+\mathrm{i}\sin\frac{4\pi}{3}\right) \) a \( c=\frac{\sqrt{2}}{2 }\left(\cos\frac{\pi}4+\mathrm{i}\sin\frac{\pi}{4}\right) \), zistite \( \frac{a\cdot b}{c} \).}
{\( 12\left(\cos\frac{7\pi}4+\mathrm{i}\sin\frac{7\pi}4\right) \)}{\( 12\left(\sin\frac{7\pi}4+\mathrm{i}\cos\frac{7\pi}4\right) \)}{\( 6\left(\cos\frac{7\pi}4+\mathrm{i}\sin\frac{7\pi}4\right) \)}{\( 12\left(\cos\frac{9\pi}4+\mathrm{i}\sin\frac{9\pi}4\right) \)}{}{}}
\LANGes{
\Question{
Dados los números complejos \( a=3\sqrt{2}\left(\cos120^{\circ}+\mathrm{i}\sin120^{\circ}\right) \),  \( b=2\left(\cos\frac{4\pi}3+\mathrm{i}\sin\frac{4\pi}{3}\right) \)  y  \( c=\frac{\sqrt{2}}{2}\left(\cos\frac{\pi}4+\mathrm{i}\sin\frac{\pi}{4}\right) \),  calcula \( \frac{a\cdot b}{c} \).}
{\( 12\left(\cos\frac{7\pi}4+\mathrm{i}\sin\frac{7\pi}4\right) \)}{\( 12\left(\sin\frac{7\pi}4+\mathrm{i}\cos\frac{7\pi}4\right) \)}{\( 6\left(\cos\frac{7\pi}4+\mathrm{i}\sin\frac{7\pi}4\right) \)}{\( 12\left(\cos\frac{9\pi}4+\mathrm{i}\sin\frac{9\pi}4\right) \)}{}{}}
\End

\Data\ID{2010013104}
\Updated{2023-01-23 10:01:18}
\Level{B}
\SubArea{Complex numbers in algebraic and trigonometric form}
\LANGen{
\Question{
Let \( [x;y]\in\mathbb{R}\times\mathbb{R} \), \( z_1   = -2 + xy\,\mathrm{i} \)  and \( z_2  = x + y + 8\,\mathrm{i}\). Find all \( [x;y] \) such that \( z_1 \) and \( z_2 \) are  the opposite numbers.}
{\( [x;y] \in\left\{[4;-2],[-2;4]\right\} \)}{\( [x;y] \in \left\{[4;-2]\right\} \)}{\( [x;y] \in\left\{[-4;2],[2;-4]\right\} \)}{\( [x;y] \in \left\{[-2;4]\right\} \)}{}{}}
\LANGpl{
\Question{
Niech  \( [x;y]\in\mathbb{R}\times\mathbb{R} \), \( z_1   = -2 + xy\,\mathrm{i} \)  i  \( z_2  = x + y + 8\,\mathrm{i}\). Znajdź wszystkie  \( [x;y] \) takie, że  \( z_1 \) i \( z_2 \) są liczbami przeciwnymi.}
{\( [x;y] \in\left\{[4;-2],[-2;4]\right\} \)}{\( [x;y] \in \left\{[4;-2]\right\} \)}{\( [x;y] \in\left\{[-4;2],[2;-4]\right\} \)}{\( [x;y] \in \left\{[-2;4]\right\} \)}{}{}}
\LANGcs{
\Question{
Určete, pro která \( [x;y]\in\mathbb{R}\times\mathbb{R} \), jsou čísla \( z_1   = -2 + xy\,\mathrm{i} \) a \( z_2  = x + y + 8\,\mathrm{i}\) opačná.}
{\( [x;y] \in\left\{[4;-2],[-2;4]\right\} \)}{\( [x;y] \in \left\{[4;-2]\right\} \)}{\( [x;y] \in\left\{[-4;2],[2;-4]\right\} \)}{\( [x;y] \in \left\{[-2;4]\right\} \)}{}{}}
\LANGsk{
\Question{
Ak \( [x;y]\in\mathbb{R}\times\mathbb{R} \), \( z_1   = -2 + xy\,\mathrm{i} \)  a \( z_2  = x + y + 8\,\mathrm{i}\). Nájdite všetky \( [x;y] \) tak že \( z_1 \) and \( z_2 \) sú opačné čísla.}
{\( [x;y] \in\left\{[4;-2],[-2;4]\right\} \)}{\( [x;y] \in \left\{[4;-2]\right\} \)}{\( [x;y] \in\left\{[-4;2],[2;-4]\right\} \)}{\( [x;y] \in \left\{[-2;4]\right\} \)}{}{}}
\LANGes{
\Question{
Sean \( [x;y]\in\mathbb{R}\times\mathbb{R} \), \( z_1   = -2 + xy\,\mathrm{i} \)  y \( z_2  = x + y + 8\,\mathrm{i}\). Halla todos los \( [x;y] \) tales que \( z_1 \) y \( z_2 \) sean números opuestos.}
{\( [x;y] \in\left\{[4;-2],[-2;4]\right\} \)}{\( [x;y] \in \left\{[4;-2]\right\} \)}{\( [x;y] \in\left\{[-4;2],[2;-4]\right\} \)}{\( [x;y] \in \left\{[-2;4]\right\} \)}{}{}}
\End

\Data\ID{2010013105}
\Updated{2023-01-23 10:01:18}
\Level{B}
\SubArea{Complex numbers in algebraic and trigonometric form}
\LANGen{
\Question{
Let \(z_1 = \sqrt3 + \mathrm{i}\), \(z_2=1 + \mathrm{i}\sqrt3\). Identify a complex number that is not equal to \(\frac{z_1}{z_2}\).}
{\(\cos \frac{5\pi}{6} + \mathrm{i} \sin \frac{5\pi}{6}\)}{\(\cos \left(-\frac{\pi}{6}\right) + \mathrm{i} \sin \left(-\frac{\pi}{6}\right)\)}{\( \frac{\sqrt{3}}{2} - \frac{\mathrm{i}}{2}\)}{\(\cos \frac{\pi}{6} - \mathrm{i} \sin \frac{\pi}{6}\)}{}{}}
\LANGpl{
\Question{
Niech  \(z_1 = \sqrt3 + \mathrm{i}\), \(z_2=1 + \mathrm{i}\sqrt3\). Wskaż liczbę, która nie jest równa  \(\frac{z_1}{z_2}\).}
{\(\cos \frac{5\pi}{6} + \mathrm{i} \sin \frac{5\pi}{6}\)}{\(\cos \left(-\frac{\pi}{6}\right) + \mathrm{i} \sin \left(-\frac{\pi}{6}\right)\)}{\( \frac{\sqrt{3}}{2} - \frac{\mathrm{i}}{2}\)}{\(\cos \frac{\pi}{6} - \mathrm{i} \sin \frac{\pi}{6}\)}{}{}}
\LANGcs{
\Question{
Jsou dána čísla \(z_1 = \sqrt3 + \mathrm{i}\), \(z_2=1 + \mathrm{i}\sqrt3\). Které z následujících komplexních čísel se nerovná \(\frac{z_1}{z_2}\).}
{\(\cos \frac{5\pi}{6} + \mathrm{i} \sin \frac{5\pi}{6}\)}{\(\cos \left(-\frac{\pi}{6}\right) + \mathrm{i} \sin \left(-\frac{\pi}{6}\right)\)}{\( \frac{\sqrt{3}}{2} - \frac{\mathrm{i}}{2}\)}{\(\cos \frac{\pi}{6} - \mathrm{i} \sin \frac{\pi}{6}\)}{}{}}
\LANGsk{
\Question{
Nech \(z_1 = \sqrt3 + \mathrm{i}\), \(z_2=1 + \mathrm{i}\sqrt3\). Identifikujte komplexné číslo, ktoré sa nerovná \(\frac{z_1}{z_2}\).}
{\(\cos \frac{5\pi}{6} + \mathrm{i} \sin \frac{5\pi}{6}\)}{\(\cos \left(-\frac{\pi}{6}\right) + \mathrm{i} \sin \left(-\frac{\pi}{6}\right)\)}{\( \frac{\sqrt{3}}{2} - \frac{\mathrm{i}}{2}\)}{\(\cos \frac{\pi}{6} - \mathrm{i} \sin \frac{\pi}{6}\)}{}{}}
\LANGes{
\Question{
Sean \(z_1 = \sqrt3 + \mathrm{i}\), \(z_2=1 + \mathrm{i}\sqrt3\). Identifica el número complejo que es distinto a \(\frac{z_1}{z_2}\).}
{\(\cos \frac{5\pi}{6} + \mathrm{i} \sin \frac{5\pi}{6}\)}{\(\cos \left(-\frac{\pi}{6}\right) + \mathrm{i} \sin \left(-\frac{\pi}{6}\right)\)}{\( \frac{\sqrt{3}}{2} - \frac{\mathrm{i}}{2}\)}{\(\cos \frac{\pi}{6} - \mathrm{i} \sin \frac{\pi}{6}\)}{}{}}
\End

\Data\ID{2010013106}
\Updated{2023-01-23 10:01:18}
\Level{B}
\SubArea{Complex numbers in algebraic and trigonometric form}
\LANGen{
\Question{
Let \(z_1 = 1 + \mathrm{i}\sqrt{3}\), \(z_2=\sqrt3 + \mathrm{i}\). Identify a complex number that is not equal to \(\frac{z_1}{z_2}\).}
{\(\cos \frac{7\pi}{6} + \mathrm{i} \sin \frac{7\pi}{6}\)}{\(\cos \frac{\pi}{6} +\mathrm{i} \sin \frac{\pi}{6}\)}{\( \frac{\sqrt{3}}{2} + \frac{\mathrm{i}}{2}\)}{\(\cos \left(-\frac{\pi}{6}\right) - \mathrm{i} \sin \left(-\frac{\pi}{6}\right)\)}{}{}}
\LANGpl{
\Question{
Niech  \(z_1 = 1 + \mathrm{i}\sqrt{3}\), \(z_2=\sqrt3 + \mathrm{i}\). Wskaż liczbę, która nie jest równa \(\frac{z_1}{z_2}\).}
{\(\cos \frac{7\pi}{6} + \mathrm{i} \sin \frac{7\pi}{6}\)}{\(\cos \frac{\pi}{6} +\mathrm{i} \sin \frac{\pi}{6}\)}{\( \frac{\sqrt{3}}{2} + \frac{\mathrm{i}}{2}\)}{\(\cos \left(-\frac{\pi}{6}\right) - \mathrm{i} \sin \left(-\frac{\pi}{6}\right)\)}{}{}}
\LANGcs{
\Question{
Jsou dána čísla \(z_1 = 1 + \mathrm{i}\sqrt{3}\), \(z_2=\sqrt3 + \mathrm{i}\). Které z následujících komplexních čísel se nerovná \(\frac{z_1}{z_2}\).}
{\(\cos \frac{7\pi}{6} + \mathrm{i} \sin \frac{7\pi}{6}\)}{\(\cos \frac{\pi}{6} +\mathrm{i} \sin \frac{\pi}{6}\)}{\( \frac{\sqrt{3}}{2} + \frac{\mathrm{i}}{2}\)}{\(\cos \left(-\frac{\pi}{6}\right) - \mathrm{i} \sin \left(-\frac{\pi}{6}\right)\)}{}{}}
\LANGsk{
\Question{
Nech \(z_1 = 1 + \mathrm{i}\sqrt{3}\), \(z_2=\sqrt3 + \mathrm{i}\). Identifikujte komplexné číslo, ktoré sa nerovná \(\frac{z_1}{z_2}\).}
{\(\cos \frac{7\pi}{6} + \mathrm{i} \sin \frac{7\pi}{6}\)}{\(\cos \frac{\pi}{6} +\mathrm{i} \sin \frac{\pi}{6}\)}{\( \frac{\sqrt{3}}{2} + \frac{\mathrm{i}}{2}\)}{\(\cos \left(-\frac{\pi}{6}\right) - \mathrm{i} \sin \left(-\frac{\pi}{6}\right)\)}{}{}}
\LANGes{
\Question{
Sean \(z_1 = 1 + \mathrm{i}\sqrt{3}\), \(z_2=\sqrt3 + \mathrm{i}\). Identifica el número complejo que es distinto a \(\frac{z_1}{z_2}\).}
{\(\cos \frac{7\pi}{6} + \mathrm{i} \sin \frac{7\pi}{6}\)}{\(\cos \frac{\pi}{6} +\mathrm{i} \sin \frac{\pi}{6}\)}{\( \frac{\sqrt{3}}{2} + \frac{\mathrm{i}}{2}\)}{\(\cos \left(-\frac{\pi}{6}\right) - \mathrm{i} \sin \left(-\frac{\pi}{6}\right)\)}{}{}}
\End

\Data\ID{2010013109}
\Updated{2023-01-23 10:01:18}
\Level{B}
\SubArea{Complex numbers in algebraic and trigonometric form}
\LANGen{
\Question{
Let \(z=\frac{1}{\cos \frac{2\pi}{3} + \mathrm{i} \sin\frac{2\pi}{3}}\). Which of the following numbers is not the value of the argument of \(z\)?}
{\(120^\circ\)}{\(-\frac{2\pi}{3}\)}{\(\frac{4\pi}{3}\)}{\(240^\circ\)}{}{}}
\LANGpl{
\Question{
Niech  \(z=\frac{1}{\cos \frac{2\pi}{3} + \mathrm{i} \sin\frac{2\pi}{3}}\). Która z poniższych liczb nie jest wartością argumentu \(z\)?}
{\(120^\circ\)}{\(-\frac{2\pi}{3}\)}{\(\frac{4\pi}{3}\)}{\(240^\circ\)}{}{}}
\LANGcs{
\Question{
Nechť \(z=\frac{1}{\cos \frac{2\pi}{3} + \mathrm{i} \sin\frac{2\pi}{3}}\). Které z následujících čísel není hodnotou argumentu \(z\)?}
{\(120^\circ\)}{\(-\frac{2\pi}{3}\)}{\(\frac{4\pi}{3}\)}{\(240^\circ\)}{}{}}
\LANGsk{
\Question{
Nech \(z=\frac{1}{\cos \frac{2\pi}{3} + \mathrm{i} \sin\frac{2\pi}{3}}\). Ktoré z nasledujúcich čísel nie je hodnotou argumentu \(z\)?}
{\(120^\circ\)}{\(-\frac{2\pi}{3}\)}{\(\frac{4\pi}{3}\)}{\(240^\circ\)}{}{}}
\LANGes{
\Question{
Sea \(z=\frac{1}{\cos \frac{2\pi}{3} + \mathrm{i} \sin\frac{2\pi}{3}}\). ¿Cuál de los siguientes valores no corresponde al argumento de \(z\)?}
{\(120^\circ\)}{\(-\frac{2\pi}{3}\)}{\(\frac{4\pi}{3}\)}{\(240^\circ\)}{}{}}
\End

\Data\ID{2010013110}
\Updated{2023-01-23 10:01:18}
\Level{B}
\SubArea{Complex numbers in algebraic and trigonometric form}
\LANGen{
\Question{
Let \(z=\frac{1}{\cos \frac{7\pi}{6} + \mathrm{i} \sin\frac{7\pi}{6}}\). Which of the following numbers is not the value of the argument of \(z\)?}
{\(210^\circ\)}{\(150^\circ\)}{\(-\frac{7\pi}{6}\)}{\(\frac{5\pi}{6}\)}{}{}}
\LANGpl{
\Question{
Niech  \(z=\frac{1}{\cos \frac{7\pi}{6} + \mathrm{i} \sin\frac{7\pi}{6}}\). Która z poniższych liczb nie jest wartością argumentu \(z\)?}
{\(210^\circ\)}{\(150^\circ\)}{\(-\frac{7\pi}{6}\)}{\(\frac{5\pi}{6}\)}{}{}}
\LANGcs{
\Question{
Nechť \(z=\frac{1}{\cos \frac{7\pi}{6} + \mathrm{i} \sin\frac{7\pi}{6}}\). Které z následujících čísel není hodnotou argumentu \(z\)?}
{\(210^\circ\)}{\(150^\circ\)}{\(-\frac{7\pi}{6}\)}{\(\frac{5\pi}{6}\)}{}{}}
\LANGsk{
\Question{
Nech \(z=\frac{1}{\cos \frac{7\pi}{6} + \mathrm{i} \sin\frac{7\pi}{6}}\). Ktoré z nasledujúcich čísel nie je hodnotou argumentu \(z\)?}
{\(210^\circ\)}{\(150^\circ\)}{\(-\frac{7\pi}{6}\)}{\(\frac{5\pi}{6}\)}{}{}}
\LANGes{
\Question{
Sea \(z=\frac{1}{\cos \frac{7\pi}{6} + \mathrm{i} \sin\frac{7\pi}{6}}\). ¿Cuál de los siguientes valores no corresponde al argumento de \(z\)?}
{\(210^\circ\)}{\(150^\circ\)}{\(-\frac{7\pi}{6}\)}{\(\frac{5\pi}{6}\)}{}{}}
\End

\Data\ID{2010013111}
\Updated{2023-01-23 10:01:18}
\Level{B}
\SubArea{Complex numbers in algebraic and trigonometric form}
\LANGen{
\Question{
Find the polar form of the complex number
\(z=\frac{\mathrm{i}^{12}+1}
 {\mathrm{i}^{11}+1} \).}
{\(\sqrt{2}\left (\cos \frac{\pi }
{4} + \mathrm{i}\sin \frac{\pi }
{4}\right )\)}{\(\cos \frac{\pi }
{2} + \mathrm{i}\sin \frac{\pi }
{2}\)}{\(\sqrt{2}\left (\cos \frac{5\pi }
{4} + \mathrm{i}\sin \frac{5\pi }
{4}\right )\)}{\(\sqrt{2}\left (\cos \frac{3\pi }
{4} + \mathrm{i}\sin \frac{3\pi }
{4}\right )\)}{}{}}
\LANGpl{
\Question{
Znajdź trygonometryczną postać liczby zespolonej
\(z=\frac{\mathrm{i}^{12}+1}
 {\mathrm{i}^{11}+1} \).}
{\(\sqrt{2}\left (\cos \frac{\pi }
{4} + \mathrm{i}\sin \frac{\pi }
{4}\right )\)}{\(\cos \frac{\pi }
{2} + \mathrm{i}\sin \frac{\pi }
{2}\)}{\(\sqrt{2}\left (\cos \frac{5\pi }
{4} + \mathrm{i}\sin \frac{5\pi }
{4}\right )\)}{\(\sqrt{2}\left (\cos \frac{3\pi }
{4} + \mathrm{i}\sin \frac{3\pi }
{4}\right )\)}{}{}}
\LANGcs{
\Question{
Goniometrický tvar komplexního čísla
\(z=\frac{\mathrm{i}^{12}+1}
 {\mathrm{i}^{11}+1} \) je:}
{\(\sqrt{2}\left (\cos \frac{\pi }
{4} + \mathrm{i}\sin \frac{\pi }
{4}\right )\)}{\(\cos \frac{\pi }
{2} + \mathrm{i}\sin \frac{\pi }
{2}\)}{\(\sqrt{2}\left (\cos \frac{5\pi }
{4} + \mathrm{i}\sin \frac{5\pi }
{4}\right )\)}{\(\sqrt{2}\left (\cos \frac{3\pi }
{4} + \mathrm{i}\sin \frac{3\pi }
{4}\right )\)}{}{}}
\LANGsk{
\Question{
Nájdite goniometrický tvar komplexného čísla
\(z=\frac{\mathrm{i}^{12}+1}
  {\mathrm{i}^{11}+1} \).}
{\(\sqrt{2}\left (\cos \frac{\pi }
{4} + \mathrm{i}\sin \frac{\pi }
{4}\right )\)}{\(\cos \frac{\pi }
{2} + \mathrm{i}\sin \frac{\pi }
{2}\)}{\(\sqrt{2}\left (\cos \frac{5\pi }
{4} + \mathrm{i}\sin \frac{5\pi }
{4}\right )\)}{\(\sqrt{2}\left (\cos \frac{3\pi }
{4} + \mathrm{i}\sin \frac{3\pi }
{4}\right )\)}{}{}}
\LANGes{
\Question{
Halla la forma polar del número complejo
\(z=\frac{\mathrm{i}^{12}+1}
 {\mathrm{i}^{11}+1} \).}
{\(\sqrt{2}\left (\cos \frac{\pi }
{4} + \mathrm{i}\sin \frac{\pi }
{4}\right )\)}{\(\cos \frac{\pi }
{2} + \mathrm{i}\sin \frac{\pi }
{2}\)}{\(\sqrt{2}\left (\cos \frac{5\pi }
{4} + \mathrm{i}\sin \frac{5\pi }
{4}\right )\)}{\(\sqrt{2}\left (\cos \frac{3\pi }
{4} + \mathrm{i}\sin \frac{3\pi }
{4}\right )\)}{}{}}
\End

\Data\ID{2010013112}
\Updated{2023-01-23 10:01:22}
\Level{B}
\SubArea{Complex numbers in algebraic and trigonometric form}
\LANGen{
\Question{
Find the polar form of the complex conjugate of \(z=-\frac{\sqrt5}{2} + \mathrm{i}\frac{\sqrt{15}}{2}\).}
{\(\sqrt{5}\left (\cos \left(-\frac{2\pi}
{3}\right) + \mathrm{i}\sin \left(-\frac{2\pi}
{3}\right )\right)\)}{\(\sqrt{5}\left (\cos \frac{2\pi }
{3} + \mathrm{i}\sin \frac{2\pi }
{3}\right )\)}{\(\sqrt{10}\left (\cos \frac{4\pi }
{3} + \mathrm{i}\sin \frac{4\pi }
{3}\right )\)}{\(\sqrt{5}\left (\cos \frac{\pi }
{3} + \mathrm{i}\sin \frac{\pi }
{3}\right )\)}{}{}}
\LANGpl{
\Question{
Znajdź trygonometryczną  postać liczby zespolonej  \(z=-\frac{\sqrt5}{2} + \mathrm{i}\frac{\sqrt{15}}{2}\).}
{\(\sqrt{5}\left (\cos \left(-\frac{2\pi}
{3}\right) + \mathrm{i}\sin \left(-\frac{2\pi}
{3}\right )\right)\)}{\(\sqrt{5}\left (\cos \frac{2\pi }
{3} + \mathrm{i}\sin \frac{2\pi }
{3}\right )\)}{\(\sqrt{10}\left (\cos \frac{4\pi }
{3} + \mathrm{i}\sin \frac{4\pi }
{3}\right )\)}{\(\sqrt{5}\left (\cos \frac{\pi }
{3} + \mathrm{i}\sin \frac{\pi }
{3}\right )\)}{}{}}
\LANGcs{
\Question{
Najděte goniometrický tvar komplexně sdruženého čísla k \(z=-\frac{\sqrt5}{2} + \mathrm{i}\frac{\sqrt{15}}{2}\).}
{\(\sqrt{5}\left (\cos \left(-\frac{2\pi}
{3}\right) + \mathrm{i}\sin \left(-\frac{2\pi}
{3}\right )\right)\)}{\(\sqrt{5}\left (\cos \frac{2\pi }
{3} + \mathrm{i}\sin \frac{2\pi }
{3}\right )\)}{\(\sqrt{10}\left (\cos \frac{4\pi }
{3} + \mathrm{i}\sin \frac{4\pi }
{3}\right )\)}{\(\sqrt{5}\left (\cos \frac{\pi }
{3} + \mathrm{i}\sin \frac{\pi }
{3}\right )\)}{}{}}
\LANGsk{
\Question{
Nájdite goniometrický tvar komplexne združeného čísla k \(z=-\frac{\sqrt5}{2} + \mathrm{i}\frac{\sqrt{15}}{2}\).}
{\(\sqrt{5}\left (\cos \left(-\frac{2\pi}
{3}\right) + \mathrm{i}\sin \left(-\frac{2\pi}
{3}\right )\right)\)}{\(\sqrt{5}\left (\cos \frac{2\pi }
{3} + \mathrm{i}\sin \frac{2\pi }
{3}\right )\)}{\(\sqrt{10}\left (\cos \frac{4\pi }
{3} + \mathrm{i}\sin \frac{4\pi }
{3}\right )\)}{\(\sqrt{5}\left (\cos \frac{\pi }
{3} + \mathrm{i}\sin \frac{\pi }
{3}\right )\)}{}{}}
\LANGes{
\Question{
Halla la forma polar del conjugado del número complejo \(z=-\frac{\sqrt5}{2} + \mathrm{i}\frac{\sqrt{15}}{2}\).}
{\(\sqrt{5}\left (\cos \left(-\frac{2\pi}
{3}\right) + \mathrm{i}\sin \left(-\frac{2\pi}
{3}\right )\right)\)}{\(\sqrt{5}\left (\cos \frac{2\pi }
{3} + \mathrm{i}\sin \frac{2\pi }
{3}\right )\)}{\(\sqrt{10}\left (\cos \frac{4\pi }
{3} + \mathrm{i}\sin \frac{4\pi }
{3}\right )\)}{\(\sqrt{5}\left (\cos \frac{\pi }
{3} + \mathrm{i}\sin \frac{\pi }
{3}\right )\)}{}{}}
\End

\Data\ID{2010013113}
\Updated{2023-01-23 10:01:22}
\Level{B}
\SubArea{Complex numbers in algebraic and trigonometric form}
\LANGen{
\Question{
Find the polar form of the opposite number of \(z=-\frac{\sqrt5}{2} + \mathrm{i}\frac{\sqrt{15}}{2}\).}
{\(\sqrt{5}\left (\cos \left(-\frac{\pi}
{3}\right) + \mathrm{i}\sin \left(-\frac{\pi}
{3}\right )\right)\)}{\(\sqrt{5}\left (\cos \frac{2\pi }
{3} + \mathrm{i}\sin \frac{2\pi }
{3}\right )\)}{\(\sqrt{10}\left (\cos \frac{5\pi }
{3} + \mathrm{i}\sin \frac{5\pi }
{3}\right )\)}{\(\sqrt{5}\left (\cos \frac{\pi }
{3} + \mathrm{i}\sin \frac{\pi }
{3}\right )\)}{}{}}
\LANGpl{
\Question{
Znajdź postać trygonometryczną  liczby przeciwnej  do liczby \(z=-\frac{\sqrt5}{2} + \mathrm{i}\frac{\sqrt{15}}{2}\).}
{\(\sqrt{5}\left (\cos \left(-\frac{\pi}
{3}\right) + \mathrm{i}\sin \left(-\frac{\pi}
{3}\right )\right)\)}{\(\sqrt{5}\left (\cos \frac{2\pi }
{3} + \mathrm{i}\sin \frac{2\pi }
{3}\right )\)}{\(\sqrt{10}\left (\cos \frac{5\pi }
{3} + \mathrm{i}\sin \frac{5\pi }
{3}\right )\)}{\(\sqrt{5}\left (\cos \frac{\pi }
{3} + \mathrm{i}\sin \frac{\pi }
{3}\right )\)}{}{}}
\LANGcs{
\Question{
Najděte goniometrický tvar komplexního čísla opačného k \(z=-\frac{\sqrt5}{2} + \mathrm{i}\frac{\sqrt{15}}{2}\).}
{\(\sqrt{5}\left (\cos \left(-\frac{\pi}
{3}\right) + \mathrm{i}\sin \left(-\frac{\pi}
{3}\right )\right)\)}{\(\sqrt{5}\left (\cos \frac{2\pi }
{3} + \mathrm{i}\sin \frac{2\pi }
{3}\right )\)}{\(\sqrt{10}\left (\cos \frac{5\pi }
{3} + \mathrm{i}\sin \frac{5\pi }
{3}\right )\)}{\(\sqrt{5}\left (\cos \frac{\pi }
{3} + \mathrm{i}\sin \frac{\pi }
{3}\right )\)}{}{}}
\LANGsk{
\Question{
Nájdite goniometrický tvar komplexného čísla opačného k \(z=-\frac{\sqrt5}{2} + \mathrm{i}\frac{\sqrt{15}}{2}\).}
{\(\sqrt{5}\left (\cos \left(-\frac{\pi}
{3}\right) + \mathrm{i}\sin \left(-\frac{\pi}
{3}\right )\right)\)}{\(\sqrt{5}\left (\cos \frac{2\pi }
{3} + \mathrm{i}\sin \frac{2\pi }
{3}\right )\)}{\(\sqrt{10}\left (\cos \frac{5\pi }
{3} + \mathrm{i}\sin \frac{5\pi }
{3}\right )\)}{\(\sqrt{5}\left (\cos \frac{\pi }
{3} + \mathrm{i}\sin \frac{\pi }
{3}\right )\)}{}{}}
\LANGes{
\Question{
Halla la forma polar del opuesto del número complejo \(z=-\frac{\sqrt5}{2} + \mathrm{i}\frac{\sqrt{15}}{2}\).}
{\(\sqrt{5}\left (\cos \left(-\frac{\pi}
{3}\right) + \mathrm{i}\sin \left(-\frac{\pi}
{3}\right )\right)\)}{\(\sqrt{5}\left (\cos \frac{2\pi }
{3} + \mathrm{i}\sin \frac{2\pi }
{3}\right )\)}{\(\sqrt{10}\left (\cos \frac{5\pi }
{3} + \mathrm{i}\sin \frac{5\pi }
{3}\right )\)}{\(\sqrt{5}\left (\cos \frac{\pi }
{3} + \mathrm{i}\sin \frac{\pi }
{3}\right )\)}{}{}}
\End

\Data\ID{9000031207}
\Updated{2020-07-18 11:07:34}
\Level{B}
\SubArea{Complex numbers in algebraic and trigonometric form}
\LANGen{
\Question{
Find the algebraic form of the complex number
\(z = 2\left (\cos \frac{3\pi }
{4} + \mathrm{i}\sin \frac{3\pi }
{4}\right )\).}
{\(-\sqrt{2} + \mathrm{i}\sqrt{2}\)}{\(\sqrt{2} + \mathrm{i}\sqrt{2}\)}{\(\sqrt{2} -\mathrm{i}\sqrt{2}\)}{\(-\sqrt{2} -\mathrm{i}\sqrt{2}\)}{}{}}
\LANGpl{
\Question{
Wskaż postać algebraiczną liczby zespolonej
\(z = 2\left (\cos \frac{3\pi }
{4} + \mathrm{i}\sin \frac{3\pi }
{4}\right )\).}
{\(-\sqrt{2} + \mathrm{i}\sqrt{2}\)}{\(\sqrt{2} + \mathrm{i}\sqrt{2}\)}{\(\sqrt{2} -\mathrm{i}\sqrt{2}\)}{\(-\sqrt{2} -\mathrm{i}\sqrt{2}\)}{}{}}
\LANGcs{
\Question{
Vyjádřete komplexní číslo
\(z =\, \) \(2\left (\cos \frac{3\pi }
{4} + \mathrm{i}\sin \frac{3\pi }
{4}\right )\) 
v algebraickém tvaru.}
{\(-\sqrt{2} + \mathrm{i}\sqrt{2}\)}{\(\sqrt{2} + \mathrm{i}\sqrt{2}\)}{\(\sqrt{2} -\mathrm{i}\sqrt{2}\)}{\(-\sqrt{2} -\mathrm{i}\sqrt{2}\)}{}{}}
\LANGsk{
\Question{
Vyjadrite komplexné číslo
\(z = 2\left (\cos \frac{3\pi }
{4} + \mathrm{i}\sin \frac{3\pi }
{4}\right )\) v algebraickom tvare.}
{\(-\sqrt{2} + \mathrm{i}\sqrt{2}\)}{\(\sqrt{2} + \mathrm{i}\sqrt{2}\)}{\(\sqrt{2} -\mathrm{i}\sqrt{2}\)}{\(-\sqrt{2} -\mathrm{i}\sqrt{2}\)}{}{}}
\LANGes{
\Question{
Determina la forma algebraica del número complejo
\(z = 2\left (\cos \frac{3\pi }
{4} + \mathrm{i}\sin \frac{3\pi }
{4}\right )\).}
{\(-\sqrt{2} + \mathrm{i}\sqrt{2}\)}{\(\sqrt{2} + \mathrm{i}\sqrt{2}\)}{\(\sqrt{2} -\mathrm{i}\sqrt{2}\)}{\(-\sqrt{2} -\mathrm{i}\sqrt{2}\)}{}{}}
\End

\Data\ID{9000031208}
\Updated{2020-12-03 06:12:17}
\Level{B}
\SubArea{Complex numbers in algebraic and trigonometric form}
\LANGen{
\Question{
Find the polar form of the complex number
\(z = -3 + 3\mathrm{i}\).}
{\(3\sqrt{2}\left (\cos \frac{3\pi }
{4} + \mathrm{i}\sin \frac{3\pi }
{4}\right )\)}{\(3\left (\cos \frac{\pi }{4} + \mathrm{i}\sin \frac{\pi }{4}\right )\)}{\(3\left (\cos \frac{5\pi }
{4} + \mathrm{i}\sin \frac{5\pi }
{4}\right )\)}{\(3\sqrt{2}\left (\cos \frac{7\pi }
{4} + \mathrm{i}\sin \frac{7\pi }
{4}\right )\)}{}{}}
\LANGpl{
\Question{
Wskaż postać biegunową  liczby zespolonej
\(z = -3 + 3\mathrm{i}\).}
{\(3\sqrt{2}\left (\cos \frac{3\pi }
{4} + \mathrm{i}\sin \frac{3\pi }
{4}\right )\)}{\(3\left (\cos \frac{\pi }{4} + \mathrm{i}\sin \frac{\pi }{4}\right )\)}{\(3\left (\cos \frac{5\pi }
{4} + \mathrm{i}\sin \frac{5\pi }
{4}\right )\)}{\(3\sqrt{2}\left (\cos \frac{7\pi }
{4} + \mathrm{i}\sin \frac{7\pi }
{4}\right )\)}{}{}}
\LANGcs{
\Question{
Vyjádřete komplexní číslo \(z = -3 + 3\mathrm{i}\) 
v goniometrickém tvaru.}
{\(3\sqrt{2}\left (\cos \frac{3\pi }
{4} + \mathrm{i}\sin \frac{3\pi }
{4}\right )\)}{\(3\left (\cos \frac{\pi }{4} + \mathrm{i}\sin \frac{\pi }{4}\right )\)}{\(3\left (\cos \frac{5\pi }
{4} + \mathrm{i}\sin \frac{5\pi }
{4}\right )\)}{\(3\sqrt{2}\left (\cos \frac{7\pi }
{4} + \mathrm{i}\sin \frac{7\pi }
{4}\right )\)}{}{}}
\LANGsk{
\Question{
Vyjadrite komplexné číslo
\(z = -3 + 3\mathrm{i}\) v goniometrickom tvare.}
{\(3\sqrt{2}\left (\cos \frac{3\pi }
{4} + \mathrm{i}\sin \frac{3\pi }
{4}\right )\)}{\(3\left (\cos \frac{\pi }{4} + \mathrm{i}\sin \frac{\pi }{4}\right )\)}{\(3\left (\cos \frac{5\pi }
{4} + \mathrm{i}\sin \frac{5\pi }
{4}\right )\)}{\(3\sqrt{2}\left (\cos \frac{7\pi }
{4} + \mathrm{i}\sin \frac{7\pi }
{4}\right )\)}{}{}}
\LANGes{
\Question{
Determina la forma polar del complejo
\(z = -3 + 3\mathrm{i}\).}
{\(3\sqrt{2}\left (\cos \frac{3\pi }
{4} + \mathrm{i}\sin \frac{3\pi }
{4}\right )\)}{\(3\left (\cos \frac{\pi }{4} + \mathrm{i}\sin \frac{\pi }{4}\right )\)}{\(3\left (\cos \frac{5\pi }
{4} + \mathrm{i}\sin \frac{5\pi }
{4}\right )\)}{\(3\sqrt{2}\left (\cos \frac{7\pi }
{4} + \mathrm{i}\sin \frac{7\pi }
{4}\right )\)}{}{}}
\End

\Data\ID{9000031209}
\Updated{2020-07-18 09:07:23}
\Level{B}
\SubArea{Complex numbers in algebraic and trigonometric form}
\LANGen{
\Question{
Given complex numbers \(z_{1} = 2\sqrt{2}\left (\cos \frac{\pi }{4} + \mathrm{i}\sin \frac{\pi }{4}\right )\) 
and \(z_{2} = \sqrt{2}\left (\cos \frac{7\pi }
{4} + \mathrm{i}\sin \frac{7\pi }
{4}\right )\), find
the product \(z_{1}z_{2}\).}
{\(4\)}{\(4\mathrm{i}\)}{\(- 4\mathrm{i}\)}{\(- 4\)}{}{}}
\LANGpl{
\Question{
Dane są liczby zespolone  \(z_{1} = 2\sqrt{2}\left (\cos \frac{\pi }{4} + \mathrm{i}\sin \frac{\pi }{4}\right )\) 
i  \(z_{2} = \sqrt{2}\left (\cos \frac{7\pi }
{4} + \mathrm{i}\sin \frac{7\pi }
{4}\right )\), oblicz iloczyn
 \(z_{1}z_{2}\).}
{\(4\)}{\(4\mathrm{i}\)}{\(- 4\mathrm{i}\)}{\(- 4\)}{}{}}
\LANGcs{
\Question{
Jsou dána komplexní čísla
\(z_{1} =\, \) \(2\sqrt{2}\left (\cos \frac{\pi }{4} + \mathrm{i}\sin \frac{\pi }{4}\right )\),
\(z_{2} =\, \) \(\sqrt{2}\left (\cos \frac{7\pi }
{4} + \mathrm{i}\sin \frac{7\pi }
{4}\right )\).
Určete součin \(z_{1}z_{2}\) v algebraickém tvaru.}
{\(4\)}{\(4\mathrm{i}\)}{\(- 4\mathrm{i}\)}{\(- 4\)}{}{}}
\LANGsk{
\Question{
Sú dané komplexné čísla \(z_{1} = 2\sqrt{2}\left (\cos \frac{\pi }{4} + \mathrm{i}\sin \frac{\pi }{4}\right )\) 
a \(z_{2} = \sqrt{2}\left (\cos \frac{7\pi }
{4} + \mathrm{i}\sin \frac{7\pi }
{4}\right )\). Určte
súčin \(z_{1}z_{2}\) v algebraickom tvaru.}
{\(4\)}{\(4\mathrm{i}\)}{\(- 4\mathrm{i}\)}{\(- 4\)}{}{}}
\LANGes{
\Question{
Dados los números complejos \(z_{1} = 2\sqrt{2}\left (\cos \frac{\pi }{4} + \mathrm{i}\sin \frac{\pi }{4}\right )\) 
y \(z_{2} = \sqrt{2}\left (\cos \frac{7\pi }
{4} + \mathrm{i}\sin \frac{7\pi }
{4}\right )\), halla el producto
\(z_{1}z_{2}\).}
{\(4\)}{\(4\mathrm{i}\)}{\(- 4\mathrm{i}\)}{\(- 4\)}{}{}}
\End

\Data\ID{9000031210}
\Updated{2020-07-18 09:07:37}
\Level{B}
\SubArea{Complex numbers in algebraic and trigonometric form}
\LANGen{
\Question{
Given complex numbers \(z_{1} = 2\sqrt{3}\left (\cos \frac{\pi }{6} + \mathrm{i}\sin \frac{\pi }{6}\right )\) 
and \(z_{2} = \sqrt{3}\left (\cos \frac{4\pi }
{3} + \mathrm{i}\sin \frac{4\pi }
{3}\right )\), find the
quotient \(\frac{z_{1}} 
{z_{2}} \).}
{\(-\sqrt{3} + \mathrm{i}\)}{\(\sqrt{3} -\mathrm{i}\)}{\(\sqrt{3} + \mathrm{i}\)}{\(-\sqrt{3} -\mathrm{i}\)}{}{}}
\LANGpl{
\Question{
Dane są liczby zespolone  \(z_{1} = 2\sqrt{3}\left (\cos \frac{\pi }{6} + \mathrm{i}\sin \frac{\pi }{6}\right )\) 
i  \(z_{2} = \sqrt{3}\left (\cos \frac{4\pi }
{3} + \mathrm{i}\sin \frac{4\pi }
{3}\right )\), określ iloraz 
 \(\frac{z_{1}} 
{z_{2}} \).}
{\(-\sqrt{3} + \mathrm{i}\)}{\(\sqrt{3} -\mathrm{i}\)}{\(\sqrt{3} + \mathrm{i}\)}{\(-\sqrt{3} -\mathrm{i}\)}{}{}}
\LANGcs{
\Question{
Jsou dána komplexní čísla
\(z_{1} =\, \) \(2\sqrt{3}\left (\cos \frac{\pi }{6} + \mathrm{i}\sin \frac{\pi }{6}\right )\),
\(z_{2} =\, \) \(\sqrt{3}\left (\cos \frac{4\pi }
{3} + \mathrm{i}\sin \frac{4\pi }
{3}\right )\).
Určete jejich podíl \(\frac{z_{1}} 
{z_{2}} \) 
v algebraickém tvaru.}
{\(-\sqrt{3} + \mathrm{i}\)}{\(\sqrt{3} -\mathrm{i}\)}{\(\sqrt{3} + \mathrm{i}\)}{\(-\sqrt{3} -\mathrm{i}\)}{}{}}
\LANGsk{
\Question{
Sú dané komplexné čísla \(z_{1} = 2\sqrt{3}\left (\cos \frac{\pi }{6} + \mathrm{i}\sin \frac{\pi }{6}\right )\) 
a \(z_{2} = \sqrt{3}\left (\cos \frac{4\pi }
{3} + \mathrm{i}\sin \frac{4\pi }
{3}\right )\). Určte ich podiel \(\frac{z_{1}} 
{z_{2}} \) v algebraickom tvare.}
{\(-\sqrt{3} + \mathrm{i}\)}{\(\sqrt{3} -\mathrm{i}\)}{\(\sqrt{3} + \mathrm{i}\)}{\(-\sqrt{3} -\mathrm{i}\)}{}{}}
\LANGes{
\Question{
Dados los números complejos \(z_{1} = 2\sqrt{3}\left (\cos \frac{\pi }{6} + \mathrm{i}\sin \frac{\pi }{6}\right )\) 
y \(z_{2} = \sqrt{3}\left (\cos \frac{4\pi }
{3} + \mathrm{i}\sin \frac{4\pi }
{3}\right )\), halla el cociente
 \(\frac{z_{1}} 
{z_{2}} \).}
{\(-\sqrt{3} + \mathrm{i}\)}{\(\sqrt{3} -\mathrm{i}\)}{\(\sqrt{3} + \mathrm{i}\)}{\(-\sqrt{3} -\mathrm{i}\)}{}{}}
\End

\Data\ID{9000034807}
\Updated{2020-07-18 10:07:40}
\Level{B}
\SubArea{Complex numbers in algebraic and trigonometric form}
\LANGen{
\Question{
Find the polar form of the complex number
\(z = 2\mathrm{i}\).}
{\(2\left (\cos \frac{\pi }{2} + \mathrm{i}\sin \frac{\pi }{2}\right )\)}{\(\sqrt{2}\left (\cos \frac{\pi }{2} + \mathrm{i}\sin \frac{\pi }{2}\right )\)}{\(\cos \frac{\pi }{2} + \mathrm{i}\sin \frac{\pi }{2}\)}{\(2\left (\cos 0 + \mathrm{i}\sin 0\right )\)}{}{}}
\LANGpl{
\Question{
Wskaż postać biegunową  liczby zespolonej
\(z = 2\mathrm{i}\).}
{\(2\left (\cos \frac{\pi }{2} + \mathrm{i}\sin \frac{\pi }{2}\right )\)}{\(\sqrt{2}\left (\cos \frac{\pi }{2} + \mathrm{i}\sin \frac{\pi }{2}\right )\)}{\(\cos \frac{\pi }{2} + \mathrm{i}\sin \frac{\pi }{2}\)}{\(2\left (\cos 0 + \mathrm{i}\sin 0\right )\)}{}{}}
\LANGcs{
\Question{
Vyjádřete komplexní číslo \(z = 2\mathrm{i}\) 
v goniometrickém tvaru.}
{\(2\left (\cos \frac{\pi }{2} + \mathrm{i}\sin \frac{\pi }{2}\right )\)}{\(\sqrt{2}\left (\cos \frac{\pi }{2} + \mathrm{i}\sin \frac{\pi }{2}\right )\)}{\(\cos \frac{\pi }{2} + \mathrm{i}\sin \frac{\pi }{2}\)}{\(2\left (\cos 0 + \mathrm{i}\sin 0\right )\)}{}{}}
\LANGsk{
\Question{
Vyjadrite komplexné číslo
\(z = 2\mathrm{i}\) v goniometrickom tvare.}
{\(2\left (\cos \frac{\pi }{2} + \mathrm{i}\sin \frac{\pi }{2}\right )\)}{\(\sqrt{2}\left (\cos \frac{\pi }{2} + \mathrm{i}\sin \frac{\pi }{2}\right )\)}{\(\cos \frac{\pi }{2} + \mathrm{i}\sin \frac{\pi }{2}\)}{\(2\left (\cos 0 + \mathrm{i}\sin 0\right )\)}{}{}}
\LANGes{
\Question{
Determina la forma polar del número complejo
\(z = 2\mathrm{i}\).}
{\(2\left (\cos \frac{\pi }{2} + \mathrm{i}\sin \frac{\pi }{2}\right )\)}{\(\sqrt{2}\left (\cos \frac{\pi }{2} + \mathrm{i}\sin \frac{\pi }{2}\right )\)}{\(\cos \frac{\pi }{2} + \mathrm{i}\sin \frac{\pi }{2}\)}{\(2\left (\cos 0 + \mathrm{i}\sin 0\right )\)}{}{}}
\End

\Data\ID{9000034808}
\Updated{2020-07-18 11:07:43}
\Level{B}
\SubArea{Complex numbers in algebraic and trigonometric form}
\LANGen{
\Question{
Find the algebraic form of the complex number
\(z = 2\left (\cos \pi + \mathrm{i}\sin \pi \right )\).}
{\(- 2\)}{\(2\)}{\(- 2\mathrm{i}\)}{\(2\mathrm{i}\)}{}{}}
\LANGpl{
\Question{
Wskaż postać algebraiczną liczby zespolonej
\(z = 2\left (\cos \pi + \mathrm{i}\sin \pi \right )\).}
{\(- 2\)}{\(2\)}{\(- 2\mathrm{i}\)}{\(2\mathrm{i}\)}{}{}}
\LANGcs{
\Question{
Vyjádřete komplexní číslo \(z = 2\left (\cos \pi + \mathrm{i}\sin \pi \right )\) 
v algebraickém tvaru.}
{\(- 2\)}{\(2\)}{\(- 2\mathrm{i}\)}{\(2\mathrm{i}\)}{}{}}
\LANGsk{
\Question{
Vyjadrite komplexné číslo 
\(z = 2\left (\cos \pi + \mathrm{i}\sin \pi \right )\) v algebraickom tvare.}
{\(- 2\)}{\(2\)}{\(- 2\mathrm{i}\)}{\(2\mathrm{i}\)}{}{}}
\LANGes{
\Question{
Determina la fórmula algebraica del número complejo
\(z = 2\left (\cos \pi + \mathrm{i}\sin \pi \right )\).}
{\(- 2\)}{\(2\)}{\(- 2\mathrm{i}\)}{\(2\mathrm{i}\)}{}{}}
\End

\Data\ID{9000034809}
\Updated{2021-04-24 05:04:35}
\Level{B}
\SubArea{Complex numbers in algebraic and trigonometric form}
\LANGen{
\Question{
Given complex numbers \(z_{1} = 2\left (\cos \frac{\pi }{6} + \mathrm{i}\sin \frac{\pi }{6}\right )\) 
and \(z_{2} = \sqrt{3}\left (\cos \frac{4\pi }
{3} + \mathrm{i}\sin \frac{4\pi }
{3}\right )\), find the angle in the
polar form of the product \(z_{1}z_{2}\).}
{\(\frac{3\pi }
{2}\)}{\(\frac{2}
{9}\pi \)}{\(\frac{5}
{9}\pi \)}{\(3\pi \)}{}{}}
\LANGpl{
\Question{
Dane są liczby zespolone  \(z_{1} = 2\left (\cos \frac{\pi }{6} + \mathrm{i}\sin \frac{\pi }{6}\right )\) 
i  \(z_{2} = \sqrt{3}\left (\cos \frac{4\pi }
{3} + \mathrm{i}\sin \frac{4\pi }
{3}\right )\), wyznacz kąt w postaci biegunowej iloczynu liczb zespolonych
\(z_{1}z_{2}\).}
{\(\frac{3\pi }
{2}\)}{\(\frac{2}
{9}\pi \)}{\(\frac{5}
{9}\pi \)}{\(3\pi \)}{}{}}
\LANGcs{
\Question{
Jsou dána komplexní čísla
\(z_{1} =\) \(2\left (\cos \frac{\pi }{6} + \mathrm{i}\sin \frac{\pi }{6}\right )\),
\(z_{2} =\) \(\sqrt{3}\left (\cos \frac{4\pi }
{3} + \mathrm{i}\sin \frac{4\pi }
{3}\right )\).
Určete hlavní hodnotu argumentu jejich součinu.}
{\(\frac{3\pi }
{2}\)}{\(\frac{2}
{9}\pi \)}{\(\frac{5}
{9}\pi \)}{\(3\pi \)}{}{}}
\LANGsk{
\Question{
Sú dané komplexné čísla \(z_{1} = 2\left (\cos \frac{\pi }{6} + \mathrm{i}\sin \frac{\pi }{6}\right )\) 
a \(z_{2} = \sqrt{3}\left (\cos \frac{4\pi }
{3} + \mathrm{i}\sin \frac{4\pi }
{3}\right )\). Určte základnú hodnotu argumentu ich súčinu.}
{\(\frac{3\pi }
{2}\)}{\(\frac{2}
{9}\pi \)}{\(\frac{5}
{9}\pi \)}{\(3\pi \)}{}{}}
\LANGes{
\Question{
Dados los números complejos \(z_{1} = 2\left (\cos \frac{\pi }{6} + \mathrm{i}\sin \frac{\pi }{6}\right )\) 
y \(z_{2} = \sqrt{3}\left (\cos \frac{4\pi }
{3} + \mathrm{i}\sin \frac{4\pi }
{3}\right )\), determina el ángulo del producto \(z_{1}z_{2}\) en forma polar.}
{\(\frac{3\pi }
{2}\)}{\(\frac{2}
{9}\pi \)}{\(\frac{5}
{9}\pi \)}{\(3\pi \)}{}{}}
\End

\Data\ID{9000034810}
\Updated{2021-04-24 05:04:20}
\Level{B}
\SubArea{Complex numbers in algebraic and trigonometric form}
\LANGen{
\Question{
Given complex numbers \(z_{1} = 2\left (\cos \frac{\pi }{4} + \mathrm{i}\sin \frac{\pi }{4}\right )\) 
and \(z_{2} = \sqrt{2}\left (\cos \frac{7\pi }
{4} + \mathrm{i}\sin \frac{7\pi }
{4}\right )\),
find the angle in the polar form of the quotient
\(\frac{z_{1}} 
{z_{2}} \).}
{\(\frac{\pi }
{2}\)}{\(- \frac{\pi } {2}\)}{\(-\frac{3} 
{2}\pi \)}{\(\frac{3}
{2}\pi \)}{}{}}
\LANGpl{
\Question{
Dane są liczby zespolone  \(z_{1} = 2\left (\cos \frac{\pi }{4} + \mathrm{i}\sin \frac{\pi }{4}\right )\) 
and \(z_{2} = \sqrt{2}\left (\cos \frac{7\pi }
{4} + \mathrm{i}\sin \frac{7\pi }
{4}\right )\),
wyznacz kąt w postaci biegunowej ilorazu liczb zespolonych
\(\frac{z_{1}} 
{z_{2}} \).}
{\(\frac{\pi }
{2}\)}{\(- \frac{\pi } {2}\)}{\(-\frac{3} 
{2}\pi \)}{\(\frac{3}
{2}\pi \)}{}{}}
\LANGcs{
\Question{
Jsou dána komplexní čísla
\(z_{1} =\) \(2\left (\cos \frac{\pi }{4} + \mathrm{i}\sin \frac{\pi }{4}\right )\),
\(z_{2} =\) \(\sqrt{2}\left (\cos \frac{7\pi }
{4} + \mathrm{i}\sin \frac{7\pi }
{4}\right )\).
Určete hlavní hodnotu argumentu jejich podílu
\(\frac{z_{1}} 
{z_{2}} \).}
{\(\frac{\pi }
{2}\)}{\(- \frac{\pi } {2}\)}{\(-\frac{3} 
{2}\pi \)}{\(\frac{3}
{2}\pi \)}{}{}}
\LANGsk{
\Question{
Sú dané komplexné čísla \(z_{1} = 2\left (\cos \frac{\pi }{4} + \mathrm{i}\sin \frac{\pi }{4}\right )\) 
a \(z_{2} = \sqrt{2}\left (\cos \frac{7\pi }
{4} + \mathrm{i}\sin \frac{7\pi }
{4}\right )\).
Určte základnú hodnotu argumentu ich podielu
\(\frac{z_{1}} 
{z_{2}} \).}
{\(\frac{\pi }
{2}\)}{\(- \frac{\pi } {2}\)}{\(-\frac{3} 
{2}\pi \)}{\(\frac{3}
{2}\pi \)}{}{}}
\LANGes{
\Question{
Dados los números complejos \(z_{1} = 2\left (\cos \frac{\pi }{4} + \mathrm{i}\sin \frac{\pi }{4}\right )\) 
y \(z_{2} = \sqrt{2}\left (\cos \frac{7\pi }
{4} + \mathrm{i}\sin \frac{7\pi }
{4}\right )\),
determina el ángulo en la forma polar del cociente
\(\frac{z_{1}} 
{z_{2}} \).}
{\(\frac{\pi }
{2}\)}{\(- \frac{\pi } {2}\)}{\(-\frac{3} 
{2}\pi \)}{\(\frac{3}
{2}\pi \)}{}{}}
\End

\Data\ID{9000035704}
\Updated{2020-07-18 11:07:55}
\Level{B}
\SubArea{Complex numbers in algebraic and trigonometric form}
\IMG{000357_04-figure0.pdf}
\LANGen{
\Question{
Find the polar form of the complex number \( A \) graphed in the complex plane as shown in the picture.}
{\(z = 2\sqrt{2}\left (\cos \frac{3\pi }
{4} + \mathrm{i}\sin \frac{3\pi }
{4}\right )\)}{\(z = 2\sqrt{2}\left (\cos \frac{\pi }{4} -\mathrm{i}\sin \frac{\pi }{4}\right )\)}{\(z = 2\sqrt{2}\left (-\cos \frac{\pi }{4} + \mathrm{i}\sin \frac{\pi }{4}\right )\)}{\(z = 2\sqrt{2}\left (\cos \frac{5\pi }
{4} + \mathrm{i}\sin \frac{5\pi }
{4}\right )\)}{}{}}
\LANGpl{
\Question{
Wskaż postać biegunową  liczby zespolonej zaznaczonej  na płaszczyźnie zespolonej.}
{\(z = 2\sqrt{2}\left (\cos \frac{3\pi }
{4} + \mathrm{i}\sin \frac{3\pi }
{4}\right )\)}{\(z = 2\sqrt{2}\left (\cos \frac{\pi }{4} -\mathrm{i}\sin \frac{\pi }{4}\right )\)}{\(z = 2\sqrt{2}\left (-\cos \frac{\pi }{4} + \mathrm{i}\sin \frac{\pi }{4}\right )\)}{\(z = 2\sqrt{2}\left (\cos \frac{5\pi }
{4} + \mathrm{i}\sin \frac{5\pi }
{4}\right )\)}{}{}}
\LANGcs{
\Question{
Nalezněte goniometrický tvar komplexního čísla \( A \) zobrazeného na obrázku:}
{\(z = 2\sqrt{2}\left (\cos \frac{3\pi }
{4} + \mathrm{i}\sin \frac{3\pi }
{4}\right )\)}{\(z = 2\sqrt{2}\left (\cos \frac{\pi }{4} -\mathrm{i}\sin \frac{\pi }{4}\right )\)}{\(z = 2\sqrt{2}\left (-\cos \frac{\pi }{4} + \mathrm{i}\sin \frac{\pi }{4}\right )\)}{\(z = 2\sqrt{2}\left (\cos \frac{5\pi }
{4} + \mathrm{i}\sin \frac{5\pi }
{4}\right )\)}{}{}}
\LANGsk{
\Question{
Bod na \( A \) obrázku je obrazom komplexného čísla:}
{\(z = 2\sqrt{2}\left (\cos \frac{3\pi }
{4} + \mathrm{i}\sin \frac{3\pi }
{4}\right )\)}{\(z = 2\sqrt{2}\left (\cos \frac{\pi }{4} -\mathrm{i}\sin \frac{\pi }{4}\right )\)}{\(z = 2\sqrt{2}\left (-\cos \frac{\pi }{4} + \mathrm{i}\sin \frac{\pi }{4}\right )\)}{\(z = 2\sqrt{2}\left (\cos \frac{5\pi }
{4} + \mathrm{i}\sin \frac{5\pi }
{4}\right )\)}{}{}}
\LANGes{
\Question{
¿Cuál es la forma polar de un número complejo representado en el plano complejo por el punto \( A \)  (mira la imagen)?}
{\(z = 2\sqrt{2}\left (\cos \frac{3\pi }
{4} + \mathrm{i}\sin \frac{3\pi }
{4}\right )\)}{\(z = 2\sqrt{2}\left (\cos \frac{\pi }{4} -\mathrm{i}\sin \frac{\pi }{4}\right )\)}{\(z = 2\sqrt{2}\left (-\cos \frac{\pi }{4} + \mathrm{i}\sin \frac{\pi }{4}\right )\)}{\(z = 2\sqrt{2}\left (\cos \frac{5\pi }
{4} + \mathrm{i}\sin \frac{5\pi }
{4}\right )\)}{}{}}
\End

\Data\ID{9000035805}
\Updated{2020-07-18 11:07:25}
\Level{B}
\SubArea{Complex numbers in algebraic and trigonometric form}
\LANGen{
\Question{
Given the complex numbers 
\[ 
 \text{$a = 2\left (\cos \frac{2\pi }
{3} + \mathrm{i}\sin \frac{2\pi }
{3}\right )$, $b = \sqrt{2}\left (\cos \frac{3\pi }
{4} + \mathrm{i}\sin \frac{3\pi }
{4}\right )$,}
\]
find the product \(ab\).}
{\(2\sqrt{2}\left (\cos \frac{17\pi }
{12} + \mathrm{i}\sin \frac{17\pi }
{12}\right )\)}{\(2\sqrt{2}\left (\cos \frac{\pi }{2} + \mathrm{i}\sin \frac{\pi }{2}\right )\)}{\(2\sqrt{2}\left (\cos \frac{5\pi }
{7} + \mathrm{i}\sin \frac{5\pi }
{7}\right )\)}{\(2\sqrt{2}\left (\cos \frac{5\pi }
{12} + \mathrm{i}\sin \frac{5\pi }
{12}\right )\)}{}{}}
\LANGpl{
\Question{
Dane są liczby zespolone
\[ 
 \text{$a = 2\left (\cos \frac{2\pi }
{3} + \mathrm{i}\sin \frac{2\pi }
{3}\right )$, $b = \sqrt{2}\left (\cos \frac{3\pi }
{4} + \mathrm{i}\sin \frac{3\pi }
{4}\right )$,}
\]
oblicz iloczyn \(ab\).}
{\(2\sqrt{2}\left (\cos \frac{17\pi }
{12} + \mathrm{i}\sin \frac{17\pi }
{12}\right )\)}{\(2\sqrt{2}\left (\cos \frac{\pi }{2} + \mathrm{i}\sin \frac{\pi }{2}\right )\)}{\(2\sqrt{2}\left (\cos \frac{5\pi }
{7} + \mathrm{i}\sin \frac{5\pi }
{7}\right )\)}{\(2\sqrt{2}\left (\cos \frac{5\pi }
{12} + \mathrm{i}\sin \frac{5\pi }
{12}\right )\)}{}{}}
\LANGcs{
\Question{
Jsou dána komplexní čísla \(a = 2\left (\cos \frac{2\pi }
{3} + \mathrm{i}\sin \frac{2\pi }
{3}\right )\),
\(b = \sqrt{2}\left (\cos \frac{3\pi }
{4} + \mathrm{i}\sin \frac{3\pi }
{4}\right )\). Součin \(a\cdot b\) 
se rovná:}
{\(2\sqrt{2}\left (\cos \frac{17\pi }
{12} + \mathrm{i}\sin \frac{17\pi }
{12}\right )\)}{\(2\sqrt{2}\left (\cos \frac{\pi }{2} + \mathrm{i}\sin \frac{\pi }{2}\right )\)}{\(2\sqrt{2}\left (\cos \frac{5\pi }
{7} + \mathrm{i}\sin \frac{5\pi }
{7}\right )\)}{\(2\sqrt{2}\left (\cos \frac{5\pi }
{12} + \mathrm{i}\sin \frac{5\pi }
{12}\right )\)}{}{}}
\LANGsk{
\Question{
Sú dané komplexné čísla 
\[ 
 \text{$a = 2\left (\cos \frac{2\pi }
{3} + \mathrm{i}\sin \frac{2\pi }
{3}\right )$, $b = \sqrt{2}\left (\cos \frac{3\pi }
{4} + \mathrm{i}\sin \frac{3\pi }
{4}\right )$.}
\] Súčin \(ab\) sa rovná:}
{\(2\sqrt{2}\left (\cos \frac{17\pi }
{12} + \mathrm{i}\sin \frac{17\pi }
{12}\right )\)}{\(2\sqrt{2}\left (\cos \frac{\pi }{2} + \mathrm{i}\sin \frac{\pi }{2}\right )\)}{\(2\sqrt{2}\left (\cos \frac{5\pi }
{7} + \mathrm{i}\sin \frac{5\pi }
{7}\right )\)}{\(2\sqrt{2}\left (\cos \frac{5\pi }
{12} + \mathrm{i}\sin \frac{5\pi }
{12}\right )\)}{}{}}
\LANGes{
\Question{
Dados los números complejos
\[ 
 \text{$a = 2\left (\cos \frac{2\pi }
{3} + \mathrm{i}\sin \frac{2\pi }
{3}\right )$, $b = \sqrt{2}\left (\cos \frac{3\pi }
{4} + \mathrm{i}\sin \frac{3\pi }
{4}\right )$,}
\]
determina el producto \(ab\).}
{\(2\sqrt{2}\left (\cos \frac{17\pi }
{12} + \mathrm{i}\sin \frac{17\pi }
{12}\right )\)}{\(2\sqrt{2}\left (\cos \frac{\pi }{2} + \mathrm{i}\sin \frac{\pi }{2}\right )\)}{\(2\sqrt{2}\left (\cos \frac{5\pi }
{7} + \mathrm{i}\sin \frac{5\pi }
{7}\right )\)}{\(2\sqrt{2}\left (\cos \frac{5\pi }
{12} + \mathrm{i}\sin \frac{5\pi }
{12}\right )\)}{}{}}
\End

\Data\ID{9000035806}
\Updated{2020-07-18 11:07:07}
\Level{B}
\SubArea{Complex numbers in algebraic and trigonometric form}
\LANGen{
\Question{
Given the complex numbers 
\[ 
 \text{ $a = 2\left (\cos \frac{5\pi }
{3} + \mathrm{i}\sin \frac{5\pi }
{3}\right )$, $b = 3\left (\cos \frac{11\pi }
 {6} + \mathrm{i}\sin \frac{11\pi }
 {6} \right )$,}
\]
find the quotient \(\frac{a}
{b}\).}
{\(\frac{2}
{3}\left (\cos \frac{11\pi }
 {6} + \mathrm{i}\sin \frac{11\pi }
 {6} \right )\)}{\(\frac{2}
{3}\left (\cos \frac{\pi }
{6} + \mathrm{i}\sin \frac{\pi }
{6}\right )\)}{\(\frac{2}
{3}\left (\cos \frac{5\pi }
{6} + \mathrm{i}\sin \frac{5\pi }
{6}\right )\)}{\(\frac{2}
{3}\left (\cos \frac{7\pi }
{6} + \mathrm{i}\sin \frac{7\pi }
{6}\right )\)}{}{}}
\LANGpl{
\Question{
Dane są liczby zespolone
\[ 
 \text{ $a = 2\left (\cos \frac{5\pi }
{3} + \mathrm{i}\sin \frac{5\pi }
{3}\right )$, $b = 3\left (\cos \frac{11\pi }
 {6} + \mathrm{i}\sin \frac{11\pi }
 {6} \right )$,}
\]
oblicz iloraz \(\frac{a}
{b}\).}
{\(\frac{2}
{3}\left (\cos \frac{11\pi }
 {6} + \mathrm{i}\sin \frac{11\pi }
 {6} \right )\)}{\(\frac{2}
{3}\left (\cos \frac{\pi }
{6} + \mathrm{i}\sin \frac{\pi }
{6}\right )\)}{\(\frac{2}
{3}\left (\cos \frac{5\pi }
{6} + \mathrm{i}\sin \frac{5\pi }
{6}\right )\)}{\(\frac{2}
{3}\left (\cos \frac{7\pi }
{6} + \mathrm{i}\sin \frac{7\pi }
{6}\right )\)}{}{}}
\LANGcs{
\Question{
Jsou dána komplexní čísla \(a = 2\left (\cos \frac{5\pi }
{3} + \mathrm{i}\sin \frac{5\pi }
{3}\right )\),
\(b = 3\left (\cos \frac{11\pi }
 {6} + \mathrm{i}\sin \frac{11\pi }
 {6} \right )\). Podíl \(\frac{a}
{b}\) 
se rovná:}
{\(\frac{2}
{3}\left (\cos \frac{11\pi }
 {6} + \mathrm{i}\sin \frac{11\pi }
 {6} \right )\)}{\(\frac{2}
{3}\left (\cos \frac{\pi }
{6} + \mathrm{i}\sin \frac{\pi }
{6}\right )\)}{\(\frac{2}
{3}\left (\cos \frac{5\pi }
{6} + \mathrm{i}\sin \frac{5\pi }
{6}\right )\)}{\(\frac{2}
{3}\left (\cos \frac{7\pi }
{6} + \mathrm{i}\sin \frac{7\pi }
{6}\right )\)}{}{}}
\LANGsk{
\Question{
Sú dané komplexné čísla 
\[ 
 \text{ $a = 2\left (\cos \frac{5\pi }
{3} + \mathrm{i}\sin \frac{5\pi }
{3}\right )$, $b = 3\left (\cos \frac{11\pi }
 {6} + \mathrm{i}\sin \frac{11\pi }
 {6} \right )$.}
\]
Podiel \(\frac{a}
{b}\) sa rovná:}
{\(\frac{2}
{3}\left (\cos \frac{11\pi }
 {6} + \mathrm{i}\sin \frac{11\pi }
 {6} \right )\)}{\(\frac{2}
{3}\left (\cos \frac{\pi }
{6} + \mathrm{i}\sin \frac{\pi }
{6}\right )\)}{\(\frac{2}
{3}\left (\cos \frac{5\pi }
{6} + \mathrm{i}\sin \frac{5\pi }
{6}\right )\)}{\(\frac{2}
{3}\left (\cos \frac{7\pi }
{6} + \mathrm{i}\sin \frac{7\pi }
{6}\right )\)}{}{}}
\LANGes{
\Question{
Dados los números complejos
\[ 
 \text{ $a = 2\left (\cos \frac{5\pi }
{3} + \mathrm{i}\sin \frac{5\pi }
{3}\right )$, $b = 3\left (\cos \frac{11\pi }
 {6} + \mathrm{i}\sin \frac{11\pi }
 {6} \right )$,}
\]
determina el cociente \(\frac{a}
{b}\).}
{\(\frac{2}
{3}\left (\cos \frac{11\pi }
 {6} + \mathrm{i}\sin \frac{11\pi }
 {6} \right )\)}{\(\frac{2}
{3}\left (\cos \frac{\pi }
{6} + \mathrm{i}\sin \frac{\pi }
{6}\right )\)}{\(\frac{2}
{3}\left (\cos \frac{5\pi }
{6} + \mathrm{i}\sin \frac{5\pi }
{6}\right )\)}{\(\frac{2}
{3}\left (\cos \frac{7\pi }
{6} + \mathrm{i}\sin \frac{7\pi }
{6}\right )\)}{}{}}
\End

\Data\ID{9000037408}
\Updated{2020-07-18 11:07:08}
\Level{B}
\SubArea{Complex numbers in algebraic and trigonometric form}
\LANGen{
\Question{
Find the polar form of the following complex number
\[z=\frac{1}
{\cos \frac{2\pi } 
{3} +\mathrm{i}\sin \frac{2\pi } 
{3} }. \]}
{\(\cos \frac{4\pi }
{3} + \mathrm{i}\sin \frac{4\pi }
{3}\)}{\(\cos \left (-\frac{4\pi }
{3}\right ) + \mathrm{i}\sin \left (-\frac{4\pi }
{3}\right )\)}{\(\cos \frac{3}
{2\pi } + \mathrm{i}\sin \frac{3}
{2\pi }\)}{\(\cos \frac{3}
{2\pi } -\mathrm{i}\sin \frac{3}
{2\pi }\)}{}{}}
\LANGpl{
\Question{
Wskaż postać biegunową  liczby zespolonej
\[z=\frac{1}
{\cos \frac{2\pi } 
{3} +\mathrm{i}\sin \frac{2\pi } 
{3} }. \]}
{\(\cos \frac{4\pi }
{3} + \mathrm{i}\sin \frac{4\pi }
{3}\)}{\(\cos \left (-\frac{4\pi }
{3}\right ) + \mathrm{i}\sin \left (-\frac{4\pi }
{3}\right )\)}{\(\cos \frac{3}
{2\pi } + \mathrm{i}\sin \frac{3}
{2\pi }\)}{\(\cos \frac{3}
{2\pi } -\mathrm{i}\sin \frac{3}
{2\pi }\)}{}{}}
\LANGcs{
\Question{
Vyjádřete v goniometrickém tvaru dané komplexní číslo
\[z=\frac{1}
{\cos \frac{2\pi } 
{3} +\mathrm{i}\sin \frac{2\pi } 
{3} }. \]}
{\(\cos \frac{4\pi }
{3} + \mathrm{i}\sin \frac{4\pi }
{3}\)}{\(\cos \left (-\frac{4\pi }
{3}\right ) + \mathrm{i}\sin \left (-\frac{4\pi }
{3}\right )\)}{\(\cos \frac{3}
{2\pi } + \mathrm{i}\sin \frac{3}
{2\pi }\)}{\(\cos \frac{3}
{2\pi } -\mathrm{i}\sin \frac{3}
{2\pi }\)}{}{}}
\LANGsk{
\Question{
Vyjadrite v goniometrickom tvare komplexné číslo 
\[z=\frac{1}
{\cos \frac{2\pi } 
{3} +\mathrm{i}\sin \frac{2\pi } 
{3} }. \]}
{\(\cos \frac{4\pi }
{3} + \mathrm{i}\sin \frac{4\pi }
{3}\)}{\(\cos \left (-\frac{4\pi }
{3}\right ) + \mathrm{i}\sin \left (-\frac{4\pi }
{3}\right )\)}{\(\cos \frac{3}
{2\pi } + \mathrm{i}\sin \frac{3}
{2\pi }\)}{\(\cos \frac{3}
{2\pi } -\mathrm{i}\sin \frac{3}
{2\pi }\)}{}{}}
\LANGes{
\Question{
Determina la forma polar del siguiente número complejo
\[z=\frac{1}
{\cos \frac{2\pi } 
{3} +\mathrm{i}\sin \frac{2\pi } 
{3} }. \]}
{\(\cos \frac{4\pi }
{3} + \mathrm{i}\sin \frac{4\pi }
{3}\)}{\(\cos \left (-\frac{4\pi }
{3}\right ) + \mathrm{i}\sin \left (-\frac{4\pi }
{3}\right )\)}{\(\cos \frac{3}
{2\pi } + \mathrm{i}\sin \frac{3}
{2\pi }\)}{\(\cos \frac{3}
{2\pi } -\mathrm{i}\sin \frac{3}
{2\pi }\)}{}{}}
\End

\Data\ID{9000037409}
\Updated{2020-07-18 11:07:48}
\Level{B}
\SubArea{Complex numbers in algebraic and trigonometric form}
\LANGen{
\Question{
Find the polar form of the complex number
\[z=\frac{1}
{\cos \frac{7\pi } 
{6} +\mathrm{i}\sin \frac{7\pi } 
{6} }. \]}
{\(\cos \frac{5\pi }
{6} + \mathrm{i}\sin \frac{5\pi }
{6}\)}{\(\cos \left (-\frac{5\pi }
{6}\right ) + \mathrm{i}\sin \left (-\frac{5\pi }
{6}\right )\)}{\(\cos \frac{\pi }{6} + \mathrm{i}\sin \frac{\pi }{6}\)}{\(\cos \left (-\frac{\pi }{6}\right ) + \mathrm{i}\sin \left (-\frac{\pi }{6}\right )\)}{}{}}
\LANGpl{
\Question{
Wskaż postać biegunową  liczby zespolonej
\[z=\frac{1}
{\cos \frac{7\pi } 
{6} +\mathrm{i}\sin \frac{7\pi } 
{6} }. \]}
{\(\cos \frac{5\pi }
{6} + \mathrm{i}\sin \frac{5\pi }
{6}\)}{\(\cos \left (-\frac{5\pi }
{6}\right ) + \mathrm{i}\sin \left (-\frac{5\pi }
{6}\right )\)}{\(\cos \frac{\pi }{6} + \mathrm{i}\sin \frac{\pi }{6}\)}{\(\cos \left (-\frac{\pi }{6}\right ) + \mathrm{i}\sin \left (-\frac{\pi }{6}\right )\)}{}{}}
\LANGcs{
\Question{
Vyjádřete v goniometrickém tvaru dané komplexní číslo
\[z=\frac{1}
{\cos \frac{7\pi } 
{6} +\mathrm{i}\sin \frac{7\pi } 
{6} }. \]}
{\(\cos \frac{5\pi }
{6} + \mathrm{i}\sin \frac{5\pi }
{6}\)}{\(\cos \left (-\frac{5\pi }
{6}\right ) + \mathrm{i}\sin \left (-\frac{5\pi }
{6}\right )\)}{\(\cos \frac{\pi }{6} + \mathrm{i}\sin \frac{\pi }{6}\)}{\(\cos \left (-\frac{\pi }{6}\right ) + \mathrm{i}\sin \left (-\frac{\pi }{6}\right )\)}{}{}}
\LANGsk{
\Question{
Vyjadrite v goniometrickom tvare dané komplexné číslo 
\[\frac{1}
{\cos \frac{7\pi } 
{6} +\mathrm{i}\sin \frac{7\pi } 
{6} }.\]}
{\(\cos \frac{5\pi }
{6} + \mathrm{i}\sin \frac{5\pi }
{6}\)}{\(\cos \left (-\frac{5\pi }
{6}\right ) + \mathrm{i}\sin \left (-\frac{5\pi }
{6}\right )\)}{\(\cos \frac{\pi }{6} + \mathrm{i}\sin \frac{\pi }{6}\)}{\(\cos \left (-\frac{\pi }{6}\right ) + \mathrm{i}\sin \left (-\frac{\pi }{6}\right )\)}{}{}}
\LANGes{
\Question{
Determina la forma polar del número complejo
\[z=\frac{1}
{\cos \frac{7\pi } 
{6} +\mathrm{i}\sin \frac{7\pi } 
{6} }. \]}
{\(\cos \frac{5\pi }
{6} + \mathrm{i}\sin \frac{5\pi }
{6}\)}{\(\cos \left (-\frac{5\pi }
{6}\right ) + \mathrm{i}\sin \left (-\frac{5\pi }
{6}\right )\)}{\(\cos \frac{\pi }{6} + \mathrm{i}\sin \frac{\pi }{6}\)}{\(\cos \left (-\frac{\pi }{6}\right ) + \mathrm{i}\sin \left (-\frac{\pi }{6}\right )\)}{}{}}
\End

\Data\ID{9000037508}
\Updated{2020-07-18 11:07:15}
\Level{B}
\SubArea{Complex numbers in algebraic and trigonometric form}
\LANGen{
\Question{
Find the absolute value of the following complex number. 
\[ 
 \sqrt{2}\left (\cos \frac{\pi }
{3} + \mathrm{i}\sin \frac{\pi }
{3}\right )
\]}
{\(\sqrt{2}\)}{\(\sqrt{2} + 2\)}{\(2\)}{\(\sqrt{2} - 2\)}{}{}}
\LANGpl{
\Question{
Określ wartość bezwzględną podanej liczby zespolonej.
\[ 
 \sqrt{2}\left (\cos \frac{\pi }
{3} + \mathrm{i}\sin \frac{\pi }
{3}\right )
\]}
{\(\sqrt{2}\)}{\(\sqrt{2} + 2\)}{\(2\)}{\(\sqrt{2} - 2\)}{}{}}
\LANGcs{
\Question{
Určete absolutní hodnotu komplexního čísla
\(a\).
\[ 
 a = \sqrt{2}\left (\cos \frac{\pi }
{3} + \mathrm{i}\sin \frac{\pi }
{3}\right )
\]}
{\(\sqrt{2}\)}{\(\sqrt{2} + 2\)}{\(2\)}{\(\sqrt{2} - 2\)}{}{}}
\LANGsk{
\Question{
Určte absolútnu hodnotu daného komplexného čísla. 
\[ 
 \sqrt{2}\left (\cos \frac{\pi }
{3} + \mathrm{i}\sin \frac{\pi }
{3}\right )
\]}
{\(\sqrt{2}\)}{\(\sqrt{2} + 2\)}{\(2\)}{\(\sqrt{2} - 2\)}{}{}}
\LANGes{
\Question{
Determina el valor absoluto del siguiente número complejo. 
\[ 
 \sqrt{2}\left (\cos \frac{\pi }
{3} + \mathrm{i}\sin \frac{\pi }
{3}\right )
\]}
{\(\sqrt{2}\)}{\(\sqrt{2} + 2\)}{\(2\)}{\(\sqrt{2} - 2\)}{}{}}
\End

\Data\ID{9000037509}
\Updated{2020-07-18 11:07:50}
\Level{B}
\SubArea{Complex numbers in algebraic and trigonometric form}
\LANGen{
\Question{
Given complex numbers 
\[ 
 a = 3\left (\cos \frac{\pi }
{3} + \mathrm{i}\sin \frac{\pi }
{3}\right ),\quad b = \sqrt{2}\left (\cos \frac{2\pi }
{3} + \mathrm{i}\sin \frac{2\pi }
{3}\right )
\]
find the product \(ab\).}
{\(- 3\sqrt{2}\)}{\(3\sqrt{2}\left (\cos \frac{\pi }{2} + \mathrm{i}\sin \frac{\pi }{2}\right )\)}{\(3\sqrt{2}\left (\cos \frac{\pi }{2} -\mathrm{i}\sin \frac{\pi }{2}\right )\)}{\(- 3\sqrt{2}\left (\cos \frac{\pi }{2} + \mathrm{i}\sin \frac{\pi }{2}\right )\)}{}{}}
\LANGpl{
\Question{
Dane są liczby zespolone
\[ 
 a = 3\left (\cos \frac{\pi }
{3} + \mathrm{i}\sin \frac{\pi }
{3}\right ),\quad b = \sqrt{2}\left (\cos \frac{2\pi }
{3} + \mathrm{i}\sin \frac{2\pi }
{3}\right ),
\]
oblicz iloczyn \(ab\).}
{\(- 3\sqrt{2}\)}{\(3\sqrt{2}\left (\cos \frac{\pi }{2} + \mathrm{i}\sin \frac{\pi }{2}\right )\)}{\(3\sqrt{2}\left (\cos \frac{\pi }{2} -\mathrm{i}\sin \frac{\pi }{2}\right )\)}{\(- 3\sqrt{2}\left (\cos \frac{\pi }{2} + \mathrm{i}\sin \frac{\pi }{2}\right )\)}{}{}}
\LANGcs{
\Question{
Určete součin komplexních čísel
\(a\),
\(b\).
\[ 
 a = 3\left (\cos \frac{\pi }
{3} + \mathrm{i}\sin \frac{\pi }
{3}\right ),\quad b = \sqrt{2}\left (\cos \frac{2\pi }
{3} + \mathrm{i}\sin \frac{2\pi }
{3}\right )
\]}
{\(- 3\sqrt{2}\)}{\(3\sqrt{2}\left (\cos \frac{\pi }{2} + \mathrm{i}\sin \frac{\pi }{2}\right )\)}{\(3\sqrt{2}\left (\cos \frac{\pi }{2} -\mathrm{i}\sin \frac{\pi }{2}\right )\)}{\(- 3\sqrt{2}\left (\cos \frac{\pi }{2} + \mathrm{i}\sin \frac{\pi }{2}\right )\)}{}{}}
\LANGsk{
\Question{
Sú dané komplexné čísla  
\[ 
 a = 3\left (\cos \frac{\pi }
{3} + \mathrm{i}\sin \frac{\pi }
{3}\right ),\quad b = \sqrt{2}\left (\cos \frac{2\pi }
{3} + \mathrm{i}\sin \frac{2\pi }
{3}\right ).
\]
Určte súčin \(ab\).}
{\(- 3\sqrt{2}\)}{\(3\sqrt{2}\left (\cos \frac{\pi }{2} + \mathrm{i}\sin \frac{\pi }{2}\right )\)}{\(3\sqrt{2}\left (\cos \frac{\pi }{2} -\mathrm{i}\sin \frac{\pi }{2}\right )\)}{\(- 3\sqrt{2}\left (\cos \frac{\pi }{2} + \mathrm{i}\sin \frac{\pi }{2}\right )\)}{}{}}
\LANGes{
\Question{
Dados los números complejos
\[ 
 a = 3\left (\cos \frac{\pi }
{3} + \mathrm{i}\sin \frac{\pi }
{3}\right ),\quad b = \sqrt{2}\left (\cos \frac{2\pi }
{3} + \mathrm{i}\sin \frac{2\pi }
{3}\right )
\]
calcula el producto \(ab\).}
{\(- 3\sqrt{2}\)}{\(3\sqrt{2}\left (\cos \frac{\pi }{2} + \mathrm{i}\sin \frac{\pi }{2}\right )\)}{\(3\sqrt{2}\left (\cos \frac{\pi }{2} -\mathrm{i}\sin \frac{\pi }{2}\right )\)}{\(- 3\sqrt{2}\left (\cos \frac{\pi }{2} + \mathrm{i}\sin \frac{\pi }{2}\right )\)}{}{}}
\End

\Data\ID{9000037510}
\Updated{2020-07-18 11:07:24}
\Level{B}
\SubArea{Complex numbers in algebraic and trigonometric form}
\LANGen{
\Question{
Given complex numbers 
\[ 
 a = \left (\cos \frac{\pi }
{3} + \mathrm{i}\sin \frac{\pi }
{3}\right ),\quad b = \sqrt{2}\left (\cos \frac{2\pi }
{3} + \mathrm{i}\sin \frac{2\pi }
{3}\right )
\]
find the quotient \(\frac{a}
{b}\).}
{\(\frac{\sqrt{2}}
 {2} \left (\cos \left (-\frac{\pi }
{3}\right ) + \mathrm{i}\sin \left (-\frac{\pi }
{3}\right )\right )\)}{\(\frac{\sqrt{2}}
 {2} \left (\cos \left (-\frac{\pi }
{3}\right ) -\mathrm{i}\sin \left (-\frac{\pi }
{3}\right )\right )\)}{\(-\frac{\sqrt{2}} 
 {2} \left (\cos \left (-\frac{\pi }
{3}\right ) -\mathrm{i}\sin \left (-\frac{\pi }
{3}\right )\right )\)}{\(-\frac{\sqrt{2}} 
 {2} \left (\cos \left (-\frac{\pi }
{3}\right ) + \mathrm{i}\sin \left (-\frac{\pi }
{3}\right )\right )\)}{}{}}
\LANGpl{
\Question{
Dane są liczby zespolone
\[ 
 a = \left (\cos \frac{\pi }
{3} + \mathrm{i}\sin \frac{\pi }
{3}\right ),\quad b = \sqrt{2}\left (\cos \frac{2\pi }
{3} + \mathrm{i}\sin \frac{2\pi }
{3}\right ),
\]
oblicz iloraz \(\frac{a}
{b}\).}
{\(\frac{\sqrt{2}}
 {2} \left (\cos \left (-\frac{\pi }
{3}\right ) + \mathrm{i}\sin \left (-\frac{\pi }
{3}\right )\right )\)}{\(\frac{\sqrt{2}}
 {2} \left (\cos \left (-\frac{\pi }
{3}\right ) -\mathrm{i}\sin \left (-\frac{\pi }
{3}\right )\right )\)}{\(-\frac{\sqrt{2}} 
 {2} \left (\cos \left (-\frac{\pi }
{3}\right ) -\mathrm{i}\sin \left (-\frac{\pi }
{3}\right )\right )\)}{\(-\frac{\sqrt{2}} 
 {2} \left (\cos \left (-\frac{\pi }
{3}\right ) + \mathrm{i}\sin \left (-\frac{\pi }
{3}\right )\right )\)}{}{}}
\LANGcs{
\Question{
Určete podíl \(\frac{a}
{b}\) 
komplexních čísel \(a\),
\(b\).
\[ 
 a = \left (\cos \frac{\pi }
{3} + \mathrm{i}\sin \frac{\pi }
{3}\right ),\quad b = \sqrt{2}\left (\cos \frac{2\pi }
{3} + \mathrm{i}\sin \frac{2\pi }
{3}\right )
\]}
{\(\frac{\sqrt{2}}
 {2} \left (\cos \left (-\frac{\pi }
{3}\right ) + \mathrm{i}\sin \left (-\frac{\pi }
{3}\right )\right )\)}{\(\frac{\sqrt{2}}
 {2} \left (\cos \left (-\frac{\pi }
{3}\right ) -\mathrm{i}\sin \left (-\frac{\pi }
{3}\right )\right )\)}{\(-\frac{\sqrt{2}} 
 {2} \left (\cos \left (-\frac{\pi }
{3}\right ) -\mathrm{i}\sin \left (-\frac{\pi }
{3}\right )\right )\)}{\(-\frac{\sqrt{2}} 
 {2} \left (\cos \left (-\frac{\pi }
{3}\right ) + \mathrm{i}\sin \left (-\frac{\pi }
{3}\right )\right )\)}{}{}}
\LANGsk{
\Question{
Sú dané komplexné čísla  
\[ 
 a = \left (\cos \frac{\pi }
{3} + \mathrm{i}\sin \frac{\pi }
{3}\right ),\quad b = \sqrt{2}\left (\cos \frac{2\pi }
{3} + \mathrm{i}\sin \frac{2\pi }
{3}\right ).
\]
Určte podiel \(\frac{a}
{b}\).}
{\(\frac{\sqrt{2}}
 {2} \left (\cos \left (-\frac{\pi }
{3}\right ) + \mathrm{i}\sin \left (-\frac{\pi }
{3}\right )\right )\)}{\(\frac{\sqrt{2}}
 {2} \left (\cos \left (-\frac{\pi }
{3}\right ) -\mathrm{i}\sin \left (-\frac{\pi }
{3}\right )\right )\)}{\(-\frac{\sqrt{2}} 
 {2} \left (\cos \left (-\frac{\pi }
{3}\right ) -\mathrm{i}\sin \left (-\frac{\pi }
{3}\right )\right )\)}{\(-\frac{\sqrt{2}} 
 {2} \left (\cos \left (-\frac{\pi }
{3}\right ) + \mathrm{i}\sin \left (-\frac{\pi }
{3}\right )\right )\)}{}{}}
\LANGes{
\Question{
Dados los números complejos
\[ 
 a = \left (\cos \frac{\pi }
{3} + \mathrm{i}\sin \frac{\pi }
{3}\right ),\quad b = \sqrt{2}\left (\cos \frac{2\pi }
{3} + \mathrm{i}\sin \frac{2\pi }
{3}\right )
\]
determina el cociente \(\frac{a}
{b}\).}
{\(\frac{\sqrt{2}}
 {2} \left (\cos \left (-\frac{\pi }
{3}\right ) + \mathrm{i}\sin \left (-\frac{\pi }
{3}\right )\right )\)}{\(\frac{\sqrt{2}}
 {2} \left (\cos \left (-\frac{\pi }
{3}\right ) -\mathrm{i}\sin \left (-\frac{\pi }
{3}\right )\right )\)}{\(-\frac{\sqrt{2}} 
 {2} \left (\cos \left (-\frac{\pi }
{3}\right ) -\mathrm{i}\sin \left (-\frac{\pi }
{3}\right )\right )\)}{\(-\frac{\sqrt{2}} 
 {2} \left (\cos \left (-\frac{\pi }
{3}\right ) + \mathrm{i}\sin \left (-\frac{\pi }
{3}\right )\right )\)}{}{}}
\End

\Data\ID{9000038601}
\Updated{2020-07-18 11:07:52}
\Level{B}
\SubArea{Complex numbers in algebraic and trigonometric form}
\LANGen{
\Question{
Find the polar form of the following complex number. 
\[ 
 -\frac{1}
{2} + \mathrm{i}\frac{\sqrt{3}} 
 {2}
\]}
{\(\cos \frac{2\pi }
{3} + \mathrm{i}\sin \frac{2\pi }
{3}\)}{\(\cos \frac{\pi }
{3} + \mathrm{i}\sin \frac{\pi }
{3}\)}{\(\cos \left (-\frac{\pi }{3}\right ) + \mathrm{i}\sin \left (-\frac{\pi }{3}\right )\)}{\(\cos \frac{3\pi }
{2} + \mathrm{i}\sin \frac{3\pi }
{2}\)}{}{}}
\LANGpl{
\Question{
Wskaż postać biegunową podanej liczby zespolonej.
\[ 
 -\frac{1}
{2} + \mathrm{i}\frac{\sqrt{3}} 
 {2}
\]}
{\(\cos \frac{2\pi }
{3} + \mathrm{i}\sin \frac{2\pi }
{3}\)}{\(\cos \frac{\pi }
{3} + \mathrm{i}\sin \frac{\pi }
{3}\)}{\(\cos \left (-\frac{\pi }{3}\right ) + \mathrm{i}\sin \left (-\frac{\pi }{3}\right )\)}{\(\cos \frac{3\pi }
{2} + \mathrm{i}\sin \frac{3\pi }
{2}\)}{}{}}
\LANGcs{
\Question{
Zapište komplexní číslo \(-\frac{1} 
{2} + \mathrm{i}\frac{\sqrt{3}} 
 {2} \) 
v goniometrickém tvaru.}
{\(\cos \frac{2\pi }
{3} + \mathrm{i}\sin \frac{2\pi }
{3}\)}{\(\cos \frac{\pi }
{3} + \mathrm{i}\sin \frac{\pi }
{3}\)}{\(\cos \left (-\frac{\pi }{3}\right ) + \mathrm{i}\sin \left (-\frac{\pi }{3}\right )\)}{\(\cos \frac{3\pi }
{2} + \mathrm{i}\sin \frac{3\pi }
{2}\)}{}{}}
\LANGsk{
\Question{
Vyjadrite v goniometrickom tvare dané komplexné číslo. 
\[ 
 -\frac{1}
{2} + \mathrm{i}\frac{\sqrt{3}} 
 {2}
\]}
{\(\cos \frac{2\pi }
{3} + \mathrm{i}\sin \frac{2\pi }
{3}\)}{\(\cos \frac{\pi }
{3} + \mathrm{i}\sin \frac{\pi }
{3}\)}{\(\cos \left (-\frac{\pi }{3}\right ) + \mathrm{i}\sin \left (-\frac{\pi }{3}\right )\)}{\(\cos \frac{3\pi }
{2} + \mathrm{i}\sin \frac{3\pi }
{2}\)}{}{}}
\LANGes{
\Question{
Determina la forma polar del siguiente número complejo. 
\[ 
 -\frac{1}
{2} + \mathrm{i}\frac{\sqrt{3}} 
 {2}
\]}
{\(\cos \frac{2\pi }
{3} + \mathrm{i}\sin \frac{2\pi }
{3}\)}{\(\cos \frac{\pi }
{3} + \mathrm{i}\sin \frac{\pi }
{3}\)}{\(\cos \left (-\frac{\pi }{3}\right ) + \mathrm{i}\sin \left (-\frac{\pi }{3}\right )\)}{\(\cos \frac{3\pi }
{2} + \mathrm{i}\sin \frac{3\pi }
{2}\)}{}{}}
\End

\Data\ID{9000038602}
\Updated{2020-07-18 11:07:09}
\Level{B}
\SubArea{Complex numbers in algebraic and trigonometric form}
\LANGen{
\Question{
Find the polar form of the following complex number. 
\[ 
 \frac{1} 
{2} + \mathrm{i}\frac{\sqrt{3}} 
 {2}
\]}
{\(\cos \frac{\pi }{3} + \mathrm{i}\sin \frac{\pi }{3}\)}{\(\cos \frac{2\pi }
{3} + \mathrm{i}\sin \frac{2\pi }
{3}\)}{\(\cos \frac{3\pi }
{2} + \mathrm{i}\sin \frac{3\pi }
{2}\)}{\(\cos \left (-\frac{\pi }{3}\right ) + \mathrm{i}\sin \left (-\frac{\pi }{3}\right )\)}{}{}}
\LANGpl{
\Question{
Wskaż postać biegunową  liczby zespolonej.
\[ 
 \frac{1} 
{2} + \mathrm{i}\frac{\sqrt{3}} 
 {2}
\]}
{\(\cos \frac{\pi }{3} + \mathrm{i}\sin \frac{\pi }{3}\)}{\(\cos \frac{2\pi }
{3} + \mathrm{i}\sin \frac{2\pi }
{3}\)}{\(\cos \frac{3\pi }
{2} + \mathrm{i}\sin \frac{3\pi }
{2}\)}{\(\cos \left (-\frac{\pi }{3}\right ) + \mathrm{i}\sin \left (-\frac{\pi }{3}\right )\)}{}{}}
\LANGcs{
\Question{
Zapište komplexní číslo \(\frac{1}
{2} + \mathrm{i}\frac{\sqrt{3}} 
 {2} \) 
v goniometrickém tvaru.}
{\(\cos \frac{\pi }{3} + \mathrm{i}\sin \frac{\pi }{3}\)}{\(\cos \frac{2\pi }
{3} + \mathrm{i}\sin \frac{2\pi }
{3}\)}{\(\cos \frac{3\pi }
{2} + \mathrm{i}\sin \frac{3\pi }
{2}\)}{\(\cos \left (-\frac{\pi }{3}\right ) + \mathrm{i}\sin \left (-\frac{\pi }{3}\right )\)}{}{}}
\LANGsk{
\Question{
Vyjadrite v goniometrickom tvare dané komplexné číslo. 
\[ 
 \frac{1} 
{2} + \mathrm{i}\frac{\sqrt{3}} 
 {2}
\]}
{\(\cos \frac{\pi }{3} + \mathrm{i}\sin \frac{\pi }{3}\)}{\(\cos \frac{2\pi }
{3} + \mathrm{i}\sin \frac{2\pi }
{3}\)}{\(\cos \frac{3\pi }
{2} + \mathrm{i}\sin \frac{3\pi }
{2}\)}{\(\cos \left (-\frac{\pi }{3}\right ) + \mathrm{i}\sin \left (-\frac{\pi }{3}\right )\)}{}{}}
\LANGes{
\Question{
Determina la forma polar del siguiente número complejo.
\[ 
 \frac{1} 
{2} + \mathrm{i}\frac{\sqrt{3}} 
 {2}
\]}
{\(\cos \frac{\pi }{3} + \mathrm{i}\sin \frac{\pi }{3}\)}{\(\cos \frac{2\pi }
{3} + \mathrm{i}\sin \frac{2\pi }
{3}\)}{\(\cos \frac{3\pi }
{2} + \mathrm{i}\sin \frac{3\pi }
{2}\)}{\(\cos \left (-\frac{\pi }{3}\right ) + \mathrm{i}\sin \left (-\frac{\pi }{3}\right )\)}{}{}}
\End

\Data\ID{9000038603}
\Updated{2020-07-18 11:07:24}
\Level{B}
\SubArea{Complex numbers in algebraic and trigonometric form}
\LANGen{
\Question{
Find the polar form of the following complex number. 
\[ 
 \frac{\sqrt{2}} 
 {2} + \mathrm{i}\frac{\sqrt{6}} 
 {2}
\]}
{\(\sqrt{2}\left (\cos \frac{\pi }{3} + \mathrm{i}\sin \frac{\pi }{3}\right )\)}{\(\sqrt{2}\left (\cos \frac{2\pi }
{3} + \mathrm{i}\sin \frac{2\pi }
{3}\right )\)}{\(2\left (\cos \frac{4\pi }
{3} + \mathrm{i}\sin \frac{4\pi }
{3}\right )\)}{\(2\left (\cos \frac{3\pi }
{2} + \mathrm{i}\sin \frac{3\pi }
{2}\right )\)}{}{}}
\LANGpl{
\Question{
Wskaż postać biegunową podanej  liczby zespolonej.
\[ 
 \frac{\sqrt{2}} 
 {2} + \mathrm{i}\frac{\sqrt{6}} 
 {2}
\]}
{\(\sqrt{2}\left (\cos \frac{\pi }{3} + \mathrm{i}\sin \frac{\pi }{3}\right )\)}{\(\sqrt{2}\left (\cos \frac{2\pi }
{3} + \mathrm{i}\sin \frac{2\pi }
{3}\right )\)}{\(2\left (\cos \frac{4\pi }
{3} + \mathrm{i}\sin \frac{4\pi }
{3}\right )\)}{\(2\left (\cos \frac{3\pi }
{2} + \mathrm{i}\sin \frac{3\pi }
{2}\right )\)}{}{}}
\LANGcs{
\Question{
Zapište komplexní číslo \(\frac{\sqrt{2}}
 {2} + \mathrm{i}\frac{\sqrt{6}} 
 {2} \) 
v goniometrickém tvaru.}
{\(\sqrt{2}\left (\cos \frac{\pi }{3} + \mathrm{i}\sin \frac{\pi }{3}\right )\)}{\(\sqrt{2}\left (\cos \frac{2\pi }
{3} + \mathrm{i}\sin \frac{2\pi }
{3}\right )\)}{\(2\left (\cos \frac{4\pi }
{3} + \mathrm{i}\sin \frac{4\pi }
{3}\right )\)}{\(2\left (\cos \frac{3\pi }
{2} + \mathrm{i}\sin \frac{3\pi }
{2}\right )\)}{}{}}
\LANGsk{
\Question{
Vyjadrite v goniometrickom tvare dané komplexné číslo. 
\[ 
 \frac{\sqrt{2}} 
 {2} + \mathrm{i}\frac{\sqrt{6}} 
 {2}
\]}
{\(\sqrt{2}\left (\cos \frac{\pi }{3} + \mathrm{i}\sin \frac{\pi }{3}\right )\)}{\(\sqrt{2}\left (\cos \frac{2\pi }
{3} + \mathrm{i}\sin \frac{2\pi }
{3}\right )\)}{\(2\left (\cos \frac{4\pi }
{3} + \mathrm{i}\sin \frac{4\pi }
{3}\right )\)}{\(2\left (\cos \frac{3\pi }
{2} + \mathrm{i}\sin \frac{3\pi }
{2}\right )\)}{}{}}
\LANGes{
\Question{
Determina la forma polar del siguiente número complejo.
\[ 
 \frac{\sqrt{2}} 
 {2} + \mathrm{i}\frac{\sqrt{6}} 
 {2}
\]}
{\(\sqrt{2}\left (\cos \frac{\pi }{3} + \mathrm{i}\sin \frac{\pi }{3}\right )\)}{\(\sqrt{2}\left (\cos \frac{2\pi }
{3} + \mathrm{i}\sin \frac{2\pi }
{3}\right )\)}{\(2\left (\cos \frac{4\pi }
{3} + \mathrm{i}\sin \frac{4\pi }
{3}\right )\)}{\(2\left (\cos \frac{3\pi }
{2} + \mathrm{i}\sin \frac{3\pi }
{2}\right )\)}{}{}}
\End

\Data\ID{9000038604}
\Updated{2020-07-18 11:07:23}
\Level{B}
\SubArea{Complex numbers in algebraic and trigonometric form}
\LANGen{
\Question{
Find the polar form of the following complex number. 
\[ 
 \frac{\sqrt{3}} 
{\sqrt{2}} + \mathrm{i}\frac{\sqrt{3}} 
{\sqrt{2}}
\]}
{\(\sqrt{3}\left (\cos \frac{\pi }{4} + \mathrm{i}\sin \frac{\pi }{4}\right )\)}{\(\sqrt{3}\left (\cos \frac{3\pi }
{4} + \mathrm{i}\sin \frac{3\pi }
{4}\right )\)}{\(\sqrt{2}\left (\cos \frac{\pi }{3} + \mathrm{i}\sin \frac{\pi }{3}\right )\)}{\(\sqrt{2}\left (\cos \frac{2\pi }
{3} + \mathrm{i}\sin \frac{2\pi }
{3}\right )\)}{}{}}
\LANGpl{
\Question{
Wskaż postać biegunową podanej  liczby zespolonej.
\[ 
 \frac{\sqrt{3}} 
{\sqrt{2}} + \mathrm{i}\frac{\sqrt{3}} 
{\sqrt{2}}
\]}
{\(\sqrt{3}\left (\cos \frac{\pi }{4} + \mathrm{i}\sin \frac{\pi }{4}\right )\)}{\(\sqrt{3}\left (\cos \frac{3\pi }
{4} + \mathrm{i}\sin \frac{3\pi }
{4}\right )\)}{\(\sqrt{2}\left (\cos \frac{\pi }{3} + \mathrm{i}\sin \frac{\pi }{3}\right )\)}{\(\sqrt{2}\left (\cos \frac{2\pi }
{3} + \mathrm{i}\sin \frac{2\pi }
{3}\right )\)}{}{}}
\LANGcs{
\Question{
Zapište komplexní číslo \(\frac{\sqrt{3}}
{\sqrt{2}} + \mathrm{i}\frac{\sqrt{3}} 
{\sqrt{2}}\) 
v goniometrickém tvaru.}
{\(\sqrt{3}\left (\cos \frac{\pi }{4} + \mathrm{i}\sin \frac{\pi }{4}\right )\)}{\(\sqrt{3}\left (\cos \frac{3\pi }
{4} + \mathrm{i}\sin \frac{3\pi }
{4}\right )\)}{\(\sqrt{2}\left (\cos \frac{\pi }{3} + \mathrm{i}\sin \frac{\pi }{3}\right )\)}{\(\sqrt{2}\left (\cos \frac{2\pi }
{3} + \mathrm{i}\sin \frac{2\pi }
{3}\right )\)}{}{}}
\LANGsk{
\Question{
Vyjadrite v goniometrickom tvare dané komplexné číslo. 
\[ 
 \frac{\sqrt{3}} 
{\sqrt{2}} + \mathrm{i}\frac{\sqrt{3}} 
{\sqrt{2}}
\]}
{\(\sqrt{3}\left (\cos \frac{\pi }{4} + \mathrm{i}\sin \frac{\pi }{4}\right )\)}{\(\sqrt{3}\left (\cos \frac{3\pi }
{4} + \mathrm{i}\sin \frac{3\pi }
{4}\right )\)}{\(\sqrt{2}\left (\cos \frac{\pi }{3} + \mathrm{i}\sin \frac{\pi }{3}\right )\)}{\(\sqrt{2}\left (\cos \frac{2\pi }
{3} + \mathrm{i}\sin \frac{2\pi }
{3}\right )\)}{}{}}
\LANGes{
\Question{
Determina la forma polar del siguiente número complejo.
\[ 
 \frac{\sqrt{3}} 
{\sqrt{2}} + \mathrm{i}\frac{\sqrt{3}} 
{\sqrt{2}}
\]}
{\(\sqrt{3}\left (\cos \frac{\pi }{4} + \mathrm{i}\sin \frac{\pi }{4}\right )\)}{\(\sqrt{3}\left (\cos \frac{3\pi }
{4} + \mathrm{i}\sin \frac{3\pi }
{4}\right )\)}{\(\sqrt{2}\left (\cos \frac{\pi }{3} + \mathrm{i}\sin \frac{\pi }{3}\right )\)}{\(\sqrt{2}\left (\cos \frac{2\pi }
{3} + \mathrm{i}\sin \frac{2\pi }
{3}\right )\)}{}{}}
\End

\Data\ID{9000038605}
\Updated{2020-07-18 11:07:33}
\Level{B}
\SubArea{Complex numbers in algebraic and trigonometric form}
\LANGen{
\Question{
Find the polar form of the following complex number. 
\[ 
 -\frac{\sqrt{5}}
 {2} + \mathrm{i}\frac{\sqrt{15}} 
 {2}
\]}
{\(\sqrt{5}\left (\cos \frac{2\pi }
{3} + \mathrm{i}\sin \frac{2\pi }
{3}\right )\)}{\(\sqrt{5}\left (\cos \frac{\pi }{3} + \mathrm{i}\sin \frac{\pi }{3}\right )\)}{\(\sqrt{5}\left (\cos \frac{2\pi }
{5} + \mathrm{i}\sin \frac{2\pi }
{5}\right )\)}{\(\sqrt{5}\left (\cos \frac{3\pi }
{2} + \mathrm{i}\sin \frac{3\pi }
{2}\right )\)}{}{}}
\LANGpl{
\Question{
Wskaż postać biegunową podanej  liczby zespolonej.
\[ 
 -\frac{\sqrt{5}}
 {2} + \mathrm{i}\frac{\sqrt{15}} 
 {2}
\]}
{\(\sqrt{5}\left (\cos \frac{2\pi }
{3} + \mathrm{i}\sin \frac{2\pi }
{3}\right )\)}{\(\sqrt{5}\left (\cos \frac{\pi }{3} + \mathrm{i}\sin \frac{\pi }{3}\right )\)}{\(\sqrt{5}\left (\cos \frac{2\pi }
{5} + \mathrm{i}\sin \frac{2\pi }
{5}\right )\)}{\(\sqrt{5}\left (\cos \frac{3\pi }
{2} + \mathrm{i}\sin \frac{3\pi }
{2}\right )\)}{}{}}
\LANGcs{
\Question{
Zapište komplexní číslo \(-\frac{\sqrt{5}} 
 {2} + \mathrm{i}\frac{\sqrt{15}} 
 {2} \) 
v goniometrickém tvaru.}
{\(\sqrt{5}\left (\cos \frac{2\pi }
{3} + \mathrm{i}\sin \frac{2\pi }
{3}\right )\)}{\(\sqrt{5}\left (\cos \frac{\pi }{3} + \mathrm{i}\sin \frac{\pi }{3}\right )\)}{\(\sqrt{5}\left (\cos \frac{2\pi }
{5} + \mathrm{i}\sin \frac{2\pi }
{5}\right )\)}{\(\sqrt{5}\left (\cos \frac{3\pi }
{2} + \mathrm{i}\sin \frac{3\pi }
{2}\right )\)}{}{}}
\LANGsk{
\Question{
Vyjadrite v goniometrickom tvare dané komplexné číslo. 
\[ 
 -\frac{\sqrt{5}}
 {2} + \mathrm{i}\frac{\sqrt{15}} 
 {2}
\]}
{\(\sqrt{5}\left (\cos \frac{2\pi }
{3} + \mathrm{i}\sin \frac{2\pi }
{3}\right )\)}{\(\sqrt{5}\left (\cos \frac{\pi }{3} + \mathrm{i}\sin \frac{\pi }{3}\right )\)}{\(\sqrt{5}\left (\cos \frac{2\pi }
{5} + \mathrm{i}\sin \frac{2\pi }
{5}\right )\)}{\(\sqrt{5}\left (\cos \frac{3\pi }
{2} + \mathrm{i}\sin \frac{3\pi }
{2}\right )\)}{}{}}
\LANGes{
\Question{
Determina la forma polar del siguiente número complejo.
\[ 
 -\frac{\sqrt{5}}
 {2} + \mathrm{i}\frac{\sqrt{15}} 
 {2}
\]}
{\(\sqrt{5}\left (\cos \frac{2\pi }
{3} + \mathrm{i}\sin \frac{2\pi }
{3}\right )\)}{\(\sqrt{5}\left (\cos \frac{\pi }{3} + \mathrm{i}\sin \frac{\pi }{3}\right )\)}{\(\sqrt{5}\left (\cos \frac{2\pi }
{5} + \mathrm{i}\sin \frac{2\pi }
{5}\right )\)}{\(\sqrt{5}\left (\cos \frac{3\pi }
{2} + \mathrm{i}\sin \frac{3\pi }
{2}\right )\)}{}{}}
\End

\Data\ID{9000038606}
\Updated{2020-07-18 11:07:01}
\Level{B}
\SubArea{Complex numbers in algebraic and trigonometric form}
\LANGen{
\Question{
Find the algebraic form of the following complex number. 
\[ 
 \cos \frac{\pi }
{4} + \mathrm{i}\sin \frac{\pi }
{4}
\]}
{\(\frac{\sqrt{2}}
 {2} + \mathrm{i}\frac{\sqrt{2}} 
 {2} \)}{\(\frac{\sqrt{2}}
 {2} -\mathrm{i}\frac{\sqrt{2}} 
 {2} \)}{\(\frac{\sqrt{3}}
 {2} + \mathrm{i}\frac{\sqrt{3}} 
 {2} \)}{\(\frac{\sqrt{3}}
 {2} -\mathrm{i}\frac{\sqrt{3}} 
 {2} \)}{}{}}
\LANGpl{
\Question{
Wskaż postać algebraiczną podanej  liczby zespolonej.
\[ 
 \cos \frac{\pi }
{4} + \mathrm{i}\sin \frac{\pi }
{4}
\]}
{\(\frac{\sqrt{2}}
 {2} + \mathrm{i}\frac{\sqrt{2}} 
 {2} \)}{\(\frac{\sqrt{2}}
 {2} -\mathrm{i}\frac{\sqrt{2}} 
 {2} \)}{\(\frac{\sqrt{3}}
 {2} + \mathrm{i}\frac{\sqrt{3}} 
 {2} \)}{\(\frac{\sqrt{3}}
 {2} -\mathrm{i}\frac{\sqrt{3}} 
 {2} \)}{}{}}
\LANGcs{
\Question{
Zapište komplexní číslo \(\cos \frac{\pi }{4} + \mathrm{i}\sin \frac{\pi }{4}\) 
v algebraickém tvaru.}
{\(\frac{\sqrt{2}}
 {2} + \mathrm{i}\frac{\sqrt{2}} 
 {2} \)}{\(\frac{\sqrt{2}}
 {2} -\mathrm{i}\frac{\sqrt{2}} 
 {2} \)}{\(\frac{\sqrt{3}}
 {2} + \mathrm{i}\frac{\sqrt{3}} 
 {2} \)}{\(\frac{\sqrt{3}}
 {2} -\mathrm{i}\frac{\sqrt{3}} 
 {2} \)}{}{}}
\LANGsk{
\Question{
Vyjadrite v algebraickom tvare dané komplexné číslo. 
\[ 
 \cos \frac{\pi }
{4} + \mathrm{i}\sin \frac{\pi }
{4}
\]}
{\(\frac{\sqrt{2}}
 {2} + \mathrm{i}\frac{\sqrt{2}} 
 {2} \)}{\(\frac{\sqrt{2}}
 {2} -\mathrm{i}\frac{\sqrt{2}} 
 {2} \)}{\(\frac{\sqrt{3}}
 {2} + \mathrm{i}\frac{\sqrt{3}} 
 {2} \)}{\(\frac{\sqrt{3}}
 {2} -\mathrm{i}\frac{\sqrt{3}} 
 {2} \)}{}{}}
\LANGes{
\Question{
Determina la forma algebraica del siguiente número complejo.
\[ 
 \cos \frac{\pi }
{4} + \mathrm{i}\sin \frac{\pi }
{4}
\]}
{\(\frac{\sqrt{2}}
 {2} + \mathrm{i}\frac{\sqrt{2}} 
 {2} \)}{\(\frac{\sqrt{2}}
 {2} -\mathrm{i}\frac{\sqrt{2}} 
 {2} \)}{\(\frac{\sqrt{3}}
 {2} + \mathrm{i}\frac{\sqrt{3}} 
 {2} \)}{\(\frac{\sqrt{3}}
 {2} -\mathrm{i}\frac{\sqrt{3}} 
 {2} \)}{}{}}
\End

\Data\ID{9000038607}
\Updated{2020-07-18 11:07:14}
\Level{B}
\SubArea{Complex numbers in algebraic and trigonometric form}
\LANGen{
\Question{
Find the algebraic form of the following complex number. 
\[ 
 3\left (\cos \frac{\pi }
{2} + \mathrm{i}\sin \frac{\pi }
{2}\right )
\]}
{\(3\mathrm{i}\)}{\(- 3\mathrm{i}\)}{\(3 + 3\mathrm{i}\)}{\(3 - 3\mathrm{i}\)}{}{}}
\LANGpl{
\Question{
Wskaż postać algebraiczną podanej  liczby zespolonej.
\[ 
 3\left (\cos \frac{\pi }
{2} + \mathrm{i}\sin \frac{\pi }
{2}\right )
\]}
{\(3\mathrm{i}\)}{\(- 3\mathrm{i}\)}{\(3 + 3\mathrm{i}\)}{\(3 - 3\mathrm{i}\)}{}{}}
\LANGcs{
\Question{
Zapište komplexní číslo \(3\left (\cos \frac{\pi }{2} + \mathrm{i}\sin \frac{\pi }{2}\right )\) 
v algebraickém tvaru.}
{\(3\mathrm{i}\)}{\(- 3\mathrm{i}\)}{\(3 + 3\mathrm{i}\)}{\(3 - 3\mathrm{i}\)}{}{}}
\LANGsk{
\Question{
Vyjadrite v  algebraickom tvare dané komplexné číslo. 
\[ 
 3\left (\cos \frac{\pi }
{2} + \mathrm{i}\sin \frac{\pi }
{2}\right )
\]}
{\(3\mathrm{i}\)}{\(- 3\mathrm{i}\)}{\(3 + 3\mathrm{i}\)}{\(3 - 3\mathrm{i}\)}{}{}}
\LANGes{
\Question{
Determina la forma algebraica del siguiente número complejo.
\[ 
 3\left (\cos \frac{\pi }
{2} + \mathrm{i}\sin \frac{\pi }
{2}\right )
\]}
{\(3\mathrm{i}\)}{\(- 3\mathrm{i}\)}{\(3 + 3\mathrm{i}\)}{\(3 - 3\mathrm{i}\)}{}{}}
\End

\Data\ID{9000038608}
\Updated{2020-07-18 11:07:26}
\Level{B}
\SubArea{Complex numbers in algebraic and trigonometric form}
\LANGen{
\Question{
Find the algebraic form of the following complex number. 
\[ 
 8(\cos \pi + \mathrm{i}\sin \pi )
\]}
{\(- 8\)}{\(8\)}{\(8 + 8\mathrm{i}\)}{\(8 - 8\mathrm{i}\)}{}{}}
\LANGpl{
\Question{
Wskaż postać algebraiczną podanej  liczby zespolonej.
\[ 
 8(\cos \pi + \mathrm{i}\sin \pi )
\]}
{\(- 8\)}{\(8\)}{\(8 + 8\mathrm{i}\)}{\(8 - 8\mathrm{i}\)}{}{}}
\LANGcs{
\Question{
Zapište komplexní číslo \(8(\cos \pi + \mathrm{i}\sin \pi )\) 
v algebraickém tvaru.}
{\(- 8\)}{\(8\)}{\(8 + 8\mathrm{i}\)}{\(8 - 8\mathrm{i}\)}{}{}}
\LANGsk{
\Question{
Vyjadrite v algebraickom tvare dané komplexné číslo. 
\[ 
 8(\cos \pi + \mathrm{i}\sin \pi )
\]}
{\(- 8\)}{\(8\)}{\(8 + 8\mathrm{i}\)}{\(8 - 8\mathrm{i}\)}{}{}}
\LANGes{
\Question{
Determina la forma algebraica del siguiente número complejo.
\[ 
 8(\cos \pi + \mathrm{i}\sin \pi )
\]}
{\(- 8\)}{\(8\)}{\(8 + 8\mathrm{i}\)}{\(8 - 8\mathrm{i}\)}{}{}}
\End

\Data\ID{9000038609}
\Updated{2020-07-18 11:07:37}
\Level{B}
\SubArea{Complex numbers in algebraic and trigonometric form}
\LANGen{
\Question{
Find the algebraic form of the following complex number. 
\[ 
 5\left (\cos \frac{3\pi }
{4} + \mathrm{i}\sin \frac{3\pi }
{4}\right )
\]}
{\(-\frac{5\sqrt{2}} 
 {2} + \mathrm{i}\frac{5\sqrt{2}} 
 {2} \)}{\(\frac{5\sqrt{2}}
 {2} -\mathrm{i}\frac{5\sqrt{2}} 
 {2} \)}{\(\frac{5}
{2} + \mathrm{i}\frac{5} 
{2}\)}{\(\frac{5}
{2} -\mathrm{i}\frac{5} 
{2}\)}{}{}}
\LANGpl{
\Question{
Wskaż postać algebraiczną podanej  liczby zespolonej.
\[ 
 5\left (\cos \frac{3\pi }
{4} + \mathrm{i}\sin \frac{3\pi }
{4}\right )
\]}
{\(-\frac{5\sqrt{2}} 
 {2} + \mathrm{i}\frac{5\sqrt{2}} 
 {2} \)}{\(\frac{5\sqrt{2}}
 {2} -\mathrm{i}\frac{5\sqrt{2}} 
 {2} \)}{\(\frac{5}
{2} + \mathrm{i}\frac{5} 
{2}\)}{\(\frac{5}
{2} -\mathrm{i}\frac{5} 
{2}\)}{}{}}
\LANGcs{
\Question{
Zapište komplexní číslo \(5\left (\cos \frac{3\pi }
{4} + \mathrm{i}\sin \frac{3\pi }
{4}\right )\) 
v algebraickém tvaru.}
{\(-\frac{5\sqrt{2}} 
 {2} + \mathrm{i}\frac{5\sqrt{2}} 
 {2} \)}{\(\frac{5\sqrt{2}}
 {2} -\mathrm{i}\frac{5\sqrt{2}} 
 {2} \)}{\(\frac{5}
{2} + \mathrm{i}\frac{5} 
{2}\)}{\(\frac{5}
{2} -\mathrm{i}\frac{5} 
{2}\)}{}{}}
\LANGsk{
\Question{
Vyjadrite v algebraickom tvare dané komplexné číslo. 
\[ 
 5\left (\cos \frac{3\pi }
{4} + \mathrm{i}\sin \frac{3\pi }
{4}\right )
\]}
{\(-\frac{5\sqrt{2}} 
 {2} + \mathrm{i}\frac{5\sqrt{2}} 
 {2} \)}{\(\frac{5\sqrt{2}}
 {2} -\mathrm{i}\frac{5\sqrt{2}} 
 {2} \)}{\(\frac{5}
{2} + \mathrm{i}\frac{5} 
{2}\)}{\(\frac{5}
{2} -\mathrm{i}\frac{5} 
{2}\)}{}{}}
\LANGes{
\Question{
Determina la forma algebraica del siguiente número complejo.
\[ 
 5\left (\cos \frac{3\pi }
{4} + \mathrm{i}\sin \frac{3\pi }
{4}\right )
\]}
{\(-\frac{5\sqrt{2}} 
 {2} + \mathrm{i}\frac{5\sqrt{2}} 
 {2} \)}{\(\frac{5\sqrt{2}}
 {2} -\mathrm{i}\frac{5\sqrt{2}} 
 {2} \)}{\(\frac{5}
{2} + \mathrm{i}\frac{5} 
{2}\)}{\(\frac{5}
{2} -\mathrm{i}\frac{5} 
{2}\)}{}{}}
\End

\Data\ID{9000038610}
\Updated{2020-07-18 11:07:47}
\Level{B}
\SubArea{Complex numbers in algebraic and trigonometric form}
\LANGen{
\Question{
Find the algebraic form of the following complex number. 
\[ 
 2\left (\cos \frac{3\pi }
{4} + \mathrm{i}\sin \frac{3\pi }
{4}\right )
\]}
{\(-\sqrt{2} + \mathrm{i}\sqrt{2}\)}{\(\sqrt{2} -\mathrm{i}\sqrt{2}\)}{\(2 + 2\mathrm{i}\)}{\(2 - 2\mathrm{i}\)}{}{}}
\LANGpl{
\Question{
Wskaż postać algebraiczną podanej  liczby zespolonej.
\[ 
 2\left (\cos \frac{3\pi }
{4} + \mathrm{i}\sin \frac{3\pi }
{4}\right )
\]}
{\(-\sqrt{2} + \mathrm{i}\sqrt{2}\)}{\(\sqrt{2} -\mathrm{i}\sqrt{2}\)}{\(2 + 2\mathrm{i}\)}{\(2 - 2\mathrm{i}\)}{}{}}
\LANGcs{
\Question{
Zapište komplexní číslo \(2\left (\cos \frac{3\pi }
{4} + \mathrm{i}\sin \frac{3\pi }
{4}\right )\) 
v algebraickém tvaru.}
{\(-\sqrt{2} + \mathrm{i}\sqrt{2}\)}{\(\sqrt{2} -\mathrm{i}\sqrt{2}\)}{\(2 + 2\mathrm{i}\)}{\(2 - 2\mathrm{i}\)}{}{}}
\LANGsk{
\Question{
Vyjadrite v algebraickom tvare dané komplexné číslo. 
\[ 
 2\left (\cos \frac{3\pi }
{4} + \mathrm{i}\sin \frac{3\pi }
{4}\right )
\]}
{\(-\sqrt{2} + \mathrm{i}\sqrt{2}\)}{\(\sqrt{2} -\mathrm{i}\sqrt{2}\)}{\(2 + 2\mathrm{i}\)}{\(2 - 2\mathrm{i}\)}{}{}}
\LANGes{
\Question{
Determina la forma algebraica del siguiente número complejo.
\[ 
 2\left (\cos \frac{3\pi }
{4} + \mathrm{i}\sin \frac{3\pi }
{4}\right )
\]}
{\(-\sqrt{2} + \mathrm{i}\sqrt{2}\)}{\(\sqrt{2} -\mathrm{i}\sqrt{2}\)}{\(2 + 2\mathrm{i}\)}{\(2 - 2\mathrm{i}\)}{}{}}
\End

\Data\ID{9000039101}
\Updated{2022-04-23 05:04:17}
\Level{B}
\SubArea{Complex numbers in algebraic and trigonometric form}
\LANGen{
\Question{
Find the polar form of the complex number
\(z=\frac{\mathrm{i}^{14}-1}
 {\mathrm{i}^{9}+1} \).}
{\(\sqrt{2}\left (\cos \frac{3\pi }
{4} + \mathrm{i}\sin \frac{3\pi }
{4}\right )\)}{\(\sqrt{2}\left (\cos \frac{5\pi }
{4} + \mathrm{i}\sin \frac{5\pi }
{4}\right )\)}{\(\sqrt{2}\left (\cos \frac{\pi }{4} + \mathrm{i}\sin \frac{\pi }{4}\right )\)}{\(\sqrt{2}\left (\cos \frac{7\pi }
{4} + \mathrm{i}\sin \frac{7\pi }
{4}\right )\)}{}{}}
\LANGpl{
\Question{
Wskaż postać biegunową podanej  liczby zespolonej
\(z=\frac{\mathrm{i}^{14}-1}
 {\mathrm{i}^{9}+1} \).}
{\(\sqrt{2}\left (\cos \frac{3\pi }
{4} + \mathrm{i}\sin \frac{3\pi }
{4}\right )\)}{\(\sqrt{2}\left (\cos \frac{5\pi }
{4} + \mathrm{i}\sin \frac{5\pi }
{4}\right )\)}{\(\sqrt{2}\left (\cos \frac{\pi }{4} + \mathrm{i}\sin \frac{\pi }{4}\right )\)}{\(\sqrt{2}\left (\cos \frac{7\pi }
{4} + \mathrm{i}\sin \frac{7\pi }
{4}\right )\)}{}{}}
\LANGcs{
\Question{
Goniometrický tvar komplexního čísla
\(z=\frac{\mathrm{i}^{14}-1}
 {\mathrm{i}^{9}+1} \) je:}
{\(\sqrt{2}\left (\cos \frac{3\pi }
{4} + \mathrm{i}\sin \frac{3\pi }
{4}\right )\)}{\( \sqrt{2}\left (\cos \frac{5\pi }
{4} + \mathrm{i}\sin \frac{5\pi }
{4}\right )\)}{\( \sqrt{2}\left (\cos \frac{\pi }{4} + \mathrm{i}\sin \frac{\pi }{4}\right )\)}{\( \sqrt{2}\left (\cos \frac{7\pi }
{4} + \mathrm{i}\sin \frac{7\pi }
{4}\right )\)}{}{}}
\LANGsk{
\Question{
Goniometrický tvar komplexného čísla
\(z=\frac{\mathrm{i}^{14}-1} 
 {\mathrm{i}^{9}+1} \) je:}
{\(\sqrt{2}\left (\cos \frac{3\pi }
{4} + \mathrm{i}\sin \frac{3\pi }
{4}\right )\)}{\(\sqrt{2}\left (\cos \frac{5\pi }
{4} + \mathrm{i}\sin \frac{5\pi }
{4}\right )\)}{\(\sqrt{2}\left (\cos \frac{\pi }{4} + \mathrm{i}\sin \frac{\pi }{4}\right )\)}{\(\sqrt{2}\left (\cos \frac{7\pi }
{4} + \mathrm{i}\sin \frac{7\pi }
{4}\right )\)}{}{}}
\LANGes{
\Question{
Determina la forma polar del siguiente número complejo.
\(z=\frac{\mathrm{i}^{14}-1}
 {\mathrm{i}^{9}+1} \).}
{\(\sqrt{2}\left (\cos \frac{3\pi }
{4} + \mathrm{i}\sin \frac{3\pi }
{4}\right )\)}{\(\sqrt{2}\left (\cos \frac{5\pi }
{4} + \mathrm{i}\sin \frac{5\pi }
{4}\right )\)}{\(\sqrt{2}\left (\cos \frac{\pi }{4} + \mathrm{i}\sin \frac{\pi }{4}\right )\)}{\(\sqrt{2}\left (\cos \frac{7\pi }
{4} + \mathrm{i}\sin \frac{7\pi }
{4}\right )\)}{}{}}
\End

\Data\ID{9000039102}
\Updated{2020-12-03 06:12:21}
\Level{B}
\SubArea{Complex numbers in algebraic and trigonometric form}
\LANGen{
\Question{
Identify a complex number with absolute value different from
\(1\).}
{\(1 + \mathrm{i}\)}{\(\frac{1}
{2} -\frac{\sqrt{3}} 
 {2} \mathrm{i}\)}{\(-\frac{3} 
{5} -\frac{4} 
{5}\mathrm{i}\)}{\(-\mathrm{i}\)}{}{}}
\LANGpl{
\Question{
Wskaż liczbę zespoloną, której wartość bezwzględna jest różna od 
\(1\).}
{\(1 + \mathrm{i}\)}{\(\frac{1}
{2} -\frac{\sqrt{3}} 
 {2} \mathrm{i}\)}{\(-\frac{3} 
{5} -\frac{4} 
{5}\mathrm{i}\)}{\(-\mathrm{i}\)}{}{}}
\LANGcs{
\Question{
Které z následujících komplexních čísel není komplexní jednotkou? (Poznámka: Termínem komplexní jednotka se označují komplexní čísla, jejichž absolutní hodnota je rovna jedné.)}
{\(1 + \mathrm{i}\)}{\(\frac{1}
{2} -\frac{\sqrt{3}} 
 {2} \mathrm{i}\)}{\(-\frac{3} 
{5} -\frac{4} 
{5}\mathrm{i}\)}{\(-\mathrm{i}\)}{}{}}
\LANGsk{
\Question{
Ktoré z nasledujúcich komplexných čísel nie je komplexnou jednotkou? (Termínom komplexná jednotka sa označujú komplexné čísla, ktorých absolútna hodnota je rovná jednej.)}
{\(1 + \mathrm{i}\)}{\(\frac{1}
{2} -\frac{\sqrt{3}} 
 {2} \mathrm{i}\)}{\(-\frac{3} 
{5} -\frac{4} 
{5}\mathrm{i}\)}{\(-\mathrm{i}\)}{}{}}
\LANGes{
\Question{
Identifica un número complejo con un valor absoluto diferente de
\(1\).}
{\(1 + \mathrm{i}\)}{\(\frac{1}
{2} -\frac{\sqrt{3}} 
 {2} \mathrm{i}\)}{\(-\frac{3} 
{5} -\frac{4} 
{5}\mathrm{i}\)}{\(-\mathrm{i}\)}{}{}}
\End

\Data\ID{9000070110}
\Updated{2022-04-23 05:04:17}
\Level{B}
\SubArea{Complex numbers in algebraic and trigonometric form}
\LANGen{
\Question{
Given \(z_{1} = 4\left (\cos \frac{5}
{3}\pi + \mathrm{i}\sin \frac{5}
{3}\pi \right )\) 
and \(z_{2} = 2\left (\cos \frac{1}
{6}\pi + \mathrm{i}\sin \frac{1}
{6}\pi \right )\),
evaluate \(\frac{z_{1}} 
{z_{2}} \).}
{\(- 2\mathrm{i}\)}{\(4\mathrm{i}\)}{\(\mathrm{i}\)}{\(-\frac{1} 
{2}\mathrm{i}\)}{}{}}
\LANGpl{
\Question{
Podano  \(z_{1} = 4\left (\cos \frac{5}
{3}\pi + \mathrm{i}\sin \frac{5}
{3}\pi \right )\) 
i \(z_{2} = 2\left (\cos \frac{1}
{6}\pi + \mathrm{i}\sin \frac{1}
{6}\pi \right )\),
oblicz  \(\frac{z_{1}} 
{z_{2}} \).}
{\(- 2\mathrm{i}\)}{\(4\mathrm{i}\)}{\(\mathrm{i}\)}{\(-\frac{1} 
{2}\mathrm{i}\)}{}{}}
\LANGcs{
\Question{
Jsou dána komplexní čísla \(z_{1} = 4\left (\cos \frac{5}
{3}\pi + \mathrm{i}\sin \frac{5}
{3}\pi \right )\) 
a \(z_{2} = 2\left (\cos \frac{1}
{6}\pi + \mathrm{i}\sin \frac{1}
{6}\pi \right )\).
Výraz \(\frac{z_{1}} 
{z_{2}} \) 
je roven:}
{\(- 2\mathrm{i}\)}{\(4\mathrm{i}\)}{\(\mathrm{i}\)}{\(-\frac{1} 
{2}\mathrm{i}\)}{}{}}
\LANGsk{
\Question{
Sú dané \(z_{1} = 4\left (\cos \frac{5}
{3}\pi + \mathrm{i}\sin \frac{5}
{3}\pi \right )\) 
a \(z_{2} = 2\left (\cos \frac{1}
{6}\pi + \mathrm{i}\sin \frac{1}
{6}\pi \right )\).
Výraz \(\frac{z_{1}} 
{z_{2}} \) sa rovná:}
{\(- 2\mathrm{i}\)}{\(4\mathrm{i}\)}{\(\mathrm{i}\)}{\(-\frac{1} 
{2}\mathrm{i}\)}{}{}}
\LANGes{
\Question{
Dados \(z_{1} = 4\left (\cos \frac{5}
{3}\pi + \mathrm{i}\sin \frac{5}
{3}\pi \right )\) 
y  \(z_{2} = 2\left (\cos \frac{1}
{6}\pi + \mathrm{i}\sin \frac{1}
{6}\pi \right )\),
calcula \(\frac{z_{1}} 
{z_{2}} \).}
{\(- 2\mathrm{i}\)}{\(4\mathrm{i}\)}{\(\mathrm{i}\)}{\(-\frac{1} 
{2}\mathrm{i}\)}{}{}}
\End

\Data\ID{1003082304}
\Updated{2020-07-18 08:07:51}
\Level{C}
\SubArea{Complex numbers in algebraic and trigonometric form}
\LANGen{
\Question{
Let \( x \) and \( y \) be  real numbers. Find the solution of the equation.
\[ (5+2\,\mathrm{i})x+(3-5\,\mathrm{i})y=124\]}
{\( x=20\text{, }y=8 \)}{\( x=8\text{, }y=20 \)}{\( x=-20\text{, }y=-8 \)}{\( x=-8\text{, }y=-20 \)}{}{}}
\LANGpl{
\Question{
Niech  \( x \) i \( y \) będą liczbami rzeczywistymi. Wskaż rozwiązanie równania.
\[ (5+2\,\mathrm{i})x+(3-5\,\mathrm{i})y=124 \]}
{\( x=20\text{, }y=8 \)}{\( x=8\text{, }y=20 \)}{\( x=-20\text{, }y=-8 \)}{\( x=-8\text{, }y=-20 \)}{}{}}
\LANGcs{
\Question{
Určete, pro která \( x,~y \in \mathbb{R}\) platí:
\[ (5+2\,\mathrm{i})x+(3-5\,\mathrm{i})y=124\]}
{\( x=20\text{, }y=8 \)}{\( x=8\text{, }y=20 \)}{\( x=-20\text{, }y=-8 \)}{\( x=-8\text{, }y=-20 \)}{}{}}
\LANGsk{
\Question{
Nech \( x \) a \( y \) sú reálne čísla. Nájdite riešenie danej rovnice.
\[ (5+2\,\mathrm{i})x+(3-5\,\mathrm{i})y=124 \]}
{\( x=20\text{, }y=8 \)}{\( x=8\text{, }y=20 \)}{\( x=-20\text{, }y=-8 \)}{\( x=-8\text{, }y=-20 \)}{}{}}
\LANGes{
\Question{
Sean \( x \) e \( y \) números reales. Calcula
\[ (5+2\,\mathrm{i})x+(3-5\,\mathrm{i})y=124\]}
{\( x=20\text{, }y=8 \)}{\( x=8\text{, }y=20 \)}{\( x=-20\text{, }y=-8 \)}{\( x=-8\text{, }y=-20 \)}{}{}}
\End

\Data\ID{1003082306}
\Updated{2021-04-24 05:04:41}
\Level{C}
\SubArea{Complex numbers in algebraic and trigonometric form}
\LANGen{
\Question{
Let \( [x;y]\in\mathbb{R}\times\mathbb{R} \). Find all \( [x;y] \) that satisfy 
\[ (3x + 2y\,\mathrm{i})\cdot(3x - 2y\,\mathrm{i}) + y^2\,\mathrm{i} = 97 + 4\,\mathrm{i}. \]}
{\( [x;y]\in\left\{[3;2], [-3;2], [3;-2], [-3;-2]\right\} \)}{\( [x;y]\in\left\{[3;2], [-3;2]\right\} \)}{\(  [x;y]\in\left\{[3;2], [3;-2]\right\}\)}{\( [x;y]\in\left\{[3;2], [-3;-2]\right\} \)}{}{}}
\LANGpl{
\Question{
Niech  \( [x;y]\in\mathbb{R}\times\mathbb{R} \). Wskaż wszystkie  \( [x;y] \), które spełniają
\[ (3x + 2y\,\mathrm{i})\cdot(3x - 2y\,\mathrm{i}) + y^2\,\mathrm{i} = 97 + 4\,\mathrm{i}. \]}
{\( [x;y]\in\left\{[3;2], [-3;2], [3;-2], [-3;-2]\right\} \)}{\( [x;y]\in\left\{[3;2], [-3;2]\right\} \)}{\(  [x;y]\in\left\{[3;2], [3;-2]\right\}\)}{\( [x;y]\in\left\{[3;2], [-3;-2]\right\} \)}{}{}}
\LANGcs{
\Question{
Nechť \( [x;y]\in\mathbb{R}\times\mathbb{R} \). Najděte všechny \( [x;y] \), pro které platí rovnice:
\[ (3x + 2y\,\mathrm{i})\cdot(3x - 2y\,\mathrm{i}) + y^2\,\mathrm{i} = 97 + 4\,\mathrm{i} \]}
{\( [x;y]\in\left\{[3;2], [-3;2], [3;-2], [-3;-2]\right\} \)}{\( [x;y]\in\left\{[3;2], [-3;2]\right\} \)}{\(  [x;y]\in\left\{[3;2], [3;-2]\right\}\)}{\( [x;y]\in\left\{[3;2], [-3;-2]\right\} \)}{}{}}
\LANGsk{
\Question{
Nech \( [x;y]\in\mathbb{R}\times\mathbb{R} \). Nájdite všetky \( [x;y] \), pre ktoré platí rovnica:
\[ (3x + 2y\,\mathrm{i})\cdot(3x- 2y\,\mathrm{i}) + y^2\,\mathrm{i} = 97 + 4\,\mathrm{i}\]}
{\( [x;y]\in\left\{[3;2], [-3;2], [3;-2], [-3;-2]\right\} \)}{\( [x;y]\in\left\{[3;2], [-3;2]\right\} \)}{\(  [x;y]\in\left\{[3;2], [3;-2]\right\}\)}{\( [x;y]\in\left\{[3;2], [-3;-2]\right\} \)}{}{}}
\LANGes{
\Question{
Sea \( [x;y]\in\mathbb{R}\times\mathbb{R} \). Calcula todos los \( [x;y] \) suponiendo que
\[ (3x + 2y\,\mathrm{i})\cdot(3x - 2y\,\mathrm{i}) + y^2\,\mathrm{i} = 97 + 4\,\mathrm{i}. \]}
{\( [x;y]\in\left\{[3;2], [-3;2], [3;-2], [-3;-2]\right\} \)}{\( [x;y]\in\left\{[3;2], [-3;2]\right\} \)}{\(  [x;y]\in\left\{[3;2], [3;-2]\right\}\)}{\( [x;y]\in\left\{[3;2], [-3;-2]\right\} \)}{}{}}
\End

\Data\ID{1003082307}
\Updated{2022-04-23 05:04:17}
\Level{C}
\SubArea{Complex numbers in algebraic and trigonometric form}
\LANGen{
\Question{
Let \( z_1 = x^2  + 9y\,\mathrm{i}-20\,\mathrm{i} \)  and  \( z_2  = 7x-12+ y^2\,\mathrm{i} \). Find all \( [x;y] \in \mathbb{R}\times\mathbb{R} \) such that  \( z_1= z_2 \).}
{\( [x;y]\in\left\{[3;4], [3;5], [4;4], [4;5]\right\} \)}{\( [x;y]\in\left\{[4;3], [4;4], [5;3], [5;4]\right\} \)}{\( [x;y]\in\left\{[-3;-4], [-3;-5], [-4;-4], [-4;-5]\right\} \)}{\( [x;y]\in\left\{[-4;-3], [-4;-4], [-5;-3], [-5;-4]\right\} \)}{}{}}
\LANGpl{
\Question{
Niech \( z_1 = x^2  + 9y\,\mathrm{i}-20\,\mathrm{i} \)  i  \( z_2  = 7x-12+ y^2\,\mathrm{i} \). Wskaż wszystkie \( [x;y] \in \mathbb{R}\times\mathbb{R} \) tak, aby  \( z_1= z_2 \).}
{\( [x;y]\in\left\{[3;4], [3;5], [4;4], [4;5]\right\} \)}{\( [x;y]\in\left\{[4;3], [4;4], [5;3], [5;4]\right\} \)}{\( [x;y]\in\left\{[-3;-4], [-3;-5], [-4;-4], [-4;-5]\right\} \)}{\( [x;y]\in\left\{[-4;-3], [-4;-4], [-5;-3], [-5;-4]\right\} \)}{}{}}
\LANGcs{
\Question{
Nechť \( z_1 = x^2  + 9y\,\mathrm{i}-20\,\mathrm{i} \)  a  \( z_2  = 7x-12+ y^2\,\mathrm{i} \). Určete všechny  \( [x;y] \in \mathbb{R}\times\mathbb{R}\), pro které se  \( z_1= z_2 \).}
{\( [x;y]\in\left\{[3;4], [3;5], [4;4], [4;5]\right\} \)}{\( [x;y]\in\left\{[4;3], [4;4], [5;3], [5;4]\right\} \)}{\( [x;y]\in\left\{[-3;-4], [-3;-5], [-4;-4], [-4;-5]\right\} \)}{\( [x;y]\in\left\{[-4;-3], [-4;-4], [-5;-3], [-5;-4]\right\} \)}{}{}}
\LANGsk{
\Question{
Nech \( z_1 = x^2  + 9y\,\mathrm{i}-20\,\mathrm{i} \)  a  \( z_2  = 7x-12+ y^2\,\mathrm{i} \). Určte všetky  \( [x;y] \in \mathbb{R}\times\mathbb{R}\), pre ktoré sa  \( z_1= z_2 \).}
{\( [x;y]\in\left\{[3;4], [3;5], [4;4], [4;5]\right\} \)}{\( [x;y]\in\left\{[4;3], [4;4], [5;3], [5;4]\right\} \)}{\( [x;y]\in\left\{[-3;-4], [-3;-5], [-4;-4], [-4;-5]\right\} \)}{\( [x;y]\in\left\{[-4;-3], [-4;-4], [-5;-3], [-5;-4]\right\} \)}{}{}}
\LANGes{
\Question{
Sean \( z_1 = x^2  + 9y\,\mathrm{i}-20\,\mathrm{i} \)  y  \( z_2  = 7x-12+ y^2\,\mathrm{i} \). Calcula todos \( [x;y] \in \mathbb{R}\times\mathbb{R} \) suponiendo que \( z_1= z_2 \).}
{\( [x;y]\in\left\{[3;4], [3;5], [4;4], [4;5]\right\} \)}{\( [x;y]\in\left\{[4;3], [4;4], [5;3], [5;4]\right\} \)}{\( [x;y]\in\left\{[-3;-4], [-3;-5], [-4;-4], [-4;-5]\right\} \)}{\( [x;y]\in\left\{[-4;-3], [-4;-4], [-5;-3], [-5;-4]\right\} \)}{}{}}
\End

\Data\ID{1003082308}
\Updated{2020-07-18 08:07:54}
\Level{C}
\SubArea{Complex numbers in algebraic and trigonometric form}
\LANGen{
\Question{
Let \( [x;y]\in\mathbb{N}\times\mathbb{N} \). Find all \( [x;y] \), that satisfy 
\[ x(8 + 4\,\mathrm{i}) + y(1 - 4\,\mathrm{i}) + 5 = x(3 +\mathrm{i}) + 6(y - 2\,\mathrm{i}) + 9\,\mathrm{i}. \]}
{There is no \( [x;y] \). (There is no solution.)}{\( [1;0] \)}{\( [0;1] \)}{\( [-1; 0]  \)}{\( [0;-1] \)}{}}
\LANGpl{
\Question{
Niech \( [x;y]\in\mathbb{N}\times\mathbb{N} \). Wyznacz wszystkie \( [x;y] \), które spełniają 
\[ x(8 + 4\,\mathrm{i}) + y(1 - 4\,\mathrm{i}) + 5 = x(3 +\mathrm{i}) + 6(y - 2\,\mathrm{i}) + 9\,\mathrm{i}. \]}
{Brak \( [x;y] \). (Brak rozwiązania.)}{\( [1;0] \)}{\( [0;1] \)}{\( [-1; 0]  \)}{\( [0;-1] \)}{}}
\LANGcs{
\Question{
Nechť \( [x;y]\in\mathbb{N}\times\mathbb{N} \). Určete všechny \( [x;y] \), když  
\[ x(8 + 4\,\mathrm{i}) + y(1 - 4\,\mathrm{i}) + 5 = x(3 +\mathrm{i}) + 6(y - 2\,\mathrm{i}) + 9\,\mathrm{i}. \]}
{Neexistuje žádná dvojice \( [x;y] \).}{\( [1;0] \)}{\( [0;1] \)}{\( [-1; 0]  \)}{\( [0;-1] \)}{}}
\LANGsk{
\Question{
Nech \( [x;y]\in\mathbb{N}\times\mathbb{N} \). Určte všetky \( [x;y] \), ak  
\[ x(8 + 4\,\mathrm{i}) + y(1 - 4\,\mathrm{i}) + 5 = x(3 +\mathrm{i}) + 6(y - 2\,\mathrm{i}) + 9\,\mathrm{i}. \]}
{Neexistuje žiadna dvojica \( [x;y] \).}{\( [1;0] \)}{\( [0;1] \)}{\( [-1; 0]  \)}{\( [0;-1] \)}{}}
\LANGes{
\Question{
Sea \( [x;y]\in\mathbb{N}\times\mathbb{N} \). Determina todos \( [x;y] \) suponiendo que 
\[ x(8 + 4\,\mathrm{i}) + y(1 - 4\,\mathrm{i}) + 5 = x(3 +\mathrm{i}) + 6(y - 2\,\mathrm{i}) + 9\,\mathrm{i}. \]}
{No hay ninguna solución.}{\( [1;0] \)}{\( [0;1] \)}{\( [-1; 0]  \)}{\( [0;-1] \)}{}}
\End

\Data\ID{2010013107}
\Updated{2023-01-23 10:01:18}
\Level{C}
\SubArea{Complex numbers in algebraic and trigonometric form}
\LANGen{
\Question{
Let \( z_1 = x^2  + 9y\,\mathrm{i}-10\,\mathrm{i} \)  and  \( z_2  = 8x-15+ y^2\,\mathrm{i} \). Find all \( [x;y] \in \mathbb{R}\times\mathbb{R} \) such that  \( z_1= \overline{z_2} \).}
{\( [x;y]\in\left\{[3;-10], [3;1], [5;-10], [5;1]\right\} \)}{\( [x;y]\in\left\{[-10;3], [1;3], [-10;5], [1;5]\right\} \)}{\( [x;y]\in\left\{[3;10], [3;-1], [5;10], [5;-1]\right\} \)}{\( [x;y]\in\left\{[-3;-10], [-3;1], [-5;-10], [-5;1]\right\} \)}{}{}}
\LANGpl{
\Question{
Niech  \( z_1 = x^2  + 9y\,\mathrm{i}-10\,\mathrm{i} \)  oraz   \( z_2  = 8x-15+ y^2\,\mathrm{i} \). Znajdź wszystkie \( [x;y] \in \mathbb{R}\times\mathbb{R} \) takie, że   \( z_1= \overline{z_2} \).}
{\( [x;y]\in\left\{[3;-10], [3;1], [5;-10], [5;1]\right\} \)}{\( [x;y]\in\left\{[-10;3], [1;3], [-10;5], [1;5]\right\} \)}{\( [x;y]\in\left\{[3;10], [3;-1], [5;10], [5;-1]\right\} \)}{\( [x;y]\in\left\{[-3;-10], [-3;1], [-5;-10], [-5;1]\right\} \)}{}{}}
\LANGcs{
\Question{
Nechť \( z_1 = x^2  + 9y\,\mathrm{i}-10\,\mathrm{i} \) a \( z_2  = 8x-15+ y^2\,\mathrm{i} \). Určete všechny \( [x;y] \in \mathbb{R}\times\mathbb{R} \), pro které \( z_1= \overline{z_2} \).}
{\( [x;y]\in\left\{[3;-10], [3;1], [5;-10], [5;1]\right\} \)}{\( [x;y]\in\left\{[-10;3], [1;3], [-10;5], [1;5]\right\} \)}{\( [x;y]\in\left\{[3;10], [3;-1], [5;10], [5;-1]\right\} \)}{\( [x;y]\in\left\{[-3;-10], [-3;1], [-5;-10], [-5;1]\right\} \)}{}{}}
\LANGsk{
\Question{
Nech \( z_1 = x^2 + 9y\,\mathrm{i}-10\,\mathrm{i} \) a \( z_2 = 8x-15+ y^2\,\mathrm{i} \). Nájdite všetky \( [x;y] \in \mathbb{R}\times\mathbb{R} \) také, že \( z_1= \overline{z_2} \).}
{\( [x;y]\in\left\{[3;-10], [3;1], [5;-10], [5;1]\right\} \)}{\( [x;y]\in\left\{[-10;3], [1;3], [-10;5], [1;5]\right\} \)}{\( [x;y]\in\left\{[3;10], [3;-1], [5;10], [5;-1]\right\} \)}{\( [x;y]\in\left\{[-3;-10], [-3;1], [-5;-10], [-5;1]\right\} \)}{}{}}
\LANGes{
\Question{
Sean \( z_1 = x^2  + 9y\,\mathrm{i}-10\,\mathrm{i} \)  y  \( z_2  = 8x-15+ y^2\,\mathrm{i} \). Halla todos los \( [x;y] \in \mathbb{R}\times\mathbb{R} \) tales que \( z_1= \overline{z_2} \).}
{\( [x;y]\in\left\{[3;-10], [3;1], [5;-10], [5;1]\right\} \)}{\( [x;y]\in\left\{[-10;3], [1;3], [-10;5], [1;5]\right\} \)}{\( [x;y]\in\left\{[3;10], [3;-1], [5;10], [5;-1]\right\} \)}{\( [x;y]\in\left\{[-3;-10], [-3;1], [-5;-10], [-5;1]\right\} \)}{}{}}
\End

\Data\ID{2010013108}
\Updated{2023-01-23 10:01:18}
\Level{C}
\SubArea{Complex numbers in algebraic and trigonometric form}
\LANGen{
\Question{
Solve the following equation for \(z\in \mathbb{C}\). By \(\overline{z }\) the complex conjugate of \(z \) is denoted.
\[ 
 2z - 3\overline{z } = 10 - 15\mathrm{i}
\]}
{\(-10 - 3\mathrm{i}\)}{\(-10 + 15\mathrm{i}\)}{\(10 + 3\mathrm{i}\)}{\(-10 + 3\mathrm{i}\)}{}{}}
\LANGpl{
\Question{
Rozwiąż następujące równanie dla \(z\in \mathbb{C}\). Przez  \(\overline{z }\) oznaczamy sprzężenie liczby zespolonej  \(z \).
\[ 
 2z - 3\overline{z } = 10 - 15\mathrm{i}
\]}
{\(-10 - 3\mathrm{i}\)}{\(-10 + 15\mathrm{i}\)}{\(10 + 3\mathrm{i}\)}{\(-10 + 3\mathrm{i}\)}{}{}}
\LANGcs{
\Question{
Vyřešte následující rovnici pro \(z\in \mathbb{C}\). Pomocí \(\overline{z }\) označujeme číslo komplexně sdružené k číslu \(z \).
\[ 
 2z - 3\overline{z } = 10 - 15\mathrm{i}
\]}
{\(-10 - 3\mathrm{i}\)}{\(-10 + 15\mathrm{i}\)}{\(10 + 3\mathrm{i}\)}{\(-10 + 3\mathrm{i}\)}{}{}}
\LANGsk{
\Question{
Vyriešte nasledujúcu rovnicu pre \(z\in \mathbb{C}\). Pomocou \(\overline{z }\) sa označuje komplexne združené číslo ku číslu \(z \).
\[
  2z - 3\overline{z } = 10 - 15\mathrm{i}
\]}
{\(-10 - 3\mathrm{i}\)}{\(-10 + 15\mathrm{i}\)}{\(10 + 3\mathrm{i}\)}{\(-10 + 3\mathrm{i}\)}{}{}}
\LANGes{
\Question{
Resuelve la siguiente ecuación para \(z\in \mathbb{C}\). Se denota por \(\overline{z }\) al conjugado del número complejo \(z \).
\[ 
 2z - 3\overline{z } = 10 - 15\mathrm{i}
\]}
{\(-10 - 3\mathrm{i}\)}{\(-10 + 15\mathrm{i}\)}{\(10 + 3\mathrm{i}\)}{\(-10 + 3\mathrm{i}\)}{}{}}
\End

\Data\ID{9000035802}
\Updated{2022-04-23 05:04:17}
\Level{C}
\SubArea{Complex numbers in algebraic and trigonometric form}
\LANGen{
\Question{
Solve the following equation for \(z\in \mathbb{C}\). By \(\overline{z }\) the complex conjugate of \(z \) is denoted.
\[ 
 3z - 2\overline{z } = 8 - 10\mathrm{i}
\]}
{\(8 - 2\mathrm{i}\)}{\(1 + 5\mathrm{i}\)}{\(8 - 10\mathrm{i}\)}{\(2 + 2\mathrm{i}\)}{}{}}
\LANGpl{
\Question{
Rozwiąż równanie jeśli  \(z\in \mathbb{C}\).
\[ 
 3z - 2\overline{z } = 8 - 10\mathrm{i}
\]}
{\(8 - 2\mathrm{i}\)}{\(1 + 5\mathrm{i}\)}{\(8 - 10\mathrm{i}\)}{\(2 + 2\mathrm{i}\)}{}{}}
\LANGcs{
\Question{
Komplexní číslo \(z\),
pro které platí 
\[ 
 3z - 2\overline{z } = 8 - 10\mathrm{i}
\]
má tvar:}
{\(8 - 2\mathrm{i}\)}{\(1 + 5\mathrm{i}\)}{\(8 - 10\mathrm{i}\)}{\(2 + 2\mathrm{i}\)}{}{}}
\LANGsk{
\Question{
Vyriešte nasledujúcu rovnicu pre \(z\in \mathbb{C}\).
\[ 
 3z - 2\overline{z } = 8 - 10\mathrm{i}
\]}
{\(8 - 2\mathrm{i}\)}{\(1 + 5\mathrm{i}\)}{\(8 - 10\mathrm{i}\)}{\(2 + 2\mathrm{i}\)}{}{}}
\LANGes{
\Question{
¿Qué forma tiene el número complejo \(z\),
suponiendo que
\[ 
 3z - 2\overline{z } = 8 - 10\mathrm{i}
\]?}
{\(8 - 2\mathrm{i}\)}{\(1 + 5\mathrm{i}\)}{\(8 - 10\mathrm{i}\)}{\(2 + 2\mathrm{i}\)}{}{}}
\End

\Data\ID{9000039103}
\Updated{2020-07-18 11:07:33}
\Level{C}
\SubArea{Complex numbers in algebraic and trigonometric form}
\LANGen{
\Question{
Assuming \(x\in \mathbb{R}\),
\(y\in \mathbb{R}\), solve
the following equation. 
\[ 
 (2 + 5\mathrm{i})x + (1 -\mathrm{i})y = 13\mathrm{i} + 8
\]}
{\(x = 3,\ y = 2\)}{\(x = 13,\ y = 8\)}{\(x = 8,\ y = 13\)}{\(x = 2,\ y = 3\)}{}{}}
\LANGpl{
\Question{
Dla \(x\in \mathbb{R}\),
\(y\in \mathbb{R}\), rozwiąż równanie.
\[ 
 (2 + 5\mathrm{i})x + (1 -\mathrm{i})y = 13\mathrm{i} + 8
\]}
{\(x = 3,\ y = 2\)}{\(x = 13,\ y = 8\)}{\(x = 8,\ y = 13\)}{\(x = 2,\ y = 3\)}{}{}}
\LANGcs{
\Question{
Za předpokladu $x\in\mathbb{R}$, $y\in\mathbb{R}$, vyřešte danou rovnici.
 \[ (2 + 5\mathrm{i})x + (1 -\mathrm{i})y = 13\mathrm{i} + 8 \]}
{\(x = 3,\ y = 2\)}{\(x = 13,\ y = 8\)}{\(x = 8,\ y = 13\)}{\(x = 2,\ y = 3\)}{}{}}
\LANGsk{
\Question{
Za predpokladu \(x\in \mathbb{R}\),
\(y\in \mathbb{R}\), vyriešte danú rovnicu.
\[ 
 (2 + 5\mathrm{i})x + (1 -\mathrm{i})y = 13\mathrm{i} + 8
\]}
{\(x = 3,\ y = 2\)}{\(x = 13,\ y = 8\)}{\(x = 8,\ y = 13\)}{\(x = 2,\ y = 3\)}{}{}}
\LANGes{
\Question{
Suponiendo que \(x\in \mathbb{R}\),
\(y\in \mathbb{R}\), resuelve la siguiente ecuación. 
\[ 
 (2 + 5\mathrm{i})x + (1 -\mathrm{i})y = 13\mathrm{i} + 8
\]}
{\(x = 3,\ y = 2\)}{\(x = 13,\ y = 8\)}{\(x = 8,\ y = 13\)}{\(x = 2,\ y = 3\)}{}{}}
\End

\Data\ID{9000039104}
\Updated{2020-07-18 11:07:52}
\Level{C}
\SubArea{Complex numbers in algebraic and trigonometric form}
\LANGen{
\Question{
Assuming \(x\in \mathbb{R}\),
\(y\in \mathbb{R}\), solve
the following equation. 
\[ 
 (3 - 2\mathrm{i})x + (5 - 7\mathrm{i})y = 1 + 3\mathrm{i}
\]}
{\(x = 2,\ y = -1\)}{\(x = -1,\ y = 2\)}{\(x = -2,\ y = -1\)}{\(x = -1,\ y = -2\)}{}{}}
\LANGpl{
\Question{
Dla  \(x\in \mathbb{R}\),
\(y\in \mathbb{R}\), rozwiąż równanie.
\[ 
 (3 - 2\mathrm{i})x + (5 - 7\mathrm{i})y = 1 + 3\mathrm{i}
\]}
{\(x = 2,\ y = -1\)}{\(x = -1,\ y = 2\)}{\(x = -2,\ y = -1\)}{\(x = -1,\ y = -2\)}{}{}}
\LANGcs{
\Question{
Za předpokladu $x\in\mathbb{R}$, $y\in\mathbb{R}$, vyřešte danou rovnici.
 \[ (3 - 2\mathrm{i})x + (5 - 7\mathrm{i})y = 1 + 3\mathrm{i} \]}
{\(x = 2,\ y = -1\)}{\(x = -1,\ y = 2\)}{\(x = -2,\ y = -1\)}{\(x = -1,\ y = -2\)}{}{}}
\LANGsk{
\Question{
Za predpokladu \(x\in \mathbb{R}\),
\(y\in \mathbb{R}\), vyriešte
danú rovnicu.
\[ 
 (3 - 2\mathrm{i})x + (5 - 7\mathrm{i})y = 1 + 3\mathrm{i}
\]}
{\(x = 2,\ y = -1\)}{\(x = -1,\ y = 2\)}{\(x = -2,\ y = -1\)}{\(x = -1,\ y = -2\)}{}{}}
\LANGes{
\Question{
Suponiendo que \(x\in \mathbb{R}\),
\(y\in \mathbb{R}\), resuelve la siguiente ecuación. 
\[ 
 (3 - 2\mathrm{i})x + (5 - 7\mathrm{i})y = 1 + 3\mathrm{i}
\]}
{\(x = 2,\ y = -1\)}{\(x = -1,\ y = 2\)}{\(x = -2,\ y = -1\)}{\(x = -1,\ y = -2\)}{}{}}
\End

\Data\ID{9000039108}
\Updated{2020-07-18 11:07:11}
\Level{C}
\SubArea{Complex numbers in algebraic and trigonometric form}
\LANGen{
\Question{
Assuming \(z\in \mathbb{C}\),
solve the following equation. By \(\overline{z }\) the complex conjugate of \(z \) is denoted.
\[ 
 2z -\mathrm{i}\, \overline{z} = 1 -\mathrm{i}
\]}
{\(z = \frac{1} 
{3} -\frac{1} 
{3}\mathrm{i}\)}{\(z = 1 + \mathrm{i}\)}{\(z = -\frac{3}
{5} + \frac{6} 
{5}\mathrm{i}\)}{\(z = -\frac{1}
{5} -\frac{3} 
{5}\mathrm{i}\)}{}{}}
\LANGpl{
\Question{
Dla \(z\in \mathbb{C}\),
rozwiąż równanie.
\[ 
 2z -\mathrm{i}\, \overline{z} = 1 -\mathrm{i}
\]}
{\(z = \frac{1} 
{3} -\frac{1} 
{3}\mathrm{i}\)}{\(z = 1 + \mathrm{i}\)}{\(z = -\frac{3}
{5} + \frac{6} 
{5}\mathrm{i}\)}{\(z = -\frac{1}
{5} -\frac{3} 
{5}\mathrm{i}\)}{}{}}
\LANGcs{
\Question{
Rovnice \(2z -\mathrm{i}\, \overline{z} = 1 -\mathrm{i}\) 
má v \(\mathbb{C}\) 
řešení:}
{\(z = \frac{1} 
{3} -\frac{1} 
{3}\mathrm{i}\)}{\(z = 1 + \mathrm{i}\)}{\(z = -\frac{3}
{5} + \frac{6} 
{5}\mathrm{i}\)}{\(z = -\frac{1}
{5} -\frac{3} 
{5}\mathrm{i}\)}{}{}}
\LANGsk{
\Question{
Za predpokladu \(z\in \mathbb{C}\),
vyriešte danú rovnicu.
\[ 
 2z -\mathrm{i}\, \overline{z} = 1 -\mathrm{i}
\]}
{\(z = \frac{1} 
{3} -\frac{1} 
{3}\mathrm{i}\)}{\(z = 1 + \mathrm{i}\)}{\(z = -\frac{3}
{5} + \frac{6} 
{5}\mathrm{i}\)}{\(z = -\frac{1}
{5} -\frac{3} 
{5}\mathrm{i}\)}{}{}}
\LANGes{
\Question{
¿Qué solución tiene la ecuación \(2z -\mathrm{i}\, \overline{z} = 1 -\mathrm{i}\) 
en \(\mathbb{C}\)?}
{\(z = \frac{1} 
{3} -\frac{1} 
{3}\mathrm{i}\)}{\(z = 1 + \mathrm{i}\)}{\(z = -\frac{3}
{5} + \frac{6} 
{5}\mathrm{i}\)}{\(z = -\frac{1}
{5} -\frac{3} 
{5}\mathrm{i}\)}{}{}}
\End

\Data\ID{9000039109}
\Updated{2020-07-18 11:07:25}
\Level{C}
\SubArea{Complex numbers in algebraic and trigonometric form}
\LANGen{
\Question{
Assuming \(z\in \mathbb{C}\),
solve the following equation. 
\[ 
 2z -\overline{iz} = 1 -\mathrm{i}
\]}
{\(z = 1 -\mathrm{i}\)}{\(z = 1 + \mathrm{i}\)}{\(z = \frac{1} 
{3} -\frac{1} 
{3}\mathrm{i}\)}{\(z = -\frac{1}
{3} + \frac{1} 
{3}\mathrm{i}\)}{}{}}
\LANGpl{
\Question{
Dla  \(z\in \mathbb{C}\),
rozwiąż równanie. 
\[ 
 2z -\overline{iz} = 1 -\mathrm{i}
\]}
{\(z = 1 -\mathrm{i}\)}{\(z = 1 + \mathrm{i}\)}{\(z = \frac{1} 
{3} -\frac{1} 
{3}\mathrm{i}\)}{\(z = -\frac{1}
{3} + \frac{1} 
{3}\mathrm{i}\)}{}{}}
\LANGcs{
\Question{
Rovnice \(2z -\overline{iz} = 1 -\mathrm{i}\) 
má v \(\mathbb{C}\) 
řešení:}
{\(z = 1 -\mathrm{i}\)}{\(z = 1 + \mathrm{i}\)}{\(z = \frac{1} 
{3} -\frac{1} 
{3}\mathrm{i}\)}{\(z = -\frac{1}
{3} + \frac{1} 
{3}\mathrm{i}\)}{}{}}
\LANGsk{
\Question{
Za predpokladu \(z\in \mathbb{C}\),
vyriešte danú rovnicu.
\[ 
 2z -\overline{iz} = 1 -\mathrm{i} \]}
{\(z = 1 -\mathrm{i}\)}{\(z = 1 + \mathrm{i}\)}{\(z = \frac{1} 
{3} -\frac{1} 
{3}\mathrm{i}\)}{\(z = -\frac{1}
{3} + \frac{1} 
{3}\mathrm{i}\)}{}{}}
\LANGes{
\Question{
Suponiendo que \(z\in \mathbb{C}\),
resuelve la siguiente ecuación. 
\[ 
 2z -\overline{iz} = 1 -\mathrm{i}
\]}
{\(z = 1 -\mathrm{i}\)}{\(z = 1 + \mathrm{i}\)}{\(z = \frac{1} 
{3} -\frac{1} 
{3}\mathrm{i}\)}{\(z = -\frac{1}
{3} + \frac{1} 
{3}\mathrm{i}\)}{}{}}
\End

\Data\ID{9000039110}
\Updated{2020-07-18 11:07:38}
\Level{C}
\SubArea{Complex numbers in algebraic and trigonometric form}
\LANGen{
\Question{
Assuming \(z\in \mathbb{C}\),
solve the following equation. 
\[ 
 \left (1 + \mathrm{i}\sqrt{3}\right )z = 1 -\mathrm{i}\sqrt{3}
\]}
{\(z = -\frac{1}
{2} -\frac{\sqrt{3}} 
 {2} \mathrm{i}\)}{\(z = \frac{\sqrt{3}} 
 {2} + \frac{1} 
{2}\mathrm{i}\)}{\(z = -\frac{1}
{2} + \frac{\sqrt{3}} 
 {2} \mathrm{i}\)}{\(z = -\frac{\sqrt{3}}
 {2} + \frac{1} 
{2}\mathrm{i}\)}{}{}}
\LANGpl{
\Question{
Dla \(z\in \mathbb{C}\),
rozwiąż równanie. 
\[ 
 \left (1 + \mathrm{i}\sqrt{3}\right )z = 1 -\mathrm{i}\sqrt{3}
\]}
{\(z = -\frac{1}
{2} -\frac{\sqrt{3}} 
 {2} \mathrm{i}\)}{\(z = \frac{\sqrt{3}} 
 {2} + \frac{1} 
{2}\mathrm{i}\)}{\(z = -\frac{1}
{2} + \frac{\sqrt{3}} 
 {2} \mathrm{i}\)}{\(z = -\frac{\sqrt{3}}
 {2} + \frac{1} 
{2}\mathrm{i}\)}{}{}}
\LANGcs{
\Question{
Za předpokladu $z\in\mathbb{C}$, vyřešte danou rovnici. \[ \left (1 + \mathrm{i}\sqrt{3}\right )z = 1 -\mathrm{i}\sqrt{3}\]}
{\(z = -\frac{1}
{2} -\frac{\sqrt{3}} 
 {2} \mathrm{i}\)}{\(z = \frac{\sqrt{3}} 
 {2} + \frac{1} 
{2}\mathrm{i}\)}{\(z = -\frac{1}
{2} + \frac{\sqrt{3}} 
 {2} \mathrm{i}\)}{\(z = -\frac{\sqrt{3}}
 {2} + \frac{1} 
{2}\mathrm{i}\)}{}{}}
\LANGsk{
\Question{
Za predpokladu \(z\in \mathbb{C}\),
vyriešte danú rovnicu.
\[ 
 \left (1 + \mathrm{i}\sqrt{3}\right )z = 1 -\mathrm{i}\sqrt{3}
\]}
{\(z = -\frac{1}
{2} -\frac{\sqrt{3}} 
 {2} \mathrm{i}\)}{\(z = \frac{\sqrt{3}} 
 {2} + \frac{1} 
{2}\mathrm{i}\)}{\(z = -\frac{1}
{2} + \frac{\sqrt{3}} 
 {2} \mathrm{i}\)}{\(z = -\frac{\sqrt{3}}
 {2} + \frac{1} 
{2}\mathrm{i}\)}{}{}}
\LANGes{
\Question{
Suponiendo que \(z\in \mathbb{C}\),
resuelve la siguiente ecuación. 
\[ 
 \left (1 + \mathrm{i}\sqrt{3}\right )z = 1 -\mathrm{i}\sqrt{3}
\]}
{\(z = -\frac{1}
{2} -\frac{\sqrt{3}} 
 {2} \mathrm{i}\)}{\(z = \frac{\sqrt{3}} 
 {2} + \frac{1} 
{2}\mathrm{i}\)}{\(z = -\frac{1}
{2} + \frac{\sqrt{3}} 
 {2} \mathrm{i}\)}{\(z = -\frac{\sqrt{3}}
 {2} + \frac{1} 
{2}\mathrm{i}\)}{}{}}
\End
